% Le Proche et le Moyen Orient : zones de tension sans fin — Histoire Tle Tech — Cours signé OnBuch+
% Programme officiel MINESEC. Compiler avec : tectonic N03-proche-moyen-orient-zones-tension-sans-fin.tex
\newcommand{\DOCMATIERE}{Histoire}
\newcommand{\DOCNIVEAU}{Tle Tech}
\newcommand{\DOCMODULE}{Histoire – classe de Terminale technique}
\newcommand{\DOCLECON}{N03}
\newcommand{\DOCTITRE}{Le Proche et le Moyen Orient : zones de tension sans fin}
% OnBuch+ — préambule « cours signés » ENRICHI (compilé avec tectonic / XeLaTeX)
% Système visuel « Pop » d'OnBuch+ : fond crème (jamais blanc), bordures encre
% épaisses, ombres « dures » décalées, titres Archivo Black en capitales.
% Version MATHÉMATIQUES : graphes (pgfplots/tikz), boîtes pédagogiques
% (définition, propriété, méthode, attention, le savais-tu, exercice résolu,
% exercice, TP, bilan), page de garde et signature OnBuch+ ancrée en bas.
%
% Macros attendues dans main.tex AVANT \input{preamble} :
%   \DOCMATIERE  \DOCNIVEAU  \DOCMODULE  \DOCLECON (numéro)  \DOCTITRE
\documentclass[11pt]{article}

\usepackage{fontspec}
\usepackage[french]{babel}
\usepackage[a4paper,margin=2cm,top=3.4cm,bottom=2.8cm,headheight=1.4cm,footskip=1.3cm]{geometry}
\usepackage{amsmath,amssymb,amsfonts}
\usepackage{mathtools} % \xmapsto, \coloneqq... souvent utilisés par les modèles
\usepackage{cancel} % simplifications barrées (bilans de forces)
% \abs{z} (module, valeur absolue) et \norm{u} (norme), utilisés par les modèles
\providecommand{\abs}{}\let\abs\relax\DeclarePairedDelimiter{\abs}{\lvert}{\rvert}
\providecommand{\norm}{}\let\norm\relax\DeclarePairedDelimiter{\norm}{\lVert}{\rVert}
\usepackage[table]{xcolor} % [table] : fournit aussi \rowcolors (alternance de lignes)
\usepackage{pagecolor}
\usepackage{tikz}
\usetikzlibrary{arrows.meta,positioning,calc,shapes.geometric,shapes.misc,shapes.arrows,decorations.pathmorphing,decorations.pathreplacing,decorations.markings,patterns,shadows,mindmap,trees,fit,backgrounds,matrix,intersections,angles,quotes}
\usepackage{pgfplots}
\usepackage[version=4]{mhchem} % équations nucléaires \ce{^{235}_{92}U}
\usepackage[european,siunitx]{circuitikz} % schémas électriques
\pgfplotsset{compat=1.18}
\usepgfplotslibrary{fillbetween,groupplots} % aires entre courbes (calcul intégral)
\usepackage{enumitem}
\usepackage{fancyhdr}
% [most] charge toutes les bibliothèques tcolorbox (listings, minted, raster,
% documentation...) et gonfle inutilement la mémoire TeX sur un document long
% (~300 boîtes) : on ne charge que ce qui est réellement utilisé.
\usepackage[skins,breakable]{tcolorbox}
\usepackage{graphicx}
\usepackage{colortbl}
\usepackage{booktabs}
\usepackage{tabularx}
\usepackage{multirow}
\usepackage{diagbox} % case d'en-tête diagonale des tableaux à double entrée
% Colonnes extensibles centrée / gauche / droite (C, L, R), souvent utilisées par les modèles
\newcolumntype{C}{>{\centering\arraybackslash}X}
\newcolumntype{L}{>{\raggedright\arraybackslash}X}
\newcolumntype{R}{>{\raggedleft\arraybackslash}X}
\usepackage{array}
\usepackage{multicol}
\usepackage{siunitx}
% parse-numbers=false : accepte aussi \SI{2,5 \times 10^{-3}}{...}
\sisetup{locale=FR,output-decimal-marker={,},group-minimum-digits=4,parse-numbers=false}
\DeclareSIUnit\ohm{\ensuremath{\symup{\Omega}}} % U+2126 absent de la police
\DeclareSIUnit\division{div} % réglages d'oscilloscope (ms/div, V/div)
\DeclareSIUnit\tr{tr} % tours (vitesse de rotation, ex. \SI{1500}{\tr\per\minute})
\DeclareSIUnit\molar{M} % concentration molaire (ex. \SI{3}{\molar}, \SI{1}{\micro\molar})
\DeclareSIUnit\an{an} % année (le modèle confond parfois avec \year, un compteur TeX) — ex. \SI{5}{\kilo\meter\per\an}
\DeclareSIUnit\FCFA{FCFA} % franc CFA (ex. \SI{500}{\FCFA\per\kilo\gram})
\DeclareSIUnit\degreeBrix{\textdegree Brix} % teneur en sucre d'un sirop/jus (réfractométrie)
\DeclareSIUnit\month{mois} % le modèle utilise parfois \month, un compteur TeX, comme unité de durée
\DeclareSIUnit\microliter{\micro\liter} % le modèle écrit parfois \microliter en un seul token
\DeclareSIUnit\million{millions} % dénombrement (ex. \SI{15}{\million\per\mL} spermatozoïdes)
\usepackage{qrcode}
% \degree : commande fréquemment utilisée par les modèles pour un angle ou une
% température en dehors de \SI{}{} (non fournie par les paquets chargés).
\providecommand{\degree}{^{\circ}}
% \qed (fin de démonstration), utilisé par les modèles sans amsthm
\providecommand{\qed}{\hfill$\square$}
% \barre : double barre des valeurs interdites dans un tableau de variations (tkz-tab)
\providecommand{\barre}{\ensuremath{\|}}
% environnement proof (amsthm non chargé) utilisé par les modèles
\ifdefined\proof\else\newenvironment{proof}[1][Démonstration]{\par\noindent\textit{#1.}\ }{\qed\par}\fi
\usepackage{lastpage}
\usepackage{hyperref}
\hypersetup{hidelinks}

% ============================================================
% Palette « Pop » OnBuch+ — mêmes hex que la classe P (plus_ui.dart)
% ============================================================
\definecolor{poporange}{HTML}{F59321}
\definecolor{popdark}{HTML}{E07A0C}
\definecolor{popcream}{HTML}{FFF6E8}
\definecolor{popink}{HTML}{1C1714}
\definecolor{poppurple}{HTML}{7A5AE0}
\definecolor{popgreen}{HTML}{2E9E6A}
\definecolor{poppink}{HTML}{E0457B}
\definecolor{popblue}{HTML}{2F7DE1}
\definecolor{popgold}{HTML}{C98A12}
\definecolor{popmuted}{HTML}{8A7F76}
\definecolor{popline}{HTML}{EADFD2}
\definecolor{popgray}{HTML}{8A7F76}
\colorlet{popgrey}{popgray}
% Alias tolérants (le modèle nomme parfois les couleurs différemment)
\colorlet{popwhite}{white}
\colorlet{popblack}{popink}
\colorlet{popred}{poppink}
\colorlet{popyellow}{popgold}
% Teintes claires (fonds de tableaux/encadrés) — le modèle nomme parfois
% les couleurs avec un suffixe L (ex. popblueL) sans les avoir définies.
\colorlet{poporangeL}{poporange!20}
\colorlet{popdarkL}{popdark!20}
\colorlet{poppurpleL}{poppurple!20}
\colorlet{popgreenL}{popgreen!20}
\colorlet{poppinkL}{poppink!20}
\colorlet{popblueL}{popblue!20}
\colorlet{popgoldL}{popgold!20}
\colorlet{popredL}{popred!20}
\colorlet{popgrayL}{popgray!20}
% Teintes claires pour fonds de tableaux / zones
\colorlet{popblueL}{popblue!10!white}
\colorlet{poporangeL}{poporange!14!white}
\colorlet{popgreenL}{popgreen!12!white}
\colorlet{poppurpleL}{poppurple!12!white}
\colorlet{poppinkL}{poppink!12!white}
\colorlet{popgoldL}{popgold!14!white}

\pagecolor{popcream}

% ============================================================
% Typographie
% ============================================================
\defaultfontfeatures{Ligatures=TeX}
\setmainfont{PlusJakartaSans-Medium.ttf}[Path=../../../fonts/,BoldFont=PlusJakartaSans-Bold.ttf]
% Maths en sans-serif (Fira Math) avec chiffres et lettres droites de Plus
% Jakarta, pour que les formules se fondent dans le texte.
\usepackage[math-style=ISO,bold-style=ISO]{unicode-math}
\setmathfont{FiraMath-Regular.otf}[Path=../../../fonts/]
\setmathfont{PlusJakartaSans-Medium.ttf}[Path=../../../fonts/,range={up/num,up/{Latin,latin},\mathup}]
\usepackage{icomma} % virgule décimale sans espace en mode math
% Caractères mathématiques Unicode absents des polices de texte (Plus Jakarta,
% Archivo Black) : rendus par leur équivalent mathématique, en texte comme en
% formule — sinon ils s'affichent en carré vide (ex. « Arithmétique dans ℤ »).
\usepackage{newunicodechar}
\newunicodechar{ℝ}{\ensuremath{\mathbb{R}}}
\newunicodechar{ℤ}{\ensuremath{\mathbb{Z}}}
\newunicodechar{ℂ}{\ensuremath{\mathbb{C}}}
\newunicodechar{ℕ}{\ensuremath{\mathbb{N}}}
\newunicodechar{ℚ}{\ensuremath{\mathbb{Q}}}
\newunicodechar{𝔻}{\ensuremath{\mathbb{D}}}
\newunicodechar{α}{\ensuremath{\alpha}}
\newunicodechar{β}{\ensuremath{\beta}}
\newunicodechar{γ}{\ensuremath{\gamma}}
\newunicodechar{δ}{\ensuremath{\delta}}
\newunicodechar{ε}{\ensuremath{\varepsilon}}
\newunicodechar{θ}{\ensuremath{\theta}}
\newunicodechar{λ}{\ensuremath{\lambda}}
\newunicodechar{μ}{\ensuremath{\mu}}
\newunicodechar{σ}{\ensuremath{\sigma}}
\newunicodechar{τ}{\ensuremath{\tau}}
\newunicodechar{φ}{\ensuremath{\varphi}}
\newunicodechar{ψ}{\ensuremath{\psi}}
\newunicodechar{ω}{\ensuremath{\omega}}
\newunicodechar{Σ}{\ensuremath{\Sigma}}
% majuscules produites par \MakeUppercase dans les titres (α-aminés -> Α)
\newunicodechar{Α}{\ensuremath{\alpha}}
\newunicodechar{Β}{\ensuremath{\beta}}
\newunicodechar{Γ}{\ensuremath{\Gamma}}
\newunicodechar{Δ}{\ensuremath{\Delta}}
\newunicodechar{Ε}{\ensuremath{\varepsilon}}
\newunicodechar{Θ}{\ensuremath{\Theta}}
\newunicodechar{Λ}{\ensuremath{\Lambda}}
\newunicodechar{Μ}{\ensuremath{\mu}}
\newunicodechar{Τ}{\ensuremath{\tau}}
\newunicodechar{Φ}{\ensuremath{\Phi}}
\newunicodechar{Ψ}{\ensuremath{\Psi}}
\newunicodechar{Ω}{\ensuremath{\Omega}}
\newunicodechar{↦}{\ensuremath{\mapsto}}
\newunicodechar{∈}{\ensuremath{\in}}
\newunicodechar{∉}{\ensuremath{\notin}}
\newunicodechar{∧}{\ensuremath{\wedge}}
\newunicodechar{⇔}{\ensuremath{\Leftrightarrow}}
\newunicodechar{⟺}{\ensuremath{\Longleftrightarrow}}
\newunicodechar{⇒}{\ensuremath{\Rightarrow}}
\newunicodechar{⟹}{\ensuremath{\Longrightarrow}}
\newunicodechar{∘}{\ensuremath{\circ}}
\newunicodechar{∪}{\ensuremath{\cup}}
\newunicodechar{∩}{\ensuremath{\cap}}
\newunicodechar{≡}{\ensuremath{\equiv}}
\newunicodechar{⊂}{\ensuremath{\subset}}
\newunicodechar{⊥}{\ensuremath{\perp}}
\newunicodechar{∃}{\ensuremath{\exists}}
\newunicodechar{∀}{\ensuremath{\forall}}
\newunicodechar{∛}{\ensuremath{\sqrt[3]{\,}}}
\newunicodechar{∜}{\ensuremath{\sqrt[4]{\,}}}
\newunicodechar{⃗}{\ensuremath{{}^{\to}}}
\newunicodechar{ⁿ}{\ifmmode^{n}\else\textsuperscript{\ensuremath{n}}\fi}
\newunicodechar{ˣ}{\ifmmode^{x}\else\textsuperscript{\ensuremath{x}}\fi}
\newunicodechar{ᵇ}{\ifmmode^{b}\else\textsuperscript{\ensuremath{b}}\fi}
\newunicodechar{ʳ}{\ifmmode^{r}\else\textsuperscript{\ensuremath{r}}\fi}
\newunicodechar{ᵘ}{\ifmmode^{u}\else\textsuperscript{\ensuremath{u}}\fi}
\newunicodechar{ʸ}{\ifmmode^{y}\else\textsuperscript{\ensuremath{y}}\fi}
\newunicodechar{ᵃ}{\ifmmode^{a}\else\textsuperscript{\ensuremath{a}}\fi}
\newunicodechar{ᵉ}{\ifmmode^{e}\else\textsuperscript{\ensuremath{e}}\fi}
\newunicodechar{ᵏ}{\ifmmode^{k}\else\textsuperscript{\ensuremath{k}}\fi}
\newunicodechar{ᵖ}{\ifmmode^{p}\else\textsuperscript{\ensuremath{p}}\fi}
\newunicodechar{ᴺ}{\ifmmode^{N}\else\textsuperscript{\ensuremath{N}}\fi}
\newunicodechar{ᶜ}{\ifmmode^{c}\else\textsuperscript{\ensuremath{c}}\fi}
\newunicodechar{ʰ}{\ifmmode^{h}\else\textsuperscript{\ensuremath{h}}\fi}
\newunicodechar{ᵠ}{\ifmmode^{q}\else\textsuperscript{\ensuremath{q}}\fi}
\newunicodechar{ₙ}{\ifmmode_{n}\else\textsubscript{\ensuremath{n}}\fi}
\newunicodechar{ₐ}{\ifmmode_{a}\else\textsubscript{\ensuremath{a}}\fi}
\newunicodechar{ᵢ}{\ifmmode_{i}\else\textsubscript{\ensuremath{i}}\fi}
\newunicodechar{ᵣ}{\ifmmode_{r}\else\textsubscript{\ensuremath{r}}\fi}
\newunicodechar{ₖ}{\ifmmode_{k}\else\textsubscript{\ensuremath{k}}\fi}
\newunicodechar{ᵦ}{\ifmmode_{\beta}\else\textsubscript{\ensuremath{\beta}}\fi}
\newunicodechar{ₓ}{\ifmmode_{x}\else\textsubscript{\ensuremath{x}}\fi}
\newfontfamily\popdisplay{ArchivoBlack-Regular.ttf}[Path=../../../fonts/]
\newfontfamily\poplogo{SpaceGrotesk-Bold.ttf}[Path=../../../fonts/]
\newfontfamily\popsemi{PlusJakartaSans-SemiBold.ttf}[Path=../../../fonts/]
\newfontfamily\popxbold{PlusJakartaSans-ExtraBold.ttf}[Path=../../../fonts/]
\newfontfamily\popxboldls{PlusJakartaSans-ExtraBold.ttf}[Path=../../../fonts/,LetterSpace=3.0]

\setlist[enumerate]{leftmargin=1.7em,itemsep=5pt,topsep=4pt}
\setlist[itemize]{leftmargin=1.7em,itemsep=4pt,topsep=4pt,label={\color{poporange}\textbullet}}
\setlength{\parindent}{0pt}
\setlength{\parskip}{5pt}
\renewcommand{\arraystretch}{1.25}

% pgfplots : style Pop par défaut
\pgfplotsset{
  every axis/.append style={axis line style={popink,line width=1pt},tick style={popink},
    grid style={popline,line width=0.5pt},label style={font=\small},tick label style={font=\footnotesize}},
  popcurve/.style={poporange,line width=1.8pt,no markers,smooth},
}
% Même style disponible dans un \draw TikZ simple (hors pgfplots)
\tikzset{popcurve/.style={poporange,line width=1.8pt}}
\tikzset{popgrid/.style={popline,very thin}}
\colorlet{popcurve}{poporange} % le nom du style est parfois utilisé comme couleur

% ============================================================
% Pill / squiggle
% ============================================================
\newcommand{\poppill}[3]{%
  \tikz[baseline=-0.55ex]\node[draw=popink,line width=1.2pt,rounded corners=2.6pt,%
    fill=#2,inner xsep=6.5pt,inner ysep=3pt]{%
    \popxboldls\fontsize{7}{7}\selectfont\color{#3}\MakeUppercase{#1}};%
}
\newcommand{\popsquiggle}[1]{%
  \tikz[baseline=-0.4ex]{
    \draw[#1,line width=2.6pt,line cap=round]
      plot [smooth] coordinates {(0,0.09) (0.34,0.01) (0.68,0.21) (1.02,0.01) (1.36,0.10)};
  }%
}
% Mot-clé surligné (marqueur orange sous le texte)
\newcommand{\cle}[1]{{\popxbold\color{popdark}#1}}

% ============================================================
% En-tête / pied
% ============================================================
\pagestyle{fancy}
\fancyhf{}
\renewcommand{\headrulewidth}{0pt}
\renewcommand{\footrulewidth}{0pt}
\lhead{\raisebox{-0.15ex}{\poplogo\fontsize{13}{13}\selectfont\color{popink}OnBuch\color{poporange}+}}
\rhead{\poppill{\DOCMATIERE\ \textbullet\ \DOCNIVEAU\ \textbullet\ Leçon \DOCLECON}{white}{popink}}
\lfoot{\popxbold\fontsize{7.6}{7.6}\selectfont\color{popmuted}\MakeUppercase{Cours signé OnBuch+}\ \textbullet\ {\popsemi\DOCTITRE}}
\rfoot{\tikz[baseline=-0.6ex]\node[rounded corners=8.5pt,fill=popink,minimum height=17pt,minimum width=17pt,inner xsep=5pt,inner sep=0]{\popxbold\fontsize{8}{8}\selectfont\color{popcream}\thepage\,/\,\pageref*{LastPage}};}

% ============================================================
% Titres
% ============================================================
\newcommand{\chaptitre}[2]{%
  \begin{tcolorbox}[enhanced,colback=poporange,colframe=popink,boxrule=2.6pt,arc=11pt,
      drop shadow=popink,left=15pt,right=15pt,top=12pt,bottom=11pt,before skip=3pt,after skip=18pt]
    \poppill{#2}{popink}{popcream}\par\vspace{9pt}
    {\raggedright\hyphenpenalty=10000\exhyphenpenalty=10000\popdisplay\fontsize{22}{24}\selectfont\color{popcream}\MakeUppercase{#1}\par}
  \end{tcolorbox}
}
% Section numérotée : pastille orange avec le numéro + titre Archivo
\newcounter{popsec}
\newcommand{\coursec}[1]{%
  \stepcounter{popsec}\setcounter{popsub}{0}%
  \par\vspace{16pt}\noindent
  \begin{minipage}{\linewidth}
  \tikz[baseline=-0.7ex]\node[draw=popink,line width=1.6pt,fill=poporange,rounded corners=4pt,minimum size=22pt,inner sep=0]{\popdisplay\fontsize{12}{12}\selectfont\color{popcream}\thepopsec};%
  \hspace{8pt}{\popdisplay\fontsize{15}{17}\selectfont\color{popink}\MakeUppercase{#1}}\par
  \vspace{4pt}\noindent{\color{poporange}\rule{36pt}{3.6pt}}
  \end{minipage}\par\nopagebreak\vspace{8pt}
}
\newcounter{popsub}
\newcommand{\courssub}[1]{%
  \stepcounter{popsub}%
  \par\vspace{9pt}\noindent{\popxbold\fontsize{11.5}{13}\selectfont\color{popink}{\color{poporange}\thepopsec.\thepopsub}\hspace{6pt}#1}\par\nopagebreak\vspace{4pt}
}

% ============================================================
% Boîtes pédagogiques — même recette : fond blanc, bordure épaisse colorée,
% ombre dure encre, badge plein. Titre optionnel [#1].
% ============================================================
\tcbset{popbase/.style={breakable,enhanced,colback=white,boxrule=2.3pt,arc=9pt,
  shadow={3pt}{-3pt}{0pt}{popink},left=14pt,right=14pt,top=11pt,bottom=11pt,before skip=11pt,after skip=15pt,
  fontupper=\popsemi\fontsize{10.6}{14.5}\selectfont\color{popink},
  fontlower=\popsemi\fontsize{10.6}{14.5}\selectfont\color{popink}}}
% Coche dessinée (le glyphe ✓ n'existe pas dans les polices du document)
\newcommand{\popcheck}[1]{\tikz[baseline=-0.5ex]\draw[#1,line width=1.6pt,line cap=round,line join=round](-0.13,0.0)--(-0.04,-0.09)--(0.13,0.1);}
\providecommand{\coche}{\popcheck{popgreen}}
\providecommand{\resultat}[1]{\par\smallskip\noindent\fcolorbox{popgreen}{popgreenL}{\parbox{\dimexpr\linewidth-2\fboxsep-2\fboxrule\relax}{\textbf{Résultat.} #1}}\par\smallskip}
\newcommand{\popboxhead}[3]{\poppill{#1}{#2}{white}\ifx\relax#3\relax\else\hspace{6pt}{\popxbold\fontsize{10.6}{12}\selectfont\color{popink}#3}\fi\par\vspace{7pt}}

\newtcolorbox{aretenir}[1][]{popbase,colframe=popgreen,before upper={\popboxhead{À retenir}{popgreen}{#1}}}
\newtcolorbox{exemplebox}[1][]{popbase,colframe=poporange,before upper={\popboxhead{Exemple}{poporange}{#1}}}
\newtcolorbox{definition}[1][]{popbase,colframe=popblue,before upper={\popboxhead{Définition}{popblue}{#1}}}
\newtcolorbox{propriete}[1][]{popbase,colframe=poppurple,before upper={\popboxhead{Propriété}{poppurple}{#1}}}
\newtcolorbox{methode}[1][]{popbase,colframe=popgold,before upper={\popboxhead{Méthode}{popgold}{#1}}}
\newtcolorbox{attention}[1][]{popbase,colframe=poppink,before upper={\popboxhead{Attention}{poppink}{#1}}}
\newtcolorbox{experience}[1][]{popbase,colframe=popblue,colback=popblueL,before upper={\popboxhead{Expérience}{popblue}{#1}}}
% « Le savais-tu ? » : carte violette pleine (contexte, culture, Cameroun)
\newtcolorbox{savaistu}[1][]{popbase,colback=poppurple,colframe=popink,
  fontupper=\popsemi\fontsize{10.4}{14.2}\selectfont\color{white},
  before upper={\poppill{Le savais-tu\,?}{white}{poppurple}\ifx\relax#1\relax\else\hspace{6pt}{\popxbold\color{white}#1}\fi\par\vspace{7pt}}}
% Exercice résolu : énoncé en haut, \tcblower puis la solution sur fond crème
\newcounter{popexo}\newcounter{popexr}
\newtcolorbox{exoresolu}[1][]{popbase,colframe=poporange,colbacklower=popcream,
  segmentation style={popink,line width=1.2pt,dashed},
  before upper={\stepcounter{popexr}\popboxhead{Exercice résolu \thepopexr}{poporange}{#1}},
  before lower={\poppill{Solution}{popink}{popcream}\par\vspace{6pt}}}
% Exercice (non résolu en place) — la correction est dans la partie Corrigés
\newtcolorbox{exercice}[1][]{popbase,colframe=popink,
  before upper={\stepcounter{popexo}\popboxhead{Exercice \thepopexo}{popink}{#1}}}
% Corrigé
\newtcolorbox{corrige}[1][]{popbase,colframe=popgreen,colback=popgreenL,
  before upper={\popboxhead{Corrigé}{popgreen}{#1}}}
% Objectifs / prérequis / bilan
\newtcolorbox{objectifs}{popbase,colframe=popink,colback=poporangeL,
  before upper={\popboxhead{Objectifs de la leçon}{popink}{}}}
\newtcolorbox{prerequis}{popbase,colframe=popink,colback=popblueL,
  before upper={\popboxhead{Prérequis}{popblue}{}}}
\newtcolorbox{bilan}{popbase,colframe=popink,colback=white,boxrule=2.8pt,
  before upper={\popboxhead{Fiche bilan}{poporange}{L'essentiel à connaître pour l'examen}}}
% Figure encadrée avec légende
\newtcolorbox{popfigure}[1][]{enhanced,breakable=false,colback=white,colframe=popline,boxrule=1.4pt,arc=8pt,
  left=10pt,right=10pt,top=10pt,bottom=8pt,before skip=10pt,after skip=12pt,halign=center,
  lower separated=false,halign lower=center,
  fontlower=\popsemi\fontsize{9}{11.5}\selectfont\color{popmuted},
  #1}
\newcommand{\legende}[1]{\par\vspace{6pt}{\popsemi\fontsize{9}{11.5}\selectfont\color{popmuted}#1}}

% ============================================================
% Signature OnBuch+
% ============================================================
\newcommand{\poplogoinline}[1]{{\poplogo\fontsize{#1}{#1}\selectfont\color{popink}OnBuch\color{poporange}+}}
\newcommand{\signatureonbuch}{%
  \begin{tcolorbox}[enhanced,colback=popink,colframe=popink,boxrule=2.2pt,arc=10pt,
      drop shadow=poporange,left=14pt,right=14pt,top=10pt,bottom=10pt,before skip=0pt,after skip=0pt]
    \tikz[baseline=-0.7ex]\node[circle,fill=popgreen,draw=popcream,line width=1.3pt,minimum size=22pt,inner sep=0]{\popcheck{white}};%
    \hspace{9pt}\begin{minipage}[c]{0.62\linewidth}
      {\popxbold\fontsize{12}{13}\selectfont\color{popcream}Signé\ \textbullet\ {\poplogo OnBuch\color{poporange}+}}\par\vspace{2pt}
      {\raggedright\popsemi\fontsize{8}{10.5}\selectfont\color{popline}\MakeUppercase{Équipe pédagogique OnBuch+}\\\MakeUppercase{Conforme au programme officiel MINESEC}\par}
    \end{minipage}\hfill
    {\poplogo\fontsize{22}{22}\selectfont\color{popcream}OnBuch\color{poporange}+}
  \end{tcolorbox}%
}

% ============================================================
% Page de garde
% #1 titre · #2 matière · #3 niveau · #4 module · #5 n° leçon · #6 durée ·
% #7 description / objectifs (texte ou itemize)
% ============================================================
\newcommand{\pagegarde}[7]{%
  \thispagestyle{empty}
  \newgeometry{a4paper,margin=1.7cm,top=1.6cm,bottom=1.4cm}%
  \begin{tikzpicture}[remember picture,overlay]
    \fill[poporange!18] ([xshift=-3.2cm,yshift=-3.6cm]current page.north east) circle (4.6cm);
    \fill[poppurple!14] ([xshift=2.4cm,yshift=7.6cm]current page.south west) circle (3.8cm);
  \end{tikzpicture}%
  {\poplogo\fontsize{30}{30}\selectfont\color{popink}OnBuch\color{poporange}+}\hfill
  \raisebox{0.6ex}{\poppill{Cours signé \textbullet\ Premium}{popink}{popcream}}\par
  \vspace{1pt}\popsquiggle{poppurple}\par
  \vspace{8pt}
  \begin{tcolorbox}[enhanced,colback=poporange,colframe=popink,boxrule=2.8pt,arc=13pt,
      drop shadow=popink,left=18pt,right=18pt,top=12pt,bottom=13pt,after skip=10pt]
    \poppill{#2 \textbullet\ #3}{popink}{popcream}\hspace{5pt}\poppill{#4}{white}{popink}\par\vspace{8pt}
    {\popdisplay\fontsize{14}{14}\selectfont\color{popink}LEÇON #5}\par\vspace{4pt}
    {\raggedright\hyphenpenalty=10000\exhyphenpenalty=10000\popdisplay\fontsize{25}{27}\selectfont\color{popcream}\MakeUppercase{#1}\par}
    \vspace{8pt}
    {\popxbold\fontsize{9.5}{9.5}\selectfont\color{popink}\MakeUppercase{Durée conseillée : #6}}
  \end{tcolorbox}
  \begin{tcolorbox}[enhanced,colback=white,colframe=popink,boxrule=2.2pt,arc=10pt,
      drop shadow=popink,left=16pt,right=16pt,top=10pt,bottom=10pt,after skip=8pt]
    \poppill{Dans cette leçon}{poppurple}{white}\par\vspace{5pt}
    \popsemi\fontsize{9.3}{12.6}\selectfont\color{popink}#7
  \end{tcolorbox}
  \noindent\begin{minipage}[t]{0.487\linewidth}
    \begin{tcolorbox}[enhanced,colback=popgreen,colframe=popink,boxrule=2.2pt,arc=10pt,
        drop shadow=popink,left=11pt,right=11pt,top=10pt,bottom=10pt,height=3.1cm,valign=center]
      \tcbox[colback=white,colframe=popink,boxrule=1.3pt,arc=5pt,boxsep=0pt,nobeforeafter,
        left=3pt,right=3pt,top=3pt,bottom=3pt,valign=center]{\qrcode[height=1.9cm]{https://play.google.com/store/apps/details?id=com.luvvix.onbuch&hl=fr}}%
      \hspace{8pt}\begin{minipage}[c]{0.5\linewidth}
        {\popdisplay\fontsize{11}{12.5}\selectfont\color{white}\MakeUppercase{Télécharge OnBuch}}\par\vspace{4pt}
        {\raggedright\popsemi\fontsize{8.4}{11}\selectfont\color{white}Gratuit \textbullet\ Google Play\\Tuteur IA, annales, concours, résultats\par}
      \end{minipage}
    \end{tcolorbox}
  \end{minipage}\hfill
  \begin{minipage}[t]{0.487\linewidth}
    \begin{tcolorbox}[enhanced,colback=poppurple,colframe=popink,boxrule=2.2pt,arc=10pt,
        drop shadow=popink,left=11pt,right=11pt,top=10pt,bottom=10pt,height=3.1cm,valign=center]
      \tikz[baseline=-0.7ex]\node[circle,fill=white,draw=popink,line width=1.3pt,minimum size=1.6cm,inner sep=0]{\popdisplay\fontsize{24}{24}\selectfont\color{poppurple}+};%
      \hspace{8pt}\begin{minipage}[c]{0.6\linewidth}
        {\popdisplay\fontsize{11}{12.5}\selectfont\color{white}\MakeUppercase{Passe à OnBuch+}}\par\vspace{4pt}
        {\raggedright\popsemi\fontsize{8.4}{11}\selectfont\color{white}Tous les cours signés, Tuteur IA illimité, fiches PDF — onglet OnBuch+ de l'app\par}
      \end{minipage}
    \end{tcolorbox}
  \end{minipage}
  \vspace{8pt}
  \popcontenu
  \vfill
  \signatureonbuch
  \restoregeometry
  \newpage
}

% Rangée « Ce cours contient » (page de garde)
\newcommand{\poptile}[3]{% #1 couleur #2 gros texte #3 légende
  \begin{tcolorbox}[enhanced,colback=#1,colframe=popink,boxrule=2pt,arc=9pt,shadow={2.5pt}{-2.5pt}{0pt}{popink},
      left=6pt,right=6pt,top=9pt,bottom=9pt,halign=center,valign=center,height=1.75cm,width=\linewidth,nobeforeafter]
    {\popdisplay\fontsize{12.5}{13}\selectfont\color{white}\MakeUppercase{#2}}\par\vspace{3pt}
    {\centering\popsemi\fontsize{7.8}{9.5}\selectfont\color{white}#3\par}
  \end{tcolorbox}}
\newcommand{\popcontenu}{%
  \par\noindent{\popxboldls\fontsize{7.5}{7.5}\selectfont\color{popmuted}CE COURS SIGNÉ CONTIENT}\par\vspace{7pt}
  \noindent\begin{minipage}[t]{0.235\linewidth}\poptile{popblue}{Cours}{complet, illustré, pas à pas}\end{minipage}\hfill
  \begin{minipage}[t]{0.235\linewidth}\poptile{poporange}{Méthodes}{et exercices résolus}\end{minipage}\hfill
  \begin{minipage}[t]{0.235\linewidth}\poptile{poppink}{Exercices}{type examen + corrigés}\end{minipage}\hfill
  \begin{minipage}[t]{0.235\linewidth}\poptile{popgreen}{Fiche bilan}{l'essentiel à retenir}\end{minipage}\par}

% Sommaire
\newtcolorbox{sommaire}{enhanced,colback=white,colframe=popink,boxrule=2.2pt,arc=10pt,
  drop shadow=popink,left=18pt,right=18pt,top=13pt,bottom=13pt,before skip=6pt,after skip=20pt,
  before upper={\poppill{Sommaire}{poppurple}{white}\par\vspace{9pt}\popsemi\fontsize{10.6}{14.5}\selectfont\color{popink}}}

% Signature de fin de cours — poussée tout en bas de la dernière page
\newcommand{\findecours}{%
  \par\vfill\nopagebreak
  \begin{center}\popsquiggle{poporange}\end{center}\vspace{4pt}
  \signatureonbuch
}
\AtBeginDocument{\def\llbracket{\mathopen{[\mkern-3.5mu[}}\def\rrbracket{\mathclose{]\mkern-3.5mu]}}} % intervalles d'entiers (glyphe ⟦ absent de la police)
% Glyphes absents de Fira Math : redéfinis à partir de glyphes disponibles
\AtBeginDocument{%
  \def\setminus{\mathbin{\backslash}}\let\smallsetminus\setminus
  \def\vdots{\mathord{\vbox{\baselineskip3.2pt\lineskiplimit0pt\kern4pt\hbox{.}\hbox{.}\hbox{.}}}}%
  \def\ddots{\mathinner{\mkern1mu\raise6.5pt\hbox{.}\mkern2mu\raise3.6pt\hbox{.}\mkern2mu\raise0.7pt\hbox{.}\mkern1mu}}%
  \def\triangle{\mathord{\scalebox{0.9}{$\symup{\Delta}$}}}%
  \def\nabla{\mathord{\raisebox{\depth}{\scalebox{1}[-1]{$\symup{\Delta}$}}}}%
  \def\checkmark{\popcheck{popdark}}%
  \def\star{\mathbin{\ast}}\def\bigstar{\mathord{\ast}}%
  \def\diamond{\mathbin{\scalebox{0.6}{$\Diamond$}}}%
  \def\bigcup{\mathop{\mathchoice{\raisebox{-0.3em}{\scalebox{1.7}{$\cup$}}}{\scalebox{1.25}{$\cup$}}{\cup}{\cup}}\displaylimits}%
  \def\bigcap{\mathop{\mathchoice{\raisebox{-0.3em}{\scalebox{1.7}{$\cap$}}}{\scalebox{1.25}{$\cap$}}{\cap}{\cap}}\displaylimits}%
  \def\lesseqqgtr{\mathrel{\lessgtr}}\def\bigoplus{\mathop{\oplus}\displaylimits}%
}
\providecommand{\ssi}{si et seulement si} % abréviation fréquente
\AtBeginDocument{\providecommand{\wideparen}{\overparen}} % arc de cercle

%% FIGFIX-BEGIN v4 — correctifs globaux des figures (docs/onbuch-generation/tools/figfix)
\usepackage{adjustbox}\usepackage{etoolbox}
% toute figure TikZ/pgfplots plus large que la ligne est réduite à la largeur disponible
\BeforeBeginEnvironment{tikzpicture}{\begin{adjustbox}{max width=\linewidth}}
\AfterEndEnvironment{tikzpicture}{\end{adjustbox}}
% légendes pgfplots lisibles par-dessus les courbes
\pgfplotsset{every axis/.append style={legend style={fill=white,fill opacity=0.92,text opacity=1,draw=black!15,inner sep=3pt}}}
% étiquettes d'arêtes (syntaxe edge["..."]) avec fond blanc
\tikzset{every edge quotes/.append style={fill=white,inner sep=1.2pt}}
%% FIGFIX-END
\begin{document}

\pagegarde{Le Proche et le Moyen Orient : zones de tension sans fin}{Histoire}{Tle Tech}{Histoire – classe de Terminale technique}{N03}{2 séances (fiche de progression harmonisée)}{Cette leçon analyse les racines et les ressorts des conflits qui déchirent le Proche et le Moyen-Orient depuis 1948 : la question israélo-palestinienne, véritable poudrière régionale, et les guerres du Moyen-Orient et du golfe Persique liées au pétrole, aux rivalités de puissance et aux fractures religieuses. Elle montre aussi comment ces tensions lointaines affectent concrètement la vie des Camerounais (prix du carburant, diplomatie, sécurité).
\vspace{4pt}
\begin{multicols}{2}\small
\begin{itemize}
\item À la fin de cette leçon, je sais situer sur une carte le Proche-Orient, le Moyen-Orient et le golfe Persique, ainsi que les principaux États et lieux de conflit
\item À la fin de cette leçon, je sais expliquer les origines du problème israélo-palestinien (sionisme, déclaration Balfour, partage de 1947)
\item À la fin de cette leçon, je sais retracer les grandes étapes du conflit israélo-arabe (guerres de 1948, 1956, 1967, 1973) et les tentatives de paix (Camp David, Oslo)
\item À la fin de cette leçon, je sais analyser les causes des conflits du golfe Persique (guerre Iran-Irak, guerre du Koweït, intervention de 2003)
\item À la fin de cette leçon, je sais interpréter des documents (cartes, photographies, données chiffrées sur le pétrole et les réfugiés) pour établir une notion géographique et historique
\item À la fin de cette leçon, je sais construire un croquis ou un diagramme simple à partir de données (production pétrolière, chronologie des conflits)
\end{itemize}\end{multicols}}

\begin{sommaire}
\begin{multicols}{2}
\begin{enumerate}
\item Introduction : situer l'espace et poser le problème
\item Les origines du problème israélo-palestinien
\item Une poudrière : guerres, Intifadas et paix introuvable
\item Les conflits du Moyen-Orient et du golfe Persique
\item Pourquoi des tensions « sans fin » ? Répercussions jusqu'au Cameroun
\item TP guidés : construire pour comprendre
\item Activité d'intégration
\item Méthodes et exercices résolus
\item Exercices
\item Corrigés des exercices
\item Fiche bilan
\end{enumerate}
\end{multicols}
\end{sommaire}

\begin{objectifs}
\begin{itemize}
\item À la fin de cette leçon, je sais situer sur une carte le Proche-Orient, le Moyen-Orient et le golfe Persique, ainsi que les principaux États et lieux de conflit
\item À la fin de cette leçon, je sais expliquer les origines du problème israélo-palestinien (sionisme, déclaration Balfour, partage de 1947)
\item À la fin de cette leçon, je sais retracer les grandes étapes du conflit israélo-arabe (guerres de 1948, 1956, 1967, 1973) et les tentatives de paix (Camp David, Oslo)
\item À la fin de cette leçon, je sais analyser les causes des conflits du golfe Persique (guerre Iran-Irak, guerre du Koweït, intervention de 2003)
\item À la fin de cette leçon, je sais interpréter des documents (cartes, photographies, données chiffrées sur le pétrole et les réfugiés) pour établir une notion géographique et historique
\item À la fin de cette leçon, je sais construire un croquis ou un diagramme simple à partir de données (production pétrolière, chronologie des conflits)
\item À la fin de cette leçon, je sais rédiger et justifier une conclusion argumentée sur la question : « Pourquoi le Proche et le Moyen-Orient sont-ils des zones de tension sans fin ? »
\item À la fin de cette leçon, je sais relier ces conflits à des situations concrètes de la vie courante au Cameroun (prix des carburants, position diplomatique du Cameroun à l'ONU)
\end{itemize}
\end{objectifs}

\begin{prerequis}
\begin{itemize}
\item La décolonisation et la fin des mandats au Proche-Orient (notions de 3e/2de)
\item La Première Guerre mondiale et le partage de l'Empire ottoman (accords Sykes-Picot, 1916)
\item La Seconde Guerre mondiale et la Shoah, facteur de la création d'Israël
\item La guerre froide et l'intervention des grandes puissances dans le Tiers-Monde
\item Notions de base : ONU, OPEP, nationalisme, mandat international
\end{itemize}
\end{prerequis}

\newpage

\coursec{Introduction : situer l'espace et poser le problème}

\courssub{Une situation d'accroche : le Cameroun au cœur des répercussions mondiales}

En juin 2024, les élèves d'un lycée de Yaoundé débattent : pourquoi le prix du carburant à la station Total monte-t-il chaque fois qu'une guerre éclate à Gaza ou dans le golfe Persique, alors que ces lieux sont à plus de \SI{4000}{km} du Cameroun ? Pourquoi les chaînes de télévision camerounaises (CRTV, Canal~2) consacrent-elles chaque soir des reportages à une région si lointaine ? Et pourquoi le Cameroun, comme la plupart des pays africains, est-il régulièrement appelé à voter à l'ONU sur le conflit israélo-palestinien ? Comprendre ces tensions permanentes, c'est comprendre une partie essentielle de l'actualité mondiale et de ses répercussions sur notre quotidien.

Cette interrogation légitime pose d'emblée la \cle{problématique} de notre leçon : \textbf{pourquoi le Proche et le Moyen-Orient constituent-ils une « zone de tension sans fin » depuis 1948, et quelles en sont les conséquences planétaires, y compris pour le Cameroun ?} Pour y répondre, nous devons d'abord poser les cadres géographiques, géologiques et humains de cet espace complexe.

\courssub{Délimitation géographique : deux espaces imbriqués à ne pas confondre}

\begin{definition}[Proche-Orient et Moyen-Orient : deux concepts géopolitiques distincts]
Le \textbf{Proche-Orient} désigne l'espace historique et stratégique centré sur le Levant méditerranéen : Israël, Palestine, Liban, Syrie, Jordanie, Égypte (Sinaï), ainsi que la Turquie et Chypre parfois. C'est le berceau des trois monothéismes et le théâtre du conflit israélo-palestinien.
Le \textbf{Moyen-Orient} est un espace plus vaste, élargi vers l'est : il intègre le Proche-Orient et y adjoint l'Irak, l'Iran, la péninsule Arabique (Arabie saoudite, Yémen, Oman, Émirats arabes unis, Qatar, Bahreïn, Koweït), et parfois l'Afghanistan et le Pakistan. Le \textbf{Golfe Persique} (ou golfe arabo-persique) en constitue le cœur énergétique.
\emph{Attention :} l'expression « Proche et Moyen-Orient » (PMO) est souvent utilisée pour désigner l'ensemble de cette grande région stratégique.
\end{definition}

\begin{attention}[Erreur fréquente : confondre les deux espaces ou essentialiser les populations]
Ne pas confondre \cle{Proche-Orient} (foyer du conflit israélo-arabe) et \cle{Moyen-Orient} (espace élargi aux monarchies du Golfe et à l'Iran).
Ne pas placer l'\textbf{Iran} dans le « monde arabe » : l'Iran est un pays \textbf{perse} (indo-européen), majoritairement \textbf{chiite}, distinct par sa langue (le farsi), son histoire et sa culture des pays arabes sunnites qui l'entourent. De même, la Turquie est turque (sunnite, laïque historiquement) et non arabe. Les Kurdes, répartis sur quatre États (Turquie, Iran, Irak, Syrie), forment une nation sans État, ni arabe, ni perse, ni turque.
\end{attention}

\begin{popfigure}
\begin{tikzpicture}[scale=0.85, every node/.style={font=\small}]
% Coastlines simplified (polygon approximation)
\draw[popblue!50, line width=1.2pt, fill=popblue!5] 
(0,0) -- (10,0) -- (10,1) -- (11,2) -- (11,4) -- (10,5) -- (10,7) -- (9,8) -- (7,8.5) -- (5,8) -- (4,7) -- (3,6) -- (2,5) -- (1,4) -- (0,3) -- cycle; % Mediterranean / Red Sea / Arabian Sea blob
% Land masses (simplified polygons)
% Turkey (top left)
\draw[popgreen!60!popmuted, fill=popgreen!10] (1,7) -- (3,8) -- (5,8) -- (4,7) -- cycle;
\node[popdark] at (3,7.5) {\textbf{Turquie}};
% Levant coast
\draw[popgreen!60!popmuted, fill=popgreen!10] (3,5) -- (5,5.5) -- (5,4) -- (4,3.5) -- (3,4) -- cycle;
\node[popdark] at (4,4.5) {\textbf{Syrie}};
\node[popdark] at (4.5,3.8) {\textbf{Liban}};
\node[popdark] at (4.2,3.2) {\textbf{Israël}};
\node[popdark, poppink] at (3.8,3.0) {\textbf{Palestine}};
% Jordan
\draw[popgreen!60!popmuted, fill=popgreen!10] (4,3) -- (6,3.5) -- (6,2) -- (4,2) -- cycle;
\node[popdark] at (5,2.5) {\textbf{Jordanie}};
% Egypt / Sinai
\draw[popgreen!60!popmuted, fill=popgreen!10] (2,1) -- (4,2) -- (4,0.5) -- (2,0) -- cycle;
\node[popdark] at (3,1) {\textbf{Égypte}};
\node[popdark, rotate=90] at (1.5,1.5) {\textbf{Sinaï}};
% Iraq
\draw[poporange!70!popmuted, fill=poporange!10] (5,4) -- (8,5) -- (8,3) -- (6,2.5) -- cycle;
\node[popdark] at (7,4) {\textbf{Irak}};
% Iran
\draw[poppurple!70!popmuted, fill=poppurple!10] (7,5) -- (11,6) -- (11,3) -- (8,2.5) -- cycle;
\node[popdark] at (9.5,4.5) {\textbf{Iran}};
% Arabian Peninsula
\draw[popgold!70!popmuted, fill=popgold!10] (5,2) -- (10,2.5) -- (10,0) -- (4,0) -- cycle;
\node[popdark] at (7.5,1.2) {\textbf{Péninsule Arabique}};
\node[popdark] at (6.5,0.5) {\textbf{Arabie Saoudite}};
\node[popdark] at (9,0.5) {\textbf{EAU / Qatar}};
\node[popdark] at (8.5,2) {\textbf{Koweït}};
% Strategic Straits
% Suez
\draw[poppink, line width=2pt, ->] (2.5,1.5) -- (2.5,2.5);
\node[poppink, above] at (2.5,2.6) {\textbf{Détroit de Suez}};
% Bab el-Mandeb
\draw[poppink, line width=2pt, ->] (4.5,0.2) -- (5.5,-0.2);
\node[poppink, below right] at (5.5,-0.2) {\textbf{Bab el-Mandeb}};
% Ormuz
\draw[poppink, line width=2pt, ->] (9.5,2.5) -- (10.5,2.5);
\node[poppink, right] at (10.5,2.5) {\textbf{Détroit d'Ormuz}};
% Oil fields (symbols)
\foreach \x/\y in {7.5,1.5 / 8.5,1.8 / 9,2.2 / 8,3.5 / 7.5,4.5 / 7,3 / 8.5,3.5} {
    \fill[poporange] (\x,\y) circle (2.5pt);
    \draw[poporange!80!popdark] (\x,\y) circle (4pt);
}
\node[poporange!80!popdark, anchor=west] at (10.5,4) {\textbullet\ \textbf{Gisements pétroliers majeurs}};
% North Arrow
\node[popdark] at (0.5,8) {\textbf{N}};
\draw[->, popdark, line width=1pt] (0.5,7.7) -- (0.5,7.2);
% Scale bar
\draw[popdark, line width=1pt] (0.2,0.2) -- (1.2,0.2);
\node[popdark, below] at (0.7,0.1) {$\approx$ \SI{1000}{km}};
\end{tikzpicture}
\legende{Carte schématique du Proche et du Moyen-Orient : États, détroits stratégiques (Suez, Bab el-Mandeb, Ormuz) et principaux gisements pétroliers (points orange). Les couleurs distinguent le Proche-Orient (tons verts) du Moyen-Orient élargi (tons orange/violon pour Irak/Iran, or pour Péninsule Arabique).}
\end{popfigure}

\courssub{Un carrefour stratégique : la jonction de trois continents et des détroits vitaux}

L'espace PMO est un \cle{carrefour stratégique} par excellence. Il assure la jonction physique entre l'\textbf{Afrique}, l'\textbf{Asie} et l'\textbf{Europe}. Cette position centrale commande trois \cle{détroits} maritimes indispensables au commerce mondial et aux flux énergétiques :

\begin{itemize}
    \item Le \textbf{canal de Suez} (Égypte) et le \textbf{détroit de Bab el-Mandeb} (entre Yémen et Djibouti) relient la mer Méditerranée à l'océan Indien via la mer Rouge. Environ \textbf{12 \% du commerce maritime mondial} y transite.
    \item Le \textbf{détroit d'Ormuz} (entre Iran et Oman) est l'unique sortie du golfe Persique vers l'océan Indien.
\end{itemize}

\begin{savaistu}[Le détroit d'Ormuz : la jugulaire énergétique de la planète]
Chaque jour, environ \SIrange{20}{21}{\million} de barils de pétrole brut et de produits raffinés (soit \textbf{$\sim$ 20 \% de la consommation mondiale quotidienne}) transitent par le détroit d'Ormuz, large de seulement \SI{39}{km} au plus étroit (\SI{2}{km} dans le chenal de navigation). Toute tension militaire (minage, blocus, affrontement naval) y provoque une hausse immédiate des cours du brut sur les marchés de Londres (Brent) et New York (WTI), répercutée à la pompe à Yaoundé, Douala ou Paris dans les semaines qui suivent.
\end{savaistu}

\courssub{Un espace aux richesses convoitées : la rente pétrolière}

Le sous-sol du PMO recèle \textbf{environ 50 \% des réserves mondiales de pétrole prouvées} (et près de 40 \% pour le gaz naturel). Cette concentration exceptionnelle en fait l'épicentre de la géopolitique de l'énergie depuis la Première Guerre mondiale.

\begin{popfigure}
\begin{tikzpicture}
\begin{axis}[
    xbar, width=\linewidth, height=7cm,
    xmin=0, xmax=60, xlabel={Part des réserves mondiales (\%)},
    symbolic y coords={Autres,Asie-Pacifique,Europe/Ex-URSS,Afrique,Amérique Latine,Amérique du Nord,Moyen-Orient},
    ytick=data, bar width=10pt,
    nodes near coords, tick label style={font=\footnotesize, popdark},
    title={Répartition des réserves mondiales de pétrole prouvées (estimation 2023)},
    title style={font=\normalsize\bfseries, popdark}
]
\addplot[fill=popblue!70] coordinates {
    (48,Moyen-Orient) (14,Amérique du Nord) (8,Amérique Latine) (8,Afrique) (10,Europe/Ex-URSS) (7,Asie-Pacifique) (5,Autres)
};
\end{axis}
\end{tikzpicture}
\legende{Diagramme circulaire (camembert polaire) : le Moyen-Orient concentre près de la moitié des réserves mondiales prouvées (source : BP Statistical Review / OPEP 2023). Cette rente structure les économies locales (États-rentes) et les stratégies des grandes puissances.}
\end{popfigure}

\courssub{Un espace de fractures multiples : religieuses, ethniques, politiques}

Au-delà de la géographie et de la géologie, le PMO est un \cle{mosaïque} humain fragmenté par des lignes de faille profondes qui alimentent les conflits :

\begin{center}
\begin{tabularx}{\linewidth}{|l|X|X|}
\hline
\rowcolor{popblueL}
\textbf{Type de fracture} & \textbf{Principales lignes de clivage} & \textbf{Exemples de tensions} \\
\hline
\textbf{Religieuse} & \textbf{Islam sunnite} (majoritaire, Arabie saoudite, Égypte, Turquie) vs \textbf{Islam chiite} (Iran, Irak, Hezbollah au Liban, Houthis au Yémen). \textbf{Judaïsme} (Israël) vs \textbf{Islam}. Minorités chrétiennes (Coptes, Maronites, Assyriens) vulnérables. & Guerre Iran/Irak (1980-88), guerre civile syrienne (axe chiite vs sunnite), conflit israélo-palestinien. \\
\hline
\textbf{Ethnique / Nationale} & \textbf{Arabes} (majorité) vs \textbf{Juifs} (Israël) vs \textbf{Perses} (Iran) vs \textbf{Turcs} (Turquie) vs \textbf{Kurdes} (sans État, 30-40 M). & Revendications kurdes (PKK, PYD), politique turque anti-kurde, identité nationale irakienne/syrienne fracturée. \\
\hline
\textbf{Politique / Idéologique} & \textbf{Monarchies conservatrices} (Golfe) vs \textbf{Républiques nationalistes/laïques} (ex-Baas Syrie/Irak) vs \textbf{République islamique} (Iran) vs \textbf{État juif/démocratique} (Israël) vs \textbf{Mouvements islamistes} (Frères musulmans, Hamas, Daech, Al-Qaïda). & Printemps arabes (2011), guerre froide régionale Arabie/Iran, rivalité Qatar/Arabie (2017-21). \\
\hline
\end{tabularx}
\end{center}

Ces fractures ne s'alignent pas : un Kurde peut être sunnite (majorité) ou alévi/chiite ; un Arabe peut être chrétien, sunnite ou chiite ; un Israélien peut être laïc, religieux sioniste, ou arabe israélien (20 \% de la population, majoritairement musulmans sunnites). Cette complexité rend toute solution « binaire » illusoire.

\courssub{Méthode : lire et légender une carte géopolitique}

\begin{methode}[Analyser une carte géopolitique du PMO en 5 étapes]
Une carte n'est pas une illustration décorative, c'est un \textbf{document source} à interroger.
\begin{enumerate}
    \item \textbf{Identifier le titre, la date et la source} : quelle thèse la carte soutient-elle ? (ex. : « Plan de partage 1947 », « Frontières 1967 », « Zones d'influence 2024 »).
    \item \textbf{Lire la légende ordonnée} : distinguer les \emph{aires} (couleurs de fond : États, zones contrôlées), les \emph{linéaires} (frontières reconnues, lignes d'armistice, murs, barrières), les \emph{ponctuels} (capitales, lieux saints, bases militaires, gisements, détroits).
    \item \textbf{Repérer l'orientation et l'échelle} : le Nord n'est pas toujours en haut ; l'échelle permet de mesurer les distances réelles (ex. : largeur d'Israël au niveau de Tel-Aviv $\approx$ \SI{15}{km}).
    \item \textbf{Croiser les couches d'information} : superposer mentalement la carte des \emph{ressources} (pétrole, eau, nappes phréatiques), la carte des \emph{peuples} (répartition confessionnelle/ethnique) et la carte des \emph{pouvoirs} (contrôle effectif au sol).
    \item \textbf{Formuler une synthèse écrite} : « Cette carte montre que... ce qui explique que... » (lien spatial $\rightarrow$ causalité géopolitique).
\end{enumerate}
\end{methode}

\begin{exoresolu}[Application de la méthode à la carte ci-dessus]
\textbf{Énoncé :} À partir de la carte schématique (figure 1), expliquer pourquoi le contrôle du détroit d'Ormuz est un enjeu vital pour l'Iran et pour les monarchies du Golfe.
\textbf{Solution détaillée :}
\begin{enumerate}
    \item \emph{Localisation :} Le détroit d'Ormuz (flèche rose à droite) sépare l'Iran (zone violette) de la péninsule Arabique (zone or). Il est l'unique sortie maritime du golfe Persique.
    \item \emph{Ressources :} Les points orange (gisements) sont massivement concentrés au nord-est de la péninsule Arabique (Arabie saoudite, Koweït, EAU, Qatar, Irak sud) et côté iranien. 
    \item \emph{Flux :} Tout le pétrole exporté par tankers depuis ces terminaux \emph{doit} passer par ce goulet d'étranglement.
    \item \emph{Conséquence géopolitique :} L'Iran (puissance riveraine nord) peut menacer de le fermer (missiles côtiers, mines, vedettes rapides) pour faire pression sur l'Occident. Les monarchies (riveraines sud) dépendent à 100 \% de ce passage pour leur revenu principal. D'où la présence militaire américaine (5e flotte à Bahreïn) pour garantir la liberté de navigation.
\end{enumerate}
\end{exoresolu}

\courssub{Problématique et annonce du plan}

Au terme de cette introduction, les traits structurels de la région apparaissent clairement :
\begin{itemize}
    \item Une \textbf{géographie contrainte} (détroits, désert, eau rare) qui concentre les populations et les flux.
    \item Une \textbf{géologie exceptionnelle} (pétrole/gaz) qui attire les convoitises des puissances extérieures (anciennes métropoles coloniales, USA, URSS/Russie, Chine, UE).
    \item Des \textbf{fractures internes} (religieuses, ethniques, politiques) surdéterminées par l'héritage colonial (accords Sykes-Picot 1916, mandats, promesse Balfour 1917) et la création de l'État d'Israël (1948).
\end{itemize}

C'est l'interaction de ces trois niveaux — \textbf{local} (revendicatives nationales), \textbf{régional} (luttes d'hégémonie Iran/Arabie/Turquie/Israël), \textbf{international} (sécurité énergétique, lutte d'influence) — qui verrouille les conflits et justifie l'expression de « zone de tension sans fin ».

\textbf{Annonce du plan :}
\begin{enumerate}
    \item \textbf{Le problème israélo-palestinien : une poudrière au Proche-Orient} (origines, guerres, Intifadas, processus de paix avortés, situation actuelle).
    \item \textbf{Les conflits du Moyen-Orient et du golfe Persique} (guerres Irak/Iran, du Golfe, Irak 2003, printemps arabes, guerres par procuration Syrie/Yémen, nucléaire iranien).
\end{enumerate}

\coursec{Les origines du problème israélo-palestinien}

\medskip
\textit{Dans la section précédente, nous avons situ\'e l'espace g\'eographique du Proche-Orient et pos\'e la question centrale : comment ce territoire, carrefour de civilisations, est-il devenu le foyer d'un conflit persistant qui structure les relations internationales depuis plus d'un si\`ecle ? Pour le comprendre, il faut remonter aux racines historiques du diff\'erend isra\'elo-palestinien, au moment o\`u les nationalismes juif et arabe entrent en collision sur une m\^eme terre, la Palestine, sous le regard des puissances europ\'eennes.}

\courssub{La Palestine sous mandat britannique (1920-1948) : un h\'eritage ottoman et des promesses contradictoires}

Apr\`es la d\'efait de l'Empire ottoman lors de la Premi\`ere Guerre mondiale, la Soci\'et\'e des Nations (SDN) confie en 1920 \`a la Grande-Bretagne un \emph{mandat} sur la Palestine. Ce syst\`eme, pr\'evu par l'article 22 du Pacte de la SDN, vise th\'eoriquement \`a pr\'eparer les territoires \`a l'ind\'ependance. Mais le mandat palestinien comporte une particularit\'e majeure : son pr\'eambule int\`egre la \emph{D\'eclaration Balfour} (1917), par laquelle la puissance mandataire s'engage \`a favoriser l'\'etablissement en Palestine d'un \guillemotleft foyer national juif \guillemotright, tout en pr\'eservant les droits civils et religieux des communaut\'es non-juives existantes.

\begin{definition}[Foyer national juif]
Expression employ\'ee dans la D\'eclaration Balfour (1917) et reprise dans le mandat britannique. Elle d\'esigne un centre de vie nationale, culturelle et d\'emographique pour le peuple juif en Palestine, sans pr\'eciser initialement s'il s'agit d'un \'Etat souverain. L'ambigu\"{\i}t\'e du terme (foyer vs \'Etat) alimente les interpr\'etations divergentes et les tensions.
\end{definition}

La population de la Palestine mandataire est alors \`a majorit\'e arabe (musulmane et chr\'etienne), avec une minorit\'e juive ancienne (\emph{Yichouv}) repr\'esentant environ 10 \% en 1917. L'administration britannique doit g\'erer une double promesse contradictoire : l'ind\'ependance arabe (promesses de 1915-1916 \`a Hussein de La Mecque) et le foyer national juif (Balfour 1917). Cette contradiction structurelle rend la gouvernance impossible et alimente la m\'efiance mutuelle.

\courssub{Le sionisme : de l'id\'ee \`a la r\'ealisation politique}

Le sionisme politique moderne na\^it \`a la fin du XIXe si\`ecle en r\'eponse \`a l'antis\'emitisme europ\'een (affaire Dreyfus, pogroms en Russie) et \`a l'\'echec de l'assimilation. Son fondateur, \textbf{Theodor Herzl}, publie \emph{L'\'Etat des Juifs} (1896) et organise le Premier Congr\`es sioniste \`a B\^ale (1897). Le programme de B\^ale d\'efinit l'objectif : \guillemotleft cr\'eer pour le peuple juif un foyer assur\'e en Palestine, garanti par le droit public \guillemotright.

\begin{definition}[Sionisme]
Mouvement nationaliste juif n\'e \`a la fin du XIXe si\`ecle (Th\'eodor Herzl, \emph{L'\'Etat des Juifs}, 1896) qui revendique le droit \`a l'autod\'etermination du peuple juif sur sa terre historique, la Palestine (Eretz Israel). Il se d\'ecline en courants divers (politique, socialiste, r\'evisionniste, religieux) mais partage l'objectif d'un retour massif (\emph{Alyah}) et d'une souverainet\'e juive.
\end{definition}

Le mouvement s'organise : achat de terres (via le Fonds national juif), cr\'eation d'institutions pr\'e-\'etatiques (Histadrout, Agence juive, Haganah), relance de l'h\'ebreu comme langue vivante. L'immigration juive (\emph{Alyah}) s'acc\'el\`ere par vagues : troisi\`eme alyah (1919-1923), quatri\`eme (1924-1929), cinqui\`eme (1929-1939) fuyant le nazisme. La population juive passe de 56 000 (1918) \`a environ 600 000 (1947), modifiant profond\'ement l'\'equilibre d\'emographique.

\begin{savaistu}[Le congr\`es de B\^ale, acte de naissance]
Le Premier Congr\`es sioniste se tient \`a B\^ale (Suisse) du 29 au 31 ao\^ut 1897. Herzl note dans son journal : \guillemotleft \`A B\^ale, j'ai fond\'e l'\'Etat juif [...] Peut-\^etre dans cinq ans, certainement dans cinquante, tout le monde le reconna\^itra. \guillemotright La pr\'ediction se r\'ealise presque jour pour jour 50 ans plus tard (plan de partage ONU, novembre 1947).
\end{savaistu}

\courssub{L'immigration, la Shoah et la mont\'ee des violences (1920-1947)}

L'afflux de r\'efugi\'es juifs, amplifi\'e par la mont\'ee du nazisme (1933) et la fermeture des portes occidentales (conf\'erence d'\'Evian, 1938), puis par la Shoah (1941-1945), transforme le projet sioniste en urgence humanitaire et politique. La population arabe palestinienne per\c{c}oit cette immigration comme une colonisation d\'eposs\'edant les fellahs de leurs terres et mena\c{c}ant leur avenir national.

Les tensions \'eclatent en r\'evoltes ouvertes :
\begin{itemize}
    \item \textbf{1920-1921 :} \'Emeutes de J\'erusalem et de Jaffa.
    \item \textbf{1929 :} H\'ebron et J\'erusalem, affrontements autour du Mur des Lamentations / Esplanade des Mosqu\'ees.
    \item \textbf{1936-1939 :} \textbf{Grande R\'evolte arabe}. Gr\`eve g\'en\'erale, gu\'errilla, r\'epression britannique (lois d'urgence, pendaisons, exil des chefs nationalistes comme le Mufti de J\'erusalem). La commission Peel (1937) propose le premier plan de partage (rejet\'e par les deux camps). Le \emph{Livre blanc} britannique de 1939 limite drastiquement l'immigration juive (75 000 sur 5 ans) et promet un \'Etat palestinien unitaire en 10 ans, trahissant selon les sionistes la d\'eclaration Balfour.
\end{itemize}

\begin{attention}[Pi\`ege chronologique fr\'equent]
Ne pas confondre la \textbf{cr\'eation de l'\'Etat d'Isra\"el (1948)} avec la \textbf{Seconde Guerre mondiale (1939-1945)}. L'\'Etat r\'esulte d'un processus long : sionisme (1897), d\'eclaration Balfour (1917), mandat britannique, r\'esolution ONU 181 (1947). La Shoah a \emph{acc\'el\'er\'e} la l\'egitim\'e internationale et l'urgence d\'emographique, mais n'est pas la \emph{cause unique} ni l'acte juridique fondateur.
\end{attention}

\courssub{Le plan de partage de l'ONU (r\'esolution 181, 29 novembre 1947)}

Face \`a l'impasse (insurrection sioniste Irgoun/Lehi, \'epuisement britannique, contexte de guerre froide naissante), la Grande-Bretagne renvoie le mandat \`a l'ONU (f\'evrier 1947). L'UNSCOP (commission d'enqu\^ete) propose \`a la majorit\'e le partage. L'Assembl\'ee g\'en\'erale vote la \textbf{r\'esolution 181 (II)} le 29 novembre 1947.

\begin{exemplebox}[Le chiffre cl\'e du plan de partage]
La r\'esolution 181 attribue \`a l'\'Etat juif \textbf{environ 55 \% du territoire} de la Palestine mandataire (dont le N\'ege\`v d\'esertique), alors que la population juive ne repr\'esente qu'\textbf{environ un tiers (32 \%)} des habitants et ne poss\`ede que \textbf{6 \`a 7 \% des terres} priv\'ees. L'\'Etat arabe re\c{c}oit 45 \% (terres plus arables), J\'erusalem et Bethl\'eem forment un \emph{corpus separatum} sous administration internationale. Ce d\'es\'equilibre per\c{c}u explique le refus cat\'egorique des dirigeants arabes (Ligue arabe, Comit\'e supr\^eme arabe) qui y voient une injustice flagrante et une violation du droit des peuples \`a disposer d'eux-m\^emes. Le Yichouv (direction sioniste) l'accepte comme une base juridique internationale, tout en se pr\'eparant militairement \`a son extension.
\end{exemplebox}

\begin{savaistu}[Le vote de la r\'esolution 181]
Le 29 novembre 1947, l'Assembl\'ee g\'en\'erale de l'ONU vote le plan de partage : \textbf{33 voix pour} (USA, URSS, France, pays d'Am\'erique latine...), \textbf{13 contre} (pays arabes, Inde, Pakistan, Gr\`ece, Cuba...), \textbf{10 abstentions} (Grande-Bretagne, Chine, pays d'Europe de l'Ouest...). La majorit\'e des deux tiers requise est atteinte (33/56). La r\'esolution n'est pas contraignante (chapitre VI de la Charte), mais conf\`ere une l\'egitim\'e internationale \`a la cr\'eation de l'\'Etat juif.
\end{savaistu}

\courssub{La proclamation de l'\'Etat d'Isra\"el et la premi\`ere guerre isra\'elo-arabe (1948-1949)}

Le mandat britannique prend fin le 14 mai 1948 \`a minuit. Le \textbf{14 mai 1948 (5 Iyar 5708)}, \textbf{David Ben Gourion}, chef de l'Agence juive, proclame l'ind\'ependance de l'\textbf{\'Etat d'Isra\"el} au mus\'ee de Tel-Aviv. La D\'eclaration d'ind\'ependance cite la r\'esolution 181, le droit naturel et historique, et la Shoah. Les \'Etats-Unis (Truman) et l'URSS (Staline) reconnaissent imm\'ediatement le nouvel \'Etat.

Le lendemain, 15 mai 1948, les arm\'ees de cinq \'Etats arabes (\'Egypte, Transjordanie, Syrie, Liban, Irak) envahissent le territoire : c'est la \textbf{premi\`ere guerre isra\'elo-arabe} (appel\'ee \emph{Guerre d'ind\'ependance} par les Isra\'eliens, \emph{Nakba} par les Palestiniens). Malgr\'e l'inf\'eriorit\'e num\'erique initiale, les forces juives (Haganah, Palmah, Irgoun, Lehi unifi\'ees en \emph{Tsahal} le 26 mai) l'emportent gr\^ace \`a une meilleure organisation, un commandement unifi\'e, des achats d'armes (Tch\'ecoslovaquie) et la motivation existentielle.

\begin{popfigure}
\begin{tikzpicture}[scale=0.75, every node/.style={font=\small}]
    % Coastline approximation
    \draw[popblue!50!popdark, line width=1.5pt] (0,0) .. controls (1,1) and (3,2) .. (4,3.5) .. controls (5,4) and (6,4.5) .. (7,5) .. controls (8,5.5) and (9,5) .. (10,4) .. controls (11,3) and (12,2) .. (12,0);
    % Jordan River / Dead Sea / Lake Tiberias
    \draw[popblue, line width=1pt] (5,1.5) .. controls (5.5,2.5) .. (5.5,3.5); % Jordan river approx
    \fill[popblue!20] (5,1) ellipse ({0.8} and {0.4}); % Dead Sea
    \fill[popblue!20] (5.5,3.2) ellipse ({0.5} and {0.3}); % Tiberias
    % UN Partition Plan 1947 (Hatched areas)
    % Jewish State (3 zones)
    \pattern[pattern=north east lines, pattern color=popblue!60!popdark] (0.5,0.5) .. controls (2,1.5) .. (3,2.5) -- (4,2) -- (3,1) -- cycle; % North coastal
    \pattern[pattern=north east lines, pattern color=popblue!60!popdark] (6,0.5) .. controls (7,1.5) .. (7,2.5) -- (8,2) -- (7,1) -- cycle; % Center coastal
    \pattern[pattern=north east lines, pattern color=popblue!60!popdark] (4,0.5) -- (11,0.5) -- (11,1.5) .. controls (8,2) .. (4,1.5) -- cycle; % Negev
    % Arab State (3 zones) - different hatch
    \pattern[pattern=north west lines, pattern color=poporange!70!popdark] (0,3.5) -- (3,2.5) -- (3,1.5) -- (0.5,0.5) -- (0,1) -- cycle; % Galilee
    \pattern[pattern=north west lines, pattern color=poporange!70!popdark] (3,2.5) -- (6,3) -- (7,2.5) -- (4,2) -- cycle; % Central hills (West Bank approx)
    \pattern[pattern=north west lines, pattern color=poporange!70!popdark] (7,2.5) -- (10,4) -- (11,3) -- (8,2) -- cycle; % South (Gaza strip area extended)
    % Jerusalem Corpus Separatum
    \fill[poppurple!30] (5.5,2.5) circle (0.4);
    \node[poppurple!80!popdark] at (5.5,2.5) {\textbf{J\'erusalem}};
    % Green Line 1949 (Armistice) - Solid thick line
    \draw[popgreen!70!popdark, line width=2.5pt, dashed] 
        (0,3.5) .. controls (2,3) .. (3.5,2.8) .. controls (4.5,2.5) .. (5.5,2.2) .. controls (6.5,1.8) .. (7.5,1.5) .. controls (9,1) .. (11,1);
    % Legends - using inline tikz for patterns
    \node[anchor=west, align=left] at (12.5, 4.5) {\textbf{L\'egende :}};
    \node[anchor=west, align=left] at (12.5, 4.0) {\tikz\fill[pattern=north east lines, pattern color=popblue!60!popdark] (0,0) rectangle (0.5,0.3); Plan ONU 1947 : \'Etat juif (55 \%)};
    \node[anchor=west, align=left] at (12.5, 3.5) {\tikz\fill[pattern=north west lines, pattern color=poporange!70!popdark] (0,0) rectangle (0.5,0.3); Plan ONU 1947 : \'Etat arabe (45 \%)};
    \node[anchor=west, align=left] at (12.5, 3.0) {\color{poppurple!80!popdark}\rule{0.5cm}{0.3cm} J\'erusalem \emph{Corpus Separatum}};
    \node[anchor=west, align=left] at (12.5, 2.5) {\color{popgreen!70!popdark}\rule{1cm}{2pt} Ligne verte (Armistice 1949)};
    \node[anchor=west, align=left] at (12.5, 2.0) {\color{popblue!50!popdark}\rule{1cm}{2pt} C\^otes / Fronti\`eres 1947};
    % Arrow North
    \draw[->, thick] (13.5, 0.5) -- (13.5, 1.5) node[above] {N};
    % Title
    \node[align=center, font=\bfseries\small] at (6, 5.5) {Carte comparative : Plan de partage ONU (1947) vs Ligne verte (1949)};
\end{tikzpicture}
\legende{Carte sch\'ematique de la Palestine mandataire. Les hachures \textbf{bleues (NE)} montrent le territoire attribu\'e \`a l'\'Etat juif par la r\'esolution 181 (55 \%). Les hachures \textbf{oranges (NO)} montrent l'\'Etat arabe (45 \%). La zone \textbf{violette} au centre est J\'erusalem internationalis\'ee. Le trait \textbf{vert discontinu \'epais} repr\'esente la \emph{Ligne verte} issue des accords d'armistice de 1949 : Isra\"el contr\^ole 78 \% de la Palestine mandataire, la Jordanie annexe la Cisjordanie (incl. J\'erusalem-Est), l'\'Egypte administre la bande de Gaza. L'\'Etat arabe pr\'evu par l'ONU ne verra pas le jour.}
\end{popfigure}

\courssub{La Nakba : l'exode palestinien et la naissance du probl\`eme des r\'efugi\'es}

La guerre de 1948 entra\^ine un bouleversement d\'emographique majeur : la \textbf{Nakba} (\guillemotleft catastrophe \guillemotright en arabe).

\begin{definition}[Nakba]
Terme arabe d\'esignant l'exode forc\'e ou la fuite de \textbf{700 000 \`a 800 000 Palestiniens} (sur 1,3 million d'habitants arabes de la Palestine mandataire) lors de la guerre 1947-1949, et l'interdiction de leur retour par le nouvel \'Etat d'Isra\"el. Elle englobe la destruction de plus de 400 villages, la d\'epossession fonci\`ere (loi sur les biens des absents, 1950) et la dispersion des r\'efugi\'es dans les pays voisins (Jordanie, Liban, Syrie, Cisjordanie, Gaza) et \`a l'int\'erieur d'Isra\"el (d\'eplac\'es internes). C'est le c\oe{}ur du \guillemotleft probl\`eme des r\'efugi\'es \guillemotright et du droit au retour (r\'esolution ONU 194, d\'ec. 1948).
\end{definition}

Les causes de l'exode sont d\'ebattues par l'historiographie (nouveaux historiens isra\'eliens vs r\'ecit officiel) : peur des combats, ordres de fuir des chefs arabes, expulsions militaires directes (ex. Lod/Ramla, Deir Yassine), guerre psychologique. R\'esultat : \`a l'armistice de 1949 (Rhodes), Isra\"el contr\^ole \textbf{78 \% de la Palestine mandataire} (au lieu de 55 \%), la Jordanie annexe la Cisjordanie et J\'erusalem-Est, l'\'Egypte administre Gaza. \textbf{Aucun \'Etat palestinien ne na\^it.} Le conflit change de nature : d'affrontement intercommunautaire sous mandat, il devient un conflit inter\'etatique (Isra\"el vs \'Etats arabes) avec une dimension nationale palestinienne refoul\'ee mais persistante.

\begin{exoresolu}[Analyse de la r\'esolution 181 : pourquoi l'\'echec imm\'ediat ?]
\textbf{Document :} Extrait de la r\'esolution 181 (II), 29 nov. 1947. \guillemotleft L'Assembl\'ee g\'en\'erale [...] Recommande au Royaume-Uni [...] et \`a tous les autres Membres des Nations Unies l'adoption et la mise en \oe{}uvre, en ce qui concerne le futur gouvernement de la Palestine, du Plan de partage avec union \'economique [...] \guillemotright.
\textbf{Question :} En vous appuyant sur le document et vos connaissances, expliquez pourquoi ce plan, bien que vot\'e par une large majorit\'e \`a l'ONU, n'a pas pu \^etre mis en \oe{}uvre pacifiquement.
\textbf{Solution détaillée :}
\begin{enumerate}
    \item \textbf{Ill\'egitim\'e per\c{c}ue par les Arabes :} Le plan attribue la majorit\'e du territoire (55 \%) \`a une minorit\'e d\'emographique (32 \%) qui ne poss\`ede que 7 \% des terres. Cela viole le principe d'autod\'etermination (majorit\'e arabe) et la justice fonci\`ere. La Ligue arabe rejette tout partage.
    \item \textbf{Absence de m\'ecanisme d'ex\'ecution :} La r\'esolution est une recommandation (Chapitre VI), non une d\'ecision contraignante (Chapitre VII). L'ONU n'a pas de force arm\'ee pour l'imposer. La puissance mandataire (GB) refuse de coop\'erer \`a l'ex\'ecution (\'evacuation chaotique mai 1948).
    \item \textbf{Refus de la coexistence :} Les deux camps se pr\'eparent militairement. Le Yichouv (Plan Daleth) vise \`a contr\^oler le territoire attribu\'e \emph{et} les zones strat\'egiques arabes. Les milices arabes et l'Arm\'ee de lib\'eration arabe attaquent les convois et localit\'es juives d\`es le 30 nov. 1947.
    \item \textbf{G\'eographie intenable :} Les fronti\`eres propos\'ees sont tortueuses, cr\'eant des enclaves isol\'ees (Galil\'ee juive, J\'erusalem juive), impossibles \`a d\'efendre sans continuit\'e territoriale. La guerre devient le seul arbitre des fronti\`eres r\'eelles (Ligne verte).
\end{enumerate}
\emph{Conclusion :} Le plan de partage a fourni une couverture juridique \`a la cr\'eation d'Isra\"el, mais son rejet par le monde arabe et l'absence de volont\'e internationale de le faire respecter par la force ont conduit \`a la guerre et \`a la Nakba, posant les termes du conflit pour les 75 ann\'ees suivantes.
\end{exoresolu}

\courssub{Synth\`ese : les racines structurelles d'un conflit \guillemotleft sans fin \guillemotright}

\`A l'issue de cette premi\`ere phase (1897-1949), trois faits structurels sont pos\'es et conditionnent la suite :
\begin{enumerate}
    \item \textbf{Deux nationalismes l\'egitimes et exclusifs} sur un m\^eme territoire : le sionisme (r\'ealis\'e en \'Etat souverain) et le nationalisme palestinien (ni\'e, dispers\'e, r\'efugi\'e).
    \item \textbf{Une asym\'etrie de puissance} : Isra\"el dispose d'un \'Etat, d'une arm\'ee, d'institutions, de soutiens internationaux ; les Palestiniens sont apatrides, divis\'es entre pays d'accueil hostiles et territoires occup\'es (Jordanie/\'Egypte).
    \item \textbf{Des fronti\`eres provisoires (Ligne verte) non reconnues} comme fronti\`eres d\'efinitives, laissant ouverte la question des territoires (Cisjordanie, Gaza, J\'erusalem, Golan) et du retour des r\'efugi\'es.
\end{enumerate}
C'est sur ces bases que s'ouvre la p\'eriode suivante : les guerres inter\'etatiques (1956, 1967, 1973), la mont\'ee de la r\'esistance palestinienne (OLP), les Intifadas et les processus de paix inachev\'es (Oslo, Camp David), que nous \'etudierons dans la section 3.

\begin{popfigure}
\begin{tikzpicture}[scale=0.9, every node/.style={font=\small}]
    % Timeline axis
    \draw[popline, thick] (0,0) -- (16,0);
    % Ticks and events
    \foreach \x/\date/\event/\color in {
        0/1897/Congr\`es de B\^ale\\(Sionisme politique)/popblue,
        3.2/1917/D\'eclaration Balfour\\(Foyer national juif)/poporange,
        8.5/1936-39/Grande R\'evolte arabe\\(Livre blanc 1939)/popred,
        11.5/1947/R\'esolution ONU 181\\(Plan de partage)/poppurple,
        13.5/1948/Proclamation Isra\"el\\(14 mai) \& Guerre\\Nakba (700-800k r\'ef.)/popgreen,
        15.5/1949/Accords Rhodes\\(Ligne verte)/popdark} {
        \draw[\color, thick] (\x, -0.2) -- (\x, 0.2);
        \node[\color, anchor=south, align=center] at (\x, 0.5) {\textbf{\date}};
        \node[anchor=north, align=center, text width=2.2cm] at (\x, -0.7) {\event};
    }
    % Arrow
    \draw[->, thick] (16,0) -- (16.5,0) node[right] {Temps};
    % Title
    \node[align=center, font=\bfseries] at (8, 3.5) {Frise chronologique : Les origines du conflit isra\'elo-palestinien (1897-1949)};
\end{tikzpicture}
\legende{Frise chronologique des \'etapes fondatrices : 1897 (naissance du sionisme politique), 1917 (promesse britannique), 1936-39 (premier affrontement majeur), 1947 (l\'egalisme onusien), 1948 (cr\'eation d'Isra\"el et Nakba), 1949 (fronti\`eres de fait). La compression temporelle (52 ans) montre l'acc\'el\'eration brutale vers 1947-1949.}
\end{popfigure}

\begin{methode}[Comment analyser un document sur les origines du conflit (Bac) ?]
\begin{enumerate}
    \item \textbf{Identifier} : Nature (discours, trait\'e, r\'esolution, carte, t\'emoignage), auteur, date, contexte imm\'ediat.
    \item \textbf{Contextualiser} : Rattacher \`a l'une des 4 phases (Sionisme naissant, Mandat/Balfour, Mont\'ee violences, Partage/Guerre 48).
    \item \textbf{Analyser le contenu} : Vocabulaire cl\'e (foyer national, partage, Nakba, droit au retour), arguments avanc\'es, silences.
    \item \textbf{Confronter / Critiquer} : Comparer avec l'autre narratif (ex: discours Ben Gourion vs r\'esolution 194 vs m\'emoire palestinienne). Noter les biais.
    \item \textbf{Conclure} : En quoi ce document \'eclaire-t-il la responsabilit\'e partag\'ee, l'asym\'etrie, ou l'enracinement du conflit ?
\end{enumerate}
\end{methode}

\begin{exercice}[Entra\^inement Bac - Sujet type]
\textbf{Sujet :} \guillemotleft Le plan de partage de la Palestine (novembre 1947) est-il la cause ou la cons\'equence du conflit isra\'elo-palestinien ? \guillemotright Vous r\'epondrez en une vingtaine de lignes en vous appuyant sur vos connaissances et sur l'analyse de la carte comparative (figure 1) et de la frise chronologique (figure 2).
\textbf{Indications :} Structurez votre r\'eponse : 1. Le plan comme aboutissement (cons\'equence) d'un processus long (sionisme, mandat, immigration, Shoah). 2. Le plan comme d\'eclencheur imm\'ediat (cause) de la guerre 1948 et de la Nakba par son rejet arabe et son acceptation strat\'egique sioniste. 3. Le plan comme structure durable (fronti\`eres de r\'ef\'erence, r\'esolution 242, Oslo) malgr\'e son non-application.
\end{exercice}

\begin{corrige}[Exercice n°1 - \'El\'ements de correction]
\textbf{Introduction :} Le vote de la r\'esolution 181 est un \emph{tournant}, \`a la fois aboutissement et point de rupture.
\textbf{D\'eveloppement 1 (Cons\'equence) :} Il couronne 50 ans de diplomatie sioniste (B\^ale 1897), l'engagement britannique (Balfour 1917), la construction du Yichouv et l'urgence post-Shoah. La carte (Fig.1) montre qu'il l\'egalise une pr\'esence juive d\'ej\`a r\'eelle (immigration, institutions).
\textbf{D\'eveloppement 2 (Cause imm\'ediate) :} Son rejet par les Arabes (injustice d\'emographique/fonci\`ere) et son acceptation tactique par Ben Gourion d\'eclenchent la guerre civile puis inter\'etatique. La frise (Fig.2) place 1947-48 comme basculement vers la violence \'etatique. La Nakba (700-800k r\'efugi\'es) en est l'effet direct.
\textbf{D\'eveloppement 3 (Structure durable) :} Bien que non appliqu\'e (Isra\"el 78 \% vs 55 \%), le \guillemotleft partage \guillemotright reste le paradigme de r\'ef\'erence : r\'esolution 242 (1967), Oslo (1993), param\`etre Clinton (2000). La Ligne verte (Fig.1) en est l'h\'eriti\`ere d\'eform\'ee.
\textbf{Conclusion :} Ni cause unique ni simple cons\'equence, le plan de partage \emph{cristallise} les contradictions insolubles (deux peuples, une terre, deux narratifs) et institue le cadre juridique international du conflit jusqu'\`a aujourd'hui.
\end{corrige}

\coursec{Une poudrière : guerres, Intifadas et paix introuvable}

\courssub{Transition : des origines à l'embrasement}
La section précédente a montré comment la promesse contradictoire de la déclaration Balfour (1917), le mandat britannique et l'afflux sioniste ont créé un foyer de tensions en Palestine mandataire. Le plan de partage de l'ONU (novembre 1947), rejeté par les Arabes, a déclenché une guerre civile qui s'est internationalisée dès la proclamation de l'État d'Israël le 14 mai 1948. C'est le point de bascule : le conflit local devient un conflit régional, opposant un État nouveau-né à ses voisins arabes, tandis que le peuple palestinien bascule dans l'exil (\emph{Nakba}, la « catastrophe »). Nous allons suivre, chronologiquement, l'enchaînement des guerres, la montée de la résistance palestinienne, les tentatives diplomatiques et les blocages structurels qui font de cet espace une « poudrière » toujours active.

\courssub{Les grandes guerres israélo-arabes (1948--1973) : conquête, renversement, statu quo}

\begin{definition}[Guerre d'Indépendance / Nakba (1948--1949)]
Premier conflit armé opposant l'État d'Israël, proclamé le 14 mai 1948, à une coalition arabe (Égypte, Jordanie, Syrie, Liban, Irak, soutenue par la Ligue arabe). Elle se solde par la victoire israélienne, l'extension du territoire israélien au-delà du plan de partage (environ 78 \% de la Palestine mandataire), la fuite ou l'expulsion de \textbf{700 000 à 750 000 réfugiés palestiniens} (la \emph{Nakba}) et le partage du reste : la Cisjordanie et Jérusalem-Est passent sous contrôle jordanien, la bande de Gaza sous administration égyptienne. Les armistices de 1949 (île de Rhodes) ne sont pas des traités de paix : ils tracent seulement des lignes d'armistice (la \emph{Ligne verte}).
\end{definition}

\begin{definition}[Crise de Suez (1956)]
Guerre éclair déclenchée par l'attaque israélo-franco-britannique contre l'Égypte de Nasser après la nationalisation de la compagnie du canal de Suez (26 juillet 1956). Sous la pression conjointe des États-Unis et de l'URSS, les trois attaquants se retirent. Résultat : Israël évacue le Sinaï et Gaza, mais le prestige de Nasser en sort renforcé ; la cause palestinienne passe au second plan derrière le nationalisme arabe.
\end{definition}

\begin{exemplebox}[La guerre des Six Jours (5--10 juin 1967) : un tournant géopolitique majeur]
\textbf{Contexte :} fermeture du détroit de Tiran par Nasser, départ des Casques bleus de l'ONU, concentration de troupes égyptiennes au Sinaï. Israël lance une frappe préventive (\emph{Opération Moked}) qui détruit l'aviation égyptienne au sol en quelques heures.
\textbf{Bilan territorial en six jours :}
\begin{itemize}
    \item le \textbf{Sinaï} et la \textbf{bande de Gaza} (pris à l'Égypte) ;
    \item la \textbf{Cisjordanie} et \textbf{Jérusalem-Est} (pris à la Jordanie) ;
    \item le \textbf{plateau du Golan} (pris à la Syrie).
\end{itemize}
\textbf{Conséquence majeure :} la superficie contrôlée par Israël \textbf{triple} environ (d'environ 20 000 à plus de 60 000 km$^2$). Environ 300 000 nouveaux réfugiés palestiniens. Jérusalem-Est est annexée de facto (annexion officialisée par la loi de 1980, non reconnue internationalement). Début de la colonisation dans les territoires occupés.
\end{exemplebox}

\begin{popfigure}
\begin{tikzpicture}[scale=0.85, every node/.style={font=\small}]
    % Coastline simplified
    \draw[popblue!60, line width=1.5pt] (0,0) .. controls (1,1) and (3,2) .. (5,2.5) .. controls (6,3) and (7,2) .. (8,1) .. controls (9,0) and (10,-1) .. (11,-1.5);
    % Jordan River / Dead Sea
    \draw[popblue!60, line width=1pt] (4.5,2.2) -- (5.5,0.5) -- (6,-1);
    % Green Line (Armistice 1949)
    \draw[popgreen!70!black, dashed, line width=1.2pt] (2,2.5) -- (3,1.5) -- (4.5,1.2) -- (5.5,0) -- (6.5,-1.2);
    % Occupied territories shading
    \fill[poporange!30] (2,2.5) -- (3,1.5) -- (4.5,1.2) -- (5.5,0) -- (6.5,-1.2) -- (7,0) -- (6,1.5) -- (4,3) -- cycle;
    \node[poporange!80!black, align=center] at (3.5,2.2) {Cisjordanie\\(Jérusalem-Est)};
    \fill[poporange!30] (7.5,1) -- (9,0.5) -- (9.5,-0.5) -- (8,-0.5) -- cycle;
    \node[poporange!80!black] at (8.5,0.3) {Gaza};
    \fill[poporange!30] (9.5,3) -- (11,3.5) -- (11,2.5) -- (10,2) -- cycle;
    \node[poporange!80!black, rotate=-20] at (10.2,2.8) {Golan};
    % Israel core
    \fill[popblue!20] (0,0) -- (2,2.5) -- (4,3) -- (6,1.5) -- (7,0) -- (8,-0.5) -- (8,-1.5) -- (0,-1.5) -- cycle;
    \node[popblue!80!black, align=center] at (2.5,0.5) {Israël\\ (frontières 1949)};
    % Labels
    \node[anchor=east, popdark] at (-0.5,1.5) {Méditerranée};
    \node[anchor=west, popdark] at (11.5,1) {Jordanie};
    \node[anchor=west, popdark] at (11.5,-1) {Syrie};
    \node[anchor=north, popdark] at (5.5,-1.8) {Égypte (Sinaï occupé 1967--1982)};
    % Legend
    \begin{scope}[shift={(0.5,3.5)}]
        \fill[popblue!20] (0,0) rectangle (0.4,0.2); \node[anchor=west] at (0.5,0.1) {Israël (1949)};
        \fill[poporange!30] (0,0.3) rectangle (0.4,0.5); \node[anchor=west] at (0.5,0.4) {Territoires occupés 1967};
        \draw[popgreen!70!black, dashed, line width=1.2pt] (0,0.65) -- (0.4,0.65); \node[anchor=west] at (0.5,0.65) {Ligne verte (1949)};
    \end{scope}
    % North arrow
    \draw[->, popdark, thick] (10.5,3.5) -- (10.5,4);
    \node[above] at (10.5,4) {N};
\end{tikzpicture}
\legende{Carte schématique (croquis simplifié, sans échelle exacte) : territoires occupés par Israël lors de la guerre des Six Jours (juin 1967). En orange : Cisjordanie (avec Jérusalem-Est), bande de Gaza, plateau du Golan ; le Sinaï, également occupé, sera restitué à l'Égypte en 1982. La ligne verte en pointillés marque la ligne d'armistice de 1949.}
\end{popfigure}

\begin{definition}[Guerre du Kippour / du Ramadan (octobre 1973)]
Attaque surprise égypto-syrienne lancée le 6 octobre 1973 (jour du Kippour, pendant le Ramadan) pour récupérer le Sinaï et le Golan. Après un début très difficile, Israël renverse la situation (franchissement du canal de Suez par les troupes d'Ariel Sharon). Les deux supergrands interviennent (alerte nucléaire américaine) ; l'ONU adopte la résolution \textbf{338} (cessez-le-feu, application de la résolution 242). L'embargo pétrolier de l'OPEP provoque le premier choc pétrolier. Résultat politique : l'honneur arabe est restauré, ce qui ouvre la voie à la diplomatie (visite de Sadate à Jérusalem, 1977).
\end{definition}

\courssub{L'émergence de l'acteur palestinien : OLP, résistance et Intifadas}

\begin{definition}[OLP (Organisation de libération de la Palestine)]
Créée en 1964 sous l'égide de la Ligue arabe, reprise en main par les fedayins (le Fatah de Yasser Arafat) après la défaite de 1967. Sa Charte nationale (1968) proclame la lutte armée pour libérer la Palestine. Elle est reconnue comme « unique représentant légitime du peuple palestinien » (sommet arabe de Rabat, 1974 ; Arafat invité à la tribune de l'ONU, 1974). Exilée à Beyrouth (1970--1982), puis à Tunis (1982--1994).
\end{definition}

\begin{definition}[Intifada (« soulèvement » en arabe)]
Soulèvement populaire massif dans les territoires occupés.
\begin{itemize}
    \item \textbf{Première Intifada (1987--1993) :} déclenchée à Gaza (décembre 1987) après un accident de la circulation impliquant un camion israélien. « Révolte des pierres » : jets de pierres, grèves, boycotts, désobéissance civile, durement réprimés par l'armée israélienne. \textbf{Conséquence :} internationalisation de la cause palestinienne, reconnaissance mutuelle OLP--Israël (échange de lettres Arafat--Rabin, 1993), accords d'Oslo.
    \item \textbf{Seconde Intifada ou Intifada d'Al-Aqsa (2000--2005) :} déclenchée après la visite d'Ariel Sharon sur l'esplanade des Mosquées (septembre 2000). Caractère militarisé : attentats-suicides (Hamas, Jihad islamique) et ripostes lourdes (chars, aviation, assassinats ciblés, construction du mur). Bilan : environ 3 000 Palestiniens et 1 000 Israéliens tués. Échec durable du processus de paix.
\end{itemize}
\end{definition}

\begin{popfigure}
\begin{tikzpicture}[scale=0.9, every node/.style={font=\small}]
    % Wall structure
    \draw[popmuted, line width=4pt, rounded corners=2pt] (0,0) -- (12,0);
    \draw[popmuted, line width=4pt, rounded corners=2pt] (0,0.4) -- (12,0.4);
    % Concrete panels
    \foreach \x in {0.5, 2.5, 4.5, 6.5, 8.5, 10.5} {
        \draw[popmuted, line width=1pt] (\x,0) -- (\x,0.4);
        \fill[popmuted!50] (\x,0) rectangle (\x+1.8,0.4);
    }
    % Watchtower
    \fill[popdark] (3,0.4) rectangle (3.4,1.2);
    \fill[popdark] (2.8,1.2) rectangle (3.6,1.4);
    \draw[popgold, thick] (3.2,1.4) -- (3.2,1.8);
    % Graffiti
    \node[poppink, rotate=15, font=\tiny] at (1.5,0.2) {Freedom};
    \node[popgreen, rotate=-10, font=\tiny] at (5.5,0.2) {Justice};
    \node[popblue, rotate=5, font=\tiny] at (9.5,0.2) {Peace};
    % Palestinian side (East)
    \fill[poporange!20] (-1,-0.5) rectangle (0,1.5);
    \node[anchor=east, poporange!80!black, font=\tiny, align=right] at (-0.1,1.3) {Côté palestinien\\ (Cisjordanie)};
    \draw[->, poporange, thick] (-0.5,0.8) -- (-0.1,0.8);
    % Israeli side (West)
    \fill[popblue!20] (12,-0.5) rectangle (13,1.5);
    \node[anchor=west, popblue!80!black, font=\tiny, align=left] at (12.1,1.3) {Côté israélien\\ (Israël / colonies)};
    \draw[->, popblue, thick] (12.5,0.8) -- (12.9,0.8);
    % Ground
    \fill[popline!30] (-1,-0.5) rectangle (13,-0.7);
    % Scale
    \draw[|-|, popdark] (0,-1) -- (2,-1);
    \node[below] at (1,-1.1) {environ 8 m de haut};
\end{tikzpicture}
\legende{Schéma du mur de séparation en Cisjordanie (barrière construite par Israël à partir de 2002). Hauteur d'environ 8 m (blocs de béton), tours de guet, graffitis. Il ne suit pas la Ligne verte mais s'enfonce en Cisjordanie pour inclure des blocs de colonies. La Cour internationale de justice (avis de 2004) a jugé sa construction en territoire occupé contraire au droit international.}
\end{popfigure}

\courssub{Le cadre onusien et les résolutions de référence}

\begin{propriete}[Résolution 242 du Conseil de sécurité (22 novembre 1967)]
Principes fondateurs : « \emph{retrait des forces armées israéliennes de territoires occupés lors du récent conflit} » (la version anglaise dit \emph{from territories}, la version française \emph{des territoires} : ambiguïté délibérée) \textbf{en échange du} « \emph{respect de la souveraineté, de l'intégrité territoriale et de l'indépendance politique de chaque État de la région} » (donc reconnaissance d'Israël). C'est la formule « \textbf{la paix contre les territoires} ». Elle fonde toutes les négociations ultérieures (Camp David, Madrid, Oslo, initiative arabe de 2002).
\end{propriete}

\begin{propriete}[Résolution 338 (22 octobre 1973)]
Appel au cessez-le-feu immédiat (guerre du Kippour) et mise en œuvre de la résolution 242 « dans toutes ses parties ». Elle lie le cessez-le-feu militaire au règlement politique du conflit.
\end{propriete}

\courssub{La diplomatie de paix : avancées, espoirs brisés, impasses}

\begin{exemplebox}[Accords de Camp David (septembre 1978) -- Traité de paix Égypte--Israël (mars 1979)]
\textbf{Acteurs :} Anouar el-Sadate (Égypte), Menahem Begin (Israël), Jimmy Carter (États-Unis, médiateur).
\textbf{Contenu :} deux volets. 1) Traité de paix bilatéral : restitution du Sinaï à l'Égypte (achevée en 1982), reconnaissance d'Israël, démilitarisation du Sinaï. 2) Cadre pour la paix au Proche-Orient : autonomie palestinienne en Cisjordanie et à Gaza pendant cinq ans, statut final reporté.
\textbf{Résultat :} première reconnaissance d'Israël par un État arabe. Sadate et Begin reçoivent le prix Nobel de la paix (1978). Sadate est assassiné en 1981 par des islamistes égyptiens. Le volet palestinien reste lettre morte (pas de gel de la colonisation, OLP exclue).
\end{exemplebox}

\begin{exemplebox}[Accords d'Oslo (Déclaration de principes, 13 septembre 1993, Washington)]
\textbf{Acteurs :} Yasser Arafat (OLP), Yitzhak Rabin (Premier ministre israélien), Shimon Peres (ministre des Affaires étrangères), Bill Clinton (États-Unis).
\textbf{Mécanisme :} reconnaissance mutuelle (l'OLP reconnaît Israël et renonce à la violence ; Israël reconnaît l'OLP comme représentant du peuple palestinien). Transfert progressif de pouvoirs à une \textbf{Autorité palestinienne} (AP) à Gaza et Jéricho (mai 1994), puis dans les villes de Cisjordanie (zones A, B, C). Les questions du statut final (Jérusalem, réfugiés, colonies, frontières, eau) sont renvoyées à des négociations devant aboutir en cinq ans (échéance 1999).
\textbf{Suites :} prix Nobel de la paix 1994 (Arafat, Rabin, Peres). Assassinat de Rabin (4 novembre 1995) par un extrémiste juif (Yigal Amir). Le processus s'enlise ; le sommet de Camp David II (juillet 2000, Barak--Arafat--Clinton) échoue sur Jérusalem et les réfugiés. Déclenchement de la seconde Intifada.
\end{exemplebox}

\begin{savaistu}[Le Nobel de la paix 1994 : une réconciliation symbolique]
Le comité Nobel a récompensé Arafat, Rabin et Peres pour « \emph{leurs efforts en faveur de la paix au Moyen-Orient} ». Rabin a déclaré : « \emph{Nous qui avons combattu contre vous, les Palestiniens, nous vous disons aujourd'hui d'une voix forte et claire : assez de sang et de larmes. Assez.} » L'assassinat de Rabin un an plus tard a brisé cet élan.
\end{savaistu}

\courssub{Les nœuds structurels du conflit aujourd'hui : pourquoi la paix reste introuvable}

\begin{center}
\begin{tabularx}{\linewidth}{|l|X|X|}
\hline
\rowcolor{popblueL} \textbf{Nœud} & \textbf{Position israélienne (gouvernements successifs)} & \textbf{Position palestinienne (OLP / AP) et consensus international} \\ \hline
\textbf{Jérusalem} & Capitale « unifiée et éternelle » d'Israël (loi de 1980). Refus du partage. & Jérusalem-Est = capitale du futur État palestinien. Statut spécial pour les Lieux saints. \\ \hline
\textbf{Colonies de peuplement} & Maintien des « blocs de colonies » (Maale Adumim, Gush Etzion, Ariel). Justifications sécuritaires et historico-bibliques. & Illégales au regard du droit international (art. 49 de la IVe Convention de Genève). Obstacle à la continuité territoriale de l'État palestinien. Environ 700 000 colons en Cisjordanie et à Jérusalem-Est (2024). \\ \hline
\textbf{Réfugiés} & Le « droit au retour » mettrait fin à l'État juif (démographie). Solution : retour dans l'État palestinien, indemnisations, retours symboliques limités. & Droit au retour (résolution 194 de l'ONU). Choix individuel : retour, réinstallation, indemnisation. Environ 5,9 millions de réfugiés enregistrés à l'UNRWA (2023). \\ \hline
\textbf{Sécurité} & Contrôle de la vallée du Jourdain et de l'espace aérien, démilitarisation de l'État palestinien, reconnaissance d'Israël comme « État-nation du peuple juif ». & Fin de l'occupation, souveraineté pleine et entière, garanties internationales. \\ \hline
\textbf{Eau} & Contrôle de la nappe phréatique de montagne (l'essentiel des prélèvements bénéficie à Israël et aux colonies). & Partage équitable des ressources (nappe de Cisjordanie, nappe côtière de Gaza, Jourdain). \\ \hline
\textbf{Mur / barrière} & « Barrière de sécurité » anti-attentats (forte baisse des attentats-suicides après 2003). & « Mur de l'apartheid » (avis CIJ 2004) : annexion rampante, fragmentation du territoire, violation du droit international. \\ \hline
\end{tabularx}
\end{center}

\begin{definition}[Ligne verte (\emph{Green Line})]
Ligne d'armistice de 1949 (Rhodes) séparant Israël de la Cisjordanie (alors jordanienne) et de Gaza (alors égyptienne). Frontière \emph{de facto} reconnue par la communauté internationale comme base des négociations (résolution 242), mais non reconnue comme frontière souveraine par Israël. Le tracé du mur s'en écarte sur une grande partie de son parcours pour inclure colonies et points d'eau.
\end{definition}

\begin{definition}[Colonie de peuplement]
Implantation civile israélienne en territoire occupé (Cisjordanie, Jérusalem-Est, Golan). Illégale au regard du droit international humanitaire (article 49 de la IVe Convention de Genève : « \emph{la Puissance occupante ne pourra procéder à la déportation ou au transfert d'une partie de sa propre population civile dans le territoire occupé par elle} »). Israël distingue les « blocs de colonies » (qu'il entend conserver) et les « avant-postes » sauvages (souvent illégaux même en droit israélien).
\end{definition}

\courssub{La division palestinienne : deux entités, deux gouvernances (depuis 2007)}

\begin{exoresolu}[Analyse de la scission Fatah / Hamas]
\textbf{Énoncé :} Expliquer les causes et les conséquences de la division palestinienne depuis 2007.
\tcblower
\textbf{Solution détaillée :}
\begin{enumerate}
    \item \textbf{Origines idéologiques :} le Fatah (nationaliste, laïc, pilier de l'OLP, favorable à la négociation) s'oppose au Hamas (islamiste, issu des Frères musulmans ; sa charte de 1988 refuse l'existence d'Israël et prône la résistance armée ; il rejette les accords d'Oslo).
    \item \textbf{Déclencheur électoral :} aux élections législatives de janvier 2006, le Hamas remporte une victoire surprise (74 sièges sur 132). Le Quartet (États-Unis, Union européenne, ONU, Russie) conditionne son aide à trois exigences : reconnaissance d'Israël, renonciation à la violence, respect des accords antérieurs. Le Hamas refuse.
    \item \textbf{Conflit armé :} en juin 2007, après des combats à Gaza, le Hamas prend le contrôle total de la bande. Mahmoud Abbas (Fatah) dissout le gouvernement d'union et installe un gouvernement en Cisjordanie.
    \item \textbf{Conséquences :} deux gouvernements, deux administrations, deux forces de sécurité. Blocus israélo-égyptien de Gaza (contrôle des frontières, de l'espace aérien et maritime). Toutes les tentatives de réconciliation échouent (Le Caire 2011, 2014, 2017 ; Alger 2022).
    \item \textbf{Impact sur la paix :} plus d'interlocuteur palestinien unique. Israël négocie (ou refuse de négocier) avec l'AP de Cisjordanie et considère le Hamas comme une organisation terroriste. La division sert d'argument contre la création d'un État palestinien viable.
\end{enumerate}
\textbf{Conclusion :} la fracture géographique (Gaza / Cisjordanie) et idéologique (islamisme / nationalisme laïc) paralyse le projet national palestinien. \textbf{Points clés pour le Bac :} citer les dates (élections de 2006, prise de Gaza en 2007), le blocus, le Quartet, l'échec répété des réconciliations.
\end{exoresolu}

\courssub{Complément Bac : la guerre déclenchée le 7 octobre 2023 et l'embrasement régional}

\begin{aretenir}[Guerre Israël--Hamas (depuis le 7 octobre 2023) -- Points de repère pour l'épreuve]
\begin{itemize}
    \item \textbf{7 octobre 2023 :} attaque massive du Hamas (opération « Déluge d'Al-Aqsa ») : environ 1 200 morts en Israël, en majorité des civils, et environ 250 personnes prises en otage. Rupture brutale du statu quo sécuritaire.
    \item \textbf{Riposte israélienne :} opération « Épées de fer ». Objectifs déclarés : détruire les capacités militaires et gouvernementales du Hamas, libérer les otages. Bombardements massifs, offensive terrestre, déplacement forcé d'environ 1,9 million de Gazaouis (sur 2,3 millions), crise humanitaire majeure (famine, destruction des infrastructures). Bilan : plus de 35 000 morts à Gaza selon le ministère de la Santé local (mai 2024), en majorité des femmes et des enfants selon l'ONU.
    \item \textbf{Dimension régionale :} au nord, échanges de tirs quotidiens avec le Hezbollah (Liban) ; au Yémen, attaques des Houthis (pro-iraniens) contre la navigation en mer Rouge ; en avril 2024, première attaque directe de l'Iran (missiles et drones) contre Israël après la frappe contre le consulat iranien à Damas. L'« \emph{axe de la résistance} » (Iran, Hamas, Hezbollah, Houthis, milices d'Irak et de Syrie) est activé.
    \item \textbf{Diplomatie et justice internationale :} blocages au Conseil de sécurité (veto américain aux premières résolutions de cessez-le-feu) ; résolutions de l'ONU exigeant un cessez-le-feu, la libération des otages et l'acheminement de l'aide humanitaire ; saisine de la Cour pénale internationale ; procédure engagée par l'Afrique du Sud devant la Cour internationale de justice (mesures conservatoires ordonnées).
    \item \textbf{Enjeu du « jour d'après » :} qui gouvernera Gaza ? L'AP réformée ? Une force internationale ? La solution des deux États est relancée par les États arabes, les États-Unis et l'Union européenne, mais rejetée par le gouvernement Netanyahu.
\end{itemize}
\end{aretenir}

\courssub{Frise chronologique : guerres et paix (1948--2024)}

\begin{popfigure}
\begin{tikzpicture}[scale=0.95, every node/.style={font=\small}]
    % Axis
    \draw[popline, thick, ->] (0,0) -- (16,0);
    % Years ticks
    \foreach \x/\year in {0/1948, 1.5/1956, 3/1967, 4.5/1973, 6/1987, 7.5/1993, 9/2000, 10.5/2007, 12/2023} {
        \draw[popline] (\x, -0.1) -- (\x, 0.1);
        \node[below, popdark] at (\x, -0.3) {\year};
    }
    \node[below, popdark] at (16, -0.3) {2024};
    % War bars (orange)
    \draw[poporange, line width=6pt] (0, 0.5) -- (1.2, 0.5); \node[above, poporange!80!black, align=center] at (0.6, 0.9) {Guerre\\1948-49};
    \draw[poporange, line width=6pt] (1.5, 0.5) -- (1.8, 0.5); \node[above, poporange!80!black, align=center] at (1.65, 0.9) {Suez\\1956};
    \draw[poporange, line width=6pt] (3, 0.5) -- (3.1, 0.5); \node[above, poporange!80!black, align=center] at (3.05, 0.9) {6 Jours\\1967};
    \draw[poporange, line width=6pt] (4.5, 0.5) -- (4.7, 0.5); \node[above, poporange!80!black, align=center] at (4.6, 0.9) {Kippour\\1973};
    \draw[poporange, line width=6pt] (6, 0.5) -- (7.5, 0.5); \node[above, poporange!80!black, align=center] at (6.75, 0.9) {Intifada I\\1987-93};
    \draw[poporange, line width=6pt] (9, 0.5) -- (10.2, 0.5); \node[above, poporange!80!black, align=center] at (9.6, 0.9) {Intifada II\\2000-05};
    \draw[poporange, line width=6pt] (12, 0.5) -- (13.5, 0.5); \node[above, poporange!80!black, align=center] at (12.75, 0.9) {Guerre\\Gaza 2023-};
    % Peace bars (green)
    \draw[popgreen!70!black, line width=6pt] (4.8, -0.5) -- (5.2, -0.5); \node[below, popgreen!70!black, align=center] at (5, -0.9) {Camp David\\1978-79};
    \draw[popgreen!70!black, line width=6pt] (7.5, -0.5) -- (8, -0.5); \node[below, popgreen!70!black, align=center] at (7.75, -0.9) {Oslo I\\1993};
    \draw[popgreen!70!black, line width=6pt] (8.2, -0.5) -- (8.5, -0.5); \node[below, popgreen!70!black, align=center] at (8.35, -0.9) {Oslo II\\1995};
    \draw[popblue, line width=6pt] (10.5, -0.5) -- (10.7, -0.5); \node[below, popblue, align=center] at (10.6, -0.9) {Blocus\\Gaza 2007};
    % Key events markers
    \node[poppurple, anchor=south, align=center] at (0, 1.3) {\textbullet\ Création d'Israël / Nakba};
    \node[poppurple, anchor=south, align=center] at (3, 1.3) {\textbullet\ Rés. 242};
    \node[poppurple, anchor=south, align=center] at (4.5, 1.3) {\textbullet\ Rés. 338};
    \node[poppurple, anchor=south, align=center] at (7.5, 1.3) {\textbullet\ Reconnaissance mutuelle};
    \node[poppurple, anchor=south, align=center] at (8.5, 1.3) {\textbullet\ Assassinat de Rabin};
    \node[poppurple, anchor=south, align=center] at (10.5, 1.3) {\textbullet\ Scission Fatah/Hamas};
    \node[poppurple, anchor=south, align=center] at (12, 1.3) {\textbullet\ 7 octobre / embrasement};
    % Legend
    \begin{scope}[shift={(0.5, -1.8)}]
        \draw[poporange, line width=6pt] (0,0) -- (1,0); \node[anchor=west] at (1.1,0) {Guerres / violences majeures};
        \draw[popgreen!70!black, line width=6pt] (0,0.4) -- (1,0.4); \node[anchor=west] at (1.1,0.4) {Accords de paix / avancées};
        \draw[popblue, line width=6pt] (0,0.8) -- (1,0.8); \node[anchor=west] at (1.1,0.8) {Blocus / division};
        \fill[poppurple] (0,1.2) circle (2pt); \node[anchor=west] at (1.1,1.2) {Jalons diplomatiques et institutionnels};
    \end{scope}
\end{tikzpicture}
\legende{Frise chronologique synthétique : alternance des guerres (orange) et des tentatives de paix (vert). Noter la densité des conflits et la brièveté des fenêtres de négociation. Depuis 2007, le blocus de Gaza et la division palestinienne verrouillent l'horizon politique.}
\end{popfigure}

\courssub{Méthode : analyser une photographie d'actualité (mur, checkpoint, manifestation)}

\begin{methode}[Démarche en 4 étapes pour le commentaire de document iconographique]
\begin{enumerate}
    \item \textbf{Description objective (« Que vois-je ? ») :} type de document (photo, caricature, carte), auteur, date, source. Plans (avant-plan, arrière-plan). Acteurs (tenues, attitudes, nombre). Objets, symboles, inscriptions (graffitis, drapeaux, panneaux). Couleurs, cadrage, lumière.
    \item \textbf{Contextualisation (« Quand ? Où ? Pourquoi ? ») :} situer l'événement précis ou la situation structurelle (ex. : mur construit à partir de 2002, blocus de Gaza depuis 2007, « Grande marche du retour » de 2018). Rappeler les acteurs en présence, les rapports de force, le cadre juridique (occupation, droit humanitaire).
    \item \textbf{Interprétation (« Que signifie-t-il ? ») :} intention du photographe, message visuel. Symbolique (le mur : séparation, oppression ; les graffitis : résistance, espoir). Rapport de force visible (soldat contre civil, mur contre olivier). Portée politique : illustration d'un nœud du conflit (sécurité contre droits, annexion contre souveraineté).
    \item \textbf{Conclusion critique (« Quelle portée ? ») :} limites du document (instantané, subjectivité, cadrage). Mise en perspective avec d'autres sources (cartes, statistiques, textes juridiques). Lien avec la problématique du cours (poudrière, paix introuvable).
\end{enumerate}
\end{methode}

\begin{attention}[Piège majeur : le réductionnisme religieux]
\textbf{Erreur fréquente :} présenter le conflit comme une « guerre de religions » (Juifs contre Musulmans) ou un « choc de civilisations » ancestral.
\textbf{Réalité historique :} c'est \textbf{d'abord un conflit territorial et national} entre deux mouvements nationaux (sionisme et nationalisme palestinien) revendiquant la même terre. La dimension religieuse (Jérusalem et ses Lieux saints, Hamas, sionisme religieux) est un \textbf{facteur aggravant majeur} qui rigidifie les positions et sacralise le territoire, mais elle n'est pas la cause originelle. Au Bac, une analyse qui ignore les dimensions nationaliste, juridique (occupation, droit au retour) et géopolitique (eau, sécurité, alliances régionales) est incomplète et sanctionnée.
\end{attention}

\begin{exercice}[Entraînement Bac -- Sujet type]
\textbf{Sujet :} « Le Proche-Orient, une poudrière sans fin ? » À partir de vos connaissances et de la frise chronologique ci-dessus, rédigez une réponse structurée (introduction, développement en deux ou trois parties, conclusion) montrant comment l'enchaînement des guerres, l'échec des processus de paix et la fragmentation palestinienne entretiennent l'instabilité régionale.
\textbf{Indications :} mobiliser les résolutions 242 et 338, les guerres de 1967 et 1973, Oslo, les Intifadas, le mur, la division Fatah--Hamas, la guerre de 2023. Utiliser les notions : Ligne verte, colonies, réfugiés, sécurité, axe de la résistance.
\end{exercice}

\begin{corrige}[Exercice n°1 -- Éléments de corrigé]
\textbf{Introduction :} accroche (7 octobre 2023), définition de l'espace (Proche-Orient : Israël, territoires palestiniens et États voisins du Levant), problématique (pourquoi l'instabilité persiste-t-elle malgré les traités de paix ?), annonce du plan.
\textbf{I. L'enracinement du conflit par la guerre et l'occupation (1948--1993).} Guerres de 1948 (Nakba), 1967 (conquête des territoires), 1973 (choc pétrolier). Résolution 242 (« la paix contre les territoires ») inappliquée. Création de l'OLP, première Intifada : montée en puissance de l'acteur palestinien.
\textbf{II. L'illusion d'une paix négociée et ses limites (1978--2000).} Camp David (paix séparée Égypte--Israël), Oslo (reconnaissance mutuelle, Autorité palestinienne) mais pas de gel de la colonisation ni de règlement du statut final. Assassinat de Rabin, échec de Camp David II, seconde Intifada (militarisation).
\textbf{III. La fragmentation palestinienne et l'embrasement régional (2007--2024).} Scission Fatah--Hamas (2007), blocus de Gaza, guerres récurrentes (2008, 2014, 2021, 2023). Le mur comme annexion rampante. Régionalisation : Iran, Hezbollah, Houthis (axe de la résistance). Normalisations arabes (accords d'Abraham, 2020) conclues sans règlement de la question palestinienne.
\textbf{Conclusion :} la « poudrière » s'explique par l'absence de solution conforme au droit international, l'asymétrie des forces et les instrumentalismes régionaux. La solution des deux États devient improbable (colonisation continue). Une paix durable suppose l'application du droit international (résolutions 242 et 194), la sécurité mutuelle et une solution juste pour les réfugiés.
\end{corrige}

\coursec{Les conflits du Moyen-Orient et du golfe Persique}

Après avoir analysé le foyer de tension permanent que constitue le problème israélo-palestinien (voir section précédente), nous devons élargir notre regard à l'ensemble du Moyen-Orient et du golfe Persique. Cette région, cœur énergétique de la planète, est le théâtre d'une superposition de conflits : rivalités régionales, fractures religieuses, ambitions nationalistes, ingérences des grandes puissances et enjeux économiques vitaux. Comprendre cette « poudrière » élargie, c'est saisir pourquoi la paix y reste si fragile et pourquoi les secousses qui s'y produisent rejaillissent jusqu'à nos économies, y compris au Cameroun.

\begin{aretenir}[Articulation avec le dossier israélo-palestinien]
Le conflit israélo-palestinien et les conflits du Golfe ne sont pas étanches. L'Iran instrumentalise la cause palestinienne (soutien au Hamas, au Jihad islamique, au Hezbollah) pour légitimer son hégémonie régionale (« Axe de la résistance »). Les Accords d'Abraham (2020) ont normalisé les relations Israël/Émirats-Bahreïn sans régler la question palestinienne, créant une fracture au sein du monde arabe sunnite. Toute flambée à Gaza (ex. 2023-2024) risque d'embraser les fronts libanais, syrien, yéménite et irakiens.
\end{aretenir}

\courssub{Les enjeux structurels du golfe Persique : énergie, géographie et rivalités}

Le golfe Persique n'est pas seulement une mer intérieure ; c'est le \cle{carrefour énergétique mondial}. Ses fonds recèlent environ \textbf{50 \% des réserves mondiales de pétrole prouvées} et \textbf{40 \% de celles de gaz naturel}. Cinq pays riverains (Arabie saoudite, Irak, Iran, Koweït, Émirats arabes unis) figurent parmi les premiers producteurs mondiaux. Cette manne financière structure les États (États-rentes), alimente les budgets militaires et attire les convoitises.

\begin{popfigure}
\begin{tikzpicture}[scale=0.85, every node/.style={font=\small}]
% Mer
\fill[popblueL] (0,0) rectangle (14,9);
% Côtes simplifiées (polygones)
\fill[popgreenL] (0,0) -- (0,9) -- (2,9) -- (3,7) -- (4,8) -- (5,6) -- (6,7) -- (7,5) -- (8,6) -- (9,4) -- (10,5) -- (11,3) -- (12,4) -- (13,2) -- (14,2) -- (14,0) -- cycle;
% Pays - zones colorées
\fill[poporange!30] (1,1) rectangle (4,4); % Iraq (approx)
\fill[poppurple!30] (4,1) rectangle (7,5); % Iran (approx)
\fill[popgold!30] (7,1) rectangle (9,3); % Koweït (petit)
\fill[poppink!30] (8,3) rectangle (13,8); % Arabie Saoudite (grand)
\fill[popgreen!30] (11,1) rectangle (13,3); % EAU / Qatar zone
% Détroit d'Ormuz
\draw[popred, line width=2pt, dashed] (6.5,0.5) -- (7.5,0.5);
\node[popred, below] at (7,0.2) {\textbf{Détroit d'Ormuz}};
% Villes / Points
\node[circle,fill=popdark,inner sep=1.5pt] (bag) at (2.5,2.5) {};
\node[anchor=west] at (bag.east) {Bagdad};
\node[circle,fill=popdark,inner sep=1.5pt] (teh) at (5.5,3) {};
\node[anchor=west] at (teh.east) {Téhéran};
\node[circle,fill=popdark,inner sep=1.5pt] (kuw) at (8,2) {};
\node[anchor=west] at (kuw.east) {Koweït City};
\node[circle,fill=popdark,inner sep=1.5pt] (rya) at (10.5,5.5) {};
\node[anchor=west] at (rya.east) {Riyad};
\node[circle,fill=popdark,inner sep=1.5pt] (dub) at (12,2) {};
\node[anchor=west] at (dub.east) {Dubaï / Abou Dabi};
% Légendes Pays
\node[poporange, font=\bfseries] at (2.5,4.5) {IRAK};
\node[poppurple, font=\bfseries] at (5.5,5.5) {IRAN};
\node[popgold, font=\bfseries] at (8,3.5) {KOWEÏT};
\node[poppink, font=\bfseries] at (11,7) {ARABIE SAOUDITE};
\node[popgreen, font=\bfseries] at (12.5,3.5) {EAU / QATAR};
% Flèche Nord
\node[anchor=south] at (13.5,8.5) {\Large $\uparrow$ N};
\end{tikzpicture}
\legende{Carte schématique du golfe Persique : enjeux pétroliers, États riverains et point de passage stratégique du détroit d'Ormuz (par où transite ~20 \% du pétrole mondial).}
\end{popfigure}

La géographie impose un \cle{point de passage obligé} : le \textbf{détroit d'Ormuz} (large de 39 km au plus étroit). Tout pétrole exporté par mer par l'Irak, le Koweït, l'Arabie saoudite (partie est), les Émirats, le Qatar et l'Iran doit l'emprunter. Qui contrôle Ormuz tient l'économie mondiale en otage.

À cette géographie s'ajoutent des \cle{fractures identitaires et religieuses} majeures :
\begin{itemize}
    \item Le clivage \textbf{Chiites (Iran, majorité en Irak, minorités au Koweït, Bahreïn, Arabie saoudite orientale) / Sunnites (Arabie saoudite, pays du Golfe, majorité dans le monde musulman)}.
    \item Le clivage \textbf{Arabes / Perses (Iran) / Turcs (Turquie)}.
    \item La \textbf{question kurde} : peuple sans État (~30-40 millions) partagé entre Turquie, Iran, Irak, Syrie.
\end{itemize}

\begin{definition}[Guerre par procuration]
Une \textbf{guerre par procuration} (ou \emph{proxy war}) est un conflit où deux puissances rivales évitent l'affrontement direct en soutenant (armes, argent, formation, conseillers) des acteurs locaux (gouvernements, milices, groupes rebelles) qui combattent en leur nom. Elle permet de limiter les coûts politiques et humains directs pour les parrains tout en étendant leur influence.
\end{definition}

\begin{definition}[Coalition internationale]
Une \textbf{coalition internationale} est un regroupement temporaire d'États, souvent sous mandat de l'ONU (Chapitre VII de la Charte), pour mener une action militaire commune contre un État considéré comme agresseur ou menace pour la paix. Elle légitime l'intervention par le droit international (ex. : résolution 678 pour la guerre du Golfe 1991).
\end{definition}

\begin{definition}[Embargo]
Un \textbf{embargo} est une mesure de sanction économique décidée par un État ou un groupe d'États (souvent via l'ONU) interdisant tout ou partie des échanges commerciaux (importations/exportations, transfert de technologie, gel d'avoirs) avec un pays cible pour le contraindre à changer de politique.
\end{definition}

\courssub{La rupture de 1979 : la Révolution iranienne et la naissance de la République islamique}

Jusqu'en 1979, l'Iran du Shah Mohammad Reza Pahlavi est le \textbf{gendarme du Golfe} allié des États-Unis (pacte CENTO, armes américaines, pétrole contre sécurité). La \textbf{Révolution iranienne} (janvier-février 1979) renverse cette architecture.

\begin{experience}[Analyse de documents : la rupture iranienne]
\textbf{Document 1 (extrait discours Khomeiny, 1er février 1979) :} « L'Amérique est le grand Satan... Nous voulons un gouvernement islamique fondé sur le Coran. »
\textbf{Document 2 (Données) :} Rupture relations diplomatiques USA-Iran (avril 1980, crise des otages ambassade Téhéran 444 jours). Nationalisation compagnies pétrolières. Expulsion conseillers militaires US.
\textbf{Question :} En quoi 1979 constitue-t-il une \emph{rupture systémique} pour l'équilibre du Golfe ?
\textbf{Interprétation :} L'Iran passe du statut d'allié occidental à celui de \textbf{puissance révisionniste} cherchant à \emph{exporter la révolution islamique} chiite. L'Arabie saoudite (sunnite, gardienne des lieux saints) se sent menacée dans sa légitimité. Les USA perdent leur pilier régional et basculent vers un soutien accru à l'Irak de Saddam Hussein (laïc, sunnite dominant, anti-iranien).
\end{experience}

\courssub{La guerre Iran-Irak (1980-1988) : un million de morts pour un statu quo}

Saddam Hussein, profitant du chaos iranien, attaque le 22 septembre 1980 (revendication du Chatt-el-Arab, peur de la contagion chiite). Ce sera la \textbf{plus longue guerre interétatique conventionnelle du XXe siècle}.

\begin{exemplebox}[La guerre Iran-Irak (1980-1988) : chiffres-clés]
\begin{itemize}
    \item \textbf{Durée :} 8 ans (sept. 1980 - août 1988).
    \item \textbf{Pertes humaines :} \textbf{~ 1 million de morts} (tués, disparus), ~ 1,5 à 2 millions de blessés/invalidés. Proportionnellement à la population de l'époque (~55 millions pour les deux pays), l'équivalent d'environ \textbf{1,2 million de morts pour la France actuelle}.
    \item \textbf{Armes chimiques :} Utilisation massive par l'Irak (moutarde, sarin, tabun) contre troupes iraniennes (non protégées) et civils kurdes (Halabja, mars 1988 : ~5 000 morts en quelques heures).
    \item \textbf{Guerre des villes / des pétroliers :} Bombardements de Téhéran, Bagdad ; attaques de tankers dans le Golfe (affrètement de navires sous pavillon neutre).
    \item \textbf{Coût économique :} ~ 1 000 milliards \$ (cumulé des deux pays). Endettement massif de l'Irak auprès des pays du Golfe (Koweït, Arabie saoudite) et de l'Occident.
    \item \textbf{Résultat :} \textbf{Statu quo ante bellum} (frontières inchangées, Chatt-el-Arab partagé). Aucune victoire militaire, seulement l'épuisement.
\end{itemize}
\end{exemplebox}

\begin{aretenir}[Conséquences structurelles de la guerre Iran-Irak]
\begin{enumerate}
    \item \textbf{Militarisation extrême} des deux sociétés ; culture du martyr (Iran) / culte du chef guerrier (Irak).
    \item \textbf{Endettement de l'Irak} : Saddam réclame l'effacement de sa dette auprès du Koweït et de l'Arabie saoudite (« j'ai défendu la porte arabe contre les Perses »). Refus koweïtien $\rightarrow$ tension majeure.
    \item \textbf{Renforcement du régime des Mollahs} en Iran (unité nationale face à l'agresseur).
    \item \textbf{Implication des grandes puissances :} USA, URSS, France, Chine arment les deux camps (équilibre de la terreur). Affaire « Iran-Contra » (vente d'armes US à l'Iran via Israël pour financer Contras au Nicaragua).
\end{enumerate}
\end{aretenir}

\courssub{De l'invasion du Koweït (1990) à l'intervention de 2003 : l'Irak au cœur de la tourmente}

\courssub{L'invasion du Koweït et la Guerre du Golfe (1990-1991)}

Le 2 août 1990, l'Irak envahit et annexe le Koweït (19e province). Motifs : dette, accusation de forage oblique (vol de pétrole), accès à la mer, ambition hégémonique.

\begin{methode}[Construire une frise chronologique à partir d'un texte documentaire]
\textbf{Objectif :} Synthétiser une séquence événementielle dense.
\begin{enumerate}
    \item \textbf{Repérer les dates explicites} (ex : « 2 août 1990 », « 17 janvier 1991 », « 28 février 1991 »).
    \item \textbf{Identifier les phases} : Déclencheur $\rightarrow$ Réaction diplomatique (ONU) $\rightarrow$ Ultimatum $\rightarrow$ Opération militaire $\rightarrow$ Cessez-le-feu $\rightarrow$ Conséquences.
    \item \textbf{Hiérarchiser les événements} : Distinguer les \emph{causes}, les \emph{faits majeurs} (batailles, résolutions), les \emph{conséquences immédiates}.
    \item \textbf{Positionner sur l'axe} : Échelle régulière (ex : 1 cm = 1 mois). Utiliser un code couleur (Diplomatie / Militaire / Humanitaire).
    \item \textbf{Légender précisément} : Date + Événement court + Acteur principal.
\end{enumerate}
\end{methode}

\begin{popfigure}
\begin{tikzpicture}[scale=0.9, every node/.style={font=\footnotesize}]
% Axe principal
\draw[popdark, thick] (0.5,0) -- (15.5,0);
% Graduations mensuelles (août 1990 à mars 1991)
\foreach \m/\label in {0/Août 90, 1/Sept, 2/Oct, 3/Nov, 4/Déc, 5/Janv 91, 6/Fév, 7/Mars} {
    \draw (\m+0.5,0.1) -- (\m+0.5,-0.1) node[below=2pt] {\label};
}
% Événements - Diplomatie (au-dessus)
\draw[popblue, thick] (0.5,0.5) node[above=2pt, align=center] {\textbf{2 août}\\ Invasion\\ Koweït} -- (0.5,0.1);
\draw[popblue, thick] (2.5,1.2) node[above=2pt, align=center] {\textbf{Nov}\\ Résolution 678\\ (Ultimatum 15 janv)} -- (2.5,0.1);
% Événements - Militaire (au-dessous)
\draw[popred, thick] (5, -0.5) node[below=2pt, align=center] {\textbf{17 janv}\\ Opération\\ \textbf{Tempête du désert}\\ (Bombardements)} -- (5, -0.1);
\draw[popred, thick] (6, -1.2) node[below=2pt, align=center] {\textbf{24 fév}\\ Offensive\\ terrestre\\ (100 h)} -- (6, -0.1);
\draw[popred, thick] (6.5, -0.5) node[below=2pt, align=center] {\textbf{28 fév}\\ Cessez-le-feu\\ Libération Koweït} -- (6.5, -0.1);
% Conséquences
\draw[poppurple, thick] (7.5, 1.5) node[above=2pt, align=center] {\textbf{Mars 91}\\ Insurrections\\ Chiites/Kurdes\\ Répression} -- (7.5, 0.1);
\draw[poppurple, thick] (8.5, -1.5) node[below=2pt, align=center] {\textbf{Avr 91}\\ Résolution 687\\ Désarmement\\ Embargo / Zone d'exclusion} -- (8.5, -0.1);
% Légende couleurs
\node[popblue, anchor=west] at (10, 1.5) {■ Diplomatie / ONU};
\node[popred, anchor=west] at (10, 1.0) {■ Opérations militaires};
\node[poppurple, anchor=west] at (10, 0.5) {■ Conséquences / Sanctions};
\end{tikzpicture}
\legende{Frise chronologique : Invasion du Koweït (août 1990) $\rightarrow$ Guerre du Golfe (janv-fév 1991) $\rightarrow$ Sanctions et zones d'exclusion (1991+).}
\end{popfigure}

La \textbf{coalition internationale} (34 pays, 700 000 hommes, mandat ONU Rés. 678) lance l'\textbf{Opération Tempête du Désert} (17 janv - 28 fév 1991). Victoire militaire éclatante (supériorité aérienne, technologie), libération du Koweït. Mais \textbf{Saddam Hussein reste au pouvoir}. Les insurgés chiites (sud) et kurdes (nord), encouragés par Bush père, sont écrasés. L'ONU impose le \textbf{désarmement (Rés. 687)}, l'embargo pétrolier strict (\emph{Oil-for-Food} à partir de 1996, Rés. 986), et des \textbf{zones d'exclusion aérienne} (nord/sud) patrouillées par USA/UK (France initialement, retrait 1998).

\begin{savaistu}[Catastrophe écologique de 1991]
En se retirant, l'armée irakienne applique la politique de la « terre brûlée » : \textbf{plus de 600 puits de pétrole koweïtiens sont incendiés} (certains brûleront 10 mois). Fumées noires jusqu'en Inde, « pluie noire », nappes phréatiques polluées, 1 à 1,5 milliard de barils perdus. Coût extinction : ~ 1,5 milliard \$. L'une des pires catastrophes écologiques de guerre de l'histoire.
\end{savaistu}

\courssub{L'invasion de l'Irak (2003) : le chaos et l'émergence de Daech}

Après le 11 septembre 2001, l'administration Bush (fils) change de doctrine : \textbf{guerre préventive}, « axe du mal » (Irak, Iran, Corée du Nord). Accusation : \textbf{ADM (Armes de Destruction Massive)} et liens avec Al-Qaïda. \textbf{Pas de mandat ONU explicite} (Veto menacé France/Russie/Chine sur résolution de guerre). Invasion le 20 mars 2003 (coalition « des volontaires » : USA, UK, Pologne, Australie...). Bagdad tombe le 9 avril. Saddam capturé déc. 2003, pendu 2006.

\begin{attention}[Piège majeur : Confondre 1991 et 2003]
\begin{itemize}
    \item \textbf{Guerre du Golfe (1991) :} \emph{Libération} du Koweït envahi. \textbf{Mandat ONU clair}. Coalition large (monde arabe inclus). Saddam \textbf{reste au pouvoir}. Objectif : \emph{statu quo ante}.
    \item \textbf{Guerre d'Irak (2003) :} \emph{Invasion} de l'Irak souverain. \textbf{Pas de mandat ONU} (légalité contestée). Coalition réduite. Objectif : \textbf{changement de régime} (nation-building). Saddam \textbf{renversé et exécuté}.
    \item \textbf{Conséquence 2003 :} Démantèlement armée/BAAS $\rightarrow$ vide sécuritaire $\rightarrow$ guerre civile confessionnelle (chiites vs sunnites) $\rightarrow$ montée d'Al-Qaïda en Irak $\rightarrow$ \textbf{Daech (État islamique) proclamé juin 2014} (Mossoul, Raqqa). Califat territorial détruit militairement 2017-2019, mais résidus actifs.
\end{itemize}
\end{attention}

\courssub{Les guerres par procuration : la « Guerre froide » saoudo-iranienne}

Depuis 1979, l'Arabie saoudite (chef de file sunnite, allié US) et l'Iran (chef de file chiite, anti-US) s'affrontent sans guerre directe, via des \textbf{guerres par procuration} sur trois théâtres majeurs :

\begin{center}
\begin{tabularx}{\linewidth}{|l|X|X|X|}\hline
\rowcolor{popblueL}
\textbf{Théâtre} & \textbf{Acteurs soutenus par l'Iran (Chiites)} & \textbf{Acteurs soutenus par l'Arabie Saoudite (Sunnites)} & \textbf{Enjeu} \\ \hline
\textbf{Yémen (depuis 2014)} & Houthis (Ansar Allah) - Contrôle Nord/Sanaa & Gouvernement Hadi (reconnu ONU), Conseil de transition sud, coal. arabe (SA, EAU) & Contrôle Bab el-Mandeb (détroit Mer Rouge), encerclement SA \\ \hline
\textbf{Syrie (depuis 2011)} & Bachar el-Assad (Alaouites/Chiites), Hezbollah, milices irakiennes/afghanes, Russie (défenseur) & Rebelles sunnites (FSA, puis djihadistes), Turquie, Qatar, USA (initiellement) & Survie régime Assad, couloir terrestre Iran-Méditerranée (Hezbollah) \\ \hline
\textbf{Liban} & \textbf{Hezbollah} (Parti de Dieu, armée de fait, bras armé iranien) & Courants sunnites (ex-Futur Hariri), Forces libanaises (chrétiennes), Israël (ennemi commun) & Influence politique libanaise, front anti-Israël, arsenal roquettes (150 000+) \\ \hline
\end{tabularx}
\end{center}

\begin{popfigure}
\begin{tikzpicture}[scale=0.85, every node/.style={font=\small}]
% Axe temps
\draw[popdark, thick] (0.5,0) -- (15.5,0);
% Années clés
\foreach \x/\year in {0/1979, 1.5/1980, 3.5/1988, 5.5/1990, 6.5/1991, 9.5/2003, 11.5/2011, 13.5/2014} {
    \draw (\x+0.5,0.1) -- (\x+0.5,-0.1) node[below=3pt] {\year};
}
% Barres de conflits
% Iran-Irak
\draw[poporange, line width=6pt] (1.5, 1.5) -- (3.5, 1.5) node[right=5pt, align=left] {Guerre Iran-Irak (1980-88)};
% Guerre Golfe 1991
\draw[popred, line width=6pt] (5.5, 2.5) -- (6.5, 2.5) node[right=5pt, align=left] {Guerre du Golfe (1991)};
% Invasion Irak 2003
\draw[poppurple, line width=6pt] (9.5, 3.5) -- (12.5, 3.5) node[right=5pt, align=left] {Guerre d'Irak / Insurrection (2003-2011+)};
% Printemps arabe / Syrie / Yémen
\draw[popblue, line width=6pt] (11.5, 4.5) -- (15.5, 4.5) node[right=5pt, align=left] {Guerres par procuration Syrie/Yémen (2011+)};
% Daech
\draw[poppink, line width=6pt] (13.5, 5.5) -- (15.5, 5.5) node[right=5pt, align=left] {Daech / Califat (2014-2019)};
% Révolution Iranienne (point de départ)
\fill[popgold] (0.5, 0.5) circle (4pt) node[above=5pt] {Révolution iranienne (1979)};
% Légende
\node[anchor=north west] at (0.5, 6.5) {\textbf{Frise synthétique : Enchaînement des conflits du Golfe (1979-2020)}};
\end{tikzpicture}
\legende{Frise chronologique majeure : de la Révolution iranienne (1979) à la défaite territoriale de Daech (2019), en passant par la guerre Iran-Irak, les deux guerres du Golfe/Irak et les guerres par procuration.}
\end{popfigure}

\courssub{La question kurde : le plus grand peuple sans État}

Les Kurdes (~35-45 millions) sont la \cle{quatrième ethnie du Moyen-Orient} (après Arabes, Turcs, Perses). Répartis sur 4 États :
\begin{itemize}
    \item \textbf{Turquie (~15-20 M) :} PKK (Parti des travailleurs du Kurdistan, marxiste-léniniste puis autonomiste) en guérilla depuis 1984 (> 40 000 morts). Répression dure, interdiction langue/culture (assouplie puis durcie). Opérations transfrontalières en Irak/Syrie.
    \item \textbf{Irak (~5-6 M) :} \textbf{Kurdistan irakien} : autonomie \emph{de facto} depuis 1991 (zone d'exclusion), institutionnalisée Constitution 2005 (Gouvernement régional KRG, Peshmergas). Référendum indépendance 2017 (92 \% oui) $\rightarrow$ riposte Bagdad (reprise Kirkouk).
    \item \textbf{Syrie (~2-3 M) :} PYD/YPG (branche syrienne PKK) $\rightarrow$ \textbf{Administration autonome du Nord-Est (Rojava)}. Alliés USA contre Daech. Menacés par Turquie (opérations « Bouclier de l'Euphrate », « Source de Paix »).
    \item \textbf{Iran (~8-10 M) :} Répression forte, partis d'opposition (PDKI, Komala) en exil ou clandestins. Pas d'autonomie.
\end{itemize}
\begin{aretenir}[La question kurde : facteur d'instabilité régionale permanent]
Les Kurdes sont des \textbf{acteurs stratégiques} (Peshmergas, YPG/Forces démocratiques syriennes) courtisés par les USA (lutte anti-Daech) mais perçus comme \textbf{menace existentielle} par la Turquie (intégrité territoriale) et l'Iran/Irak/Syrie (sécessionnisme). Leur sort conditionne tout règlement régional.
\end{aretenir}

\courssub{Le rôle des grandes puissances et de l'OPEP : géopolitique de l'énergie}

\begin{itemize}
    \item \textbf{États-Unis :} Garants de la sécurité du Golfe (Carter Doctrine 1980 : « toute tentative de contrôle du Golfe sera repoussée par la force »). Bases massives (Qatar, Bahreïn, Koweït, EAU, Arabie Saoudite). Pivot vers l'Asie (Obama/Trump/Biden) mais maintien engagement (pétrole, Israël, containment Iran). Accords d'Abraham (2020) : normalisation Israël-EAU/Bahreïn/Soudan/Maroc $\rightarrow$ front anti-Iran.
    \item \textbf{Russie :} Retour en force depuis 2015 (intervention Syrie sauvant Assad). Partenaire OPEP+ (coordination coupes production avec Arabie Saoudite). Vente d'armes (S-400 à Turquie, S-300 à Iran). Rivalité/coopération avec USA.
    \item \textbf{Chine :} Premier importateur mondial de pétrole. Partenaire économique majeur (Nouvelles Routes de la Soie). Médiation surprise Iran-Arabie Saoudite (mars 2023, accords de Pékin). Ne veut pas de rôle militaire direct mais sécurise ses approvisionnements.
    \item \textbf{OPEP / OPEP+ :} Cartel (13 pays + 10 alliés dont Russie). Gère l'offre pour stabiliser les prix. Décisions impactent \textbf{directement le budget du Cameroun} (pétrole = ~40 \% recettes export, ~20 \% PIB). Choc pétrolier 2020 (Covid + guerre des prix SA/Russie) $\rightarrow$ prix négatifs historique (WTI avril 2020) $\rightarrow$ crise budgétaire camerounaise.
\end{itemize}

\begin{exoresolu}[Analyse d'un document statistique : Production pétrolière du Golfe (2023, millions barils/jour)]
\begin{center}
\begin{tabularx}{\linewidth}{|l|X|X|}\hline
\rowcolor{popblueL}
\textbf{Pays} & \textbf{Production (Mb/j)} & \textbf{\% Production OPEP} \\ \hline
Arabie Saoudite & 9,0 & ~30 \% \\ \hline
Irak & 4,3 & ~14 \% \\ \hline
Iran & 3,6 (sous sanctions) & ~12 \% \\ \hline
Koweït & 2,7 & ~9 \% \\ \hline
Émirats Arabes Unis & 3,2 & ~10 \% \\ \hline
\textbf{Total Golfe} & \textbf{~22,8} & \textbf{~75 \%} \\ \hline
\end{tabularx}
\end{center}
\textbf{Questions :}
\begin{enumerate}
    \item Calculer la part de la production du Golfe dans la production totale de l'OPEP (donnée : ~30 Mb/j).
    \item Expliquer pourquoi l'Arabie Saoudite est le \emph{swing producer} (producteur d'appoint) incontournable.
    \item Relier ces chiffres à la vulnérabilité économique du Cameroun.
\end{enumerate}
\tcblower
\textbf{Solution détaillée :}
\begin{enumerate}
    \item Production OPEP totale $\approx$ 30 Mb/j. Part Golfe = $22,8 / 30 = 0,76$ soit \textbf{76 \%}. Le Golfe domine l'offre OPEP.
    \item L'Arabie Saoudite a la \textbf{plus grande capacité excédentaire} (peut augmenter/baisser rapidement de 1-2 Mb/j). Elle assume le rôle de « banque centrale du pétrole » pour éviter effondrement ou envolée des prix. Son coût d'extraction est le plus bas (~3 \$/baril), lui donnant une marge de manœuvre financière immense.
    \item Le Cameroun produit ~40 000 à 60 000 b/j (peu significatif sur marché mondial mais vital pour budget). Une baisse de 10 \$/baril (décision OPEP+ ou chute demande) = perte de ~150 à 200 milliards FCFA/an pour le budget de l'État (selon prix de référence loi finances). L'instabilité du Golfe (guerre, blocus Ormuz) $\rightarrow$ volatilité prix $\rightarrow$ instabilité budgétaire à Yaoundé.
\end{enumerate}
\end{exoresolu}

\begin{popfigure}
\begin{tikzpicture}
\begin{axis}[
    ybar,
    bar width=12pt,
    width=\linewidth,
    height=8cm,
    enlarge x limits=0.2,
    ylabel={Production (millions de barils/jour)},
    xtick={1,2,3,4,5},
    xticklabels={Arabie Saoudite, Irak, Iran, EAU, Koweït},
    nodes near coords,
    nodes near coords align={vertical},
    ymin=0,
    ymax=10,
    ymajorgrids=true,
    grid style=dashed,
    title={Production pétrolière des principaux pays du Golfe (estimation 2023)},
    title style={yshift=0.5cm},
    every axis plot/.append style={fill=popblue!70},
]
\addplot coordinates {
    (1, 9.0)
    (2, 4.3)
    (3, 3.6)
    (4, 3.2)
    (5, 2.7)
};
\end{axis}
\end{tikzpicture}
\legende{Diagramme en barres : Dominance écrasante de l'Arabie Saoudite (producteur d'appoint) et poids cumulé du Golfe dans l'offre mondiale. Source : OPEP Annual Statistical Bulletin / EIA.}
\end{popfigure}

\courssub{Synthèse : pourquoi des tensions « sans fin » ?}

La région piège dans un \textbf{équilibre instable} à cinq dimensions :
\begin{enumerate}
    \item \textbf{Ressource malédiction :} La rente pétrolière finance des États autoritaires, des armées puissantes et des groupes armés non-étatiques, tout en décourageant la diversification économique et la démocratie.
    \item \textbf{Faillite de l'ordre régional :} Pas d'organisation de sécurité collective inclusive (l'Iran exclu du CCG, la Ligue arabe divisée). Rivalité hégémonique Iran vs Arabie Saoudite (religieuse + géopolitique).
    \item \textbf{Ingérence extérieure permanente :} USA (sécurité Israël + pétrole), Russie (retour puissance), Chine (énergie). Aucune ne veut la guerre générale, toutes alimentent le statu quo par ventes d'armes et alliances.
    \item \textbf{Fragmentation identitaire :} Frontières Sykes-Picot (1916) artificielles $\neq$ réalités ethno-confessionnelles (Kurdes, Chiites/Sunnites, minorités chrétiennes/yézidies). États-nations fragiles (Irak, Syrie, Liban, Yémen).
    \item \textbf{Non-résolution du dossier palestinien :} Cause mobilisatrice pour l'opinion arabe/musulmane, instrumentalisée par l'Iran (axe de la résistance) et obstacle à la normalisation totale (Accords d'Abraham incomplets).
\end{enumerate}

\begin{methode}[Rédiger une analyse géopolitique structurée (type Bac)]
\textbf{Sujet type :} « Les conflits du Golfe sont-ils condamnés à durer ? »
\begin{enumerate}
    \item \textbf{Analyser le sujet :} Mots-clés : « conflits du Golfe », « condamnés », « durer ». Problématique : \emph{En quoi la superposition des enjeux énergétiques, identitaires et stratégiques verrouille-t-elle les issues de sortie des conflits du Golfe Persique ?}
    \item \textbf{Plan dialectique :}
    \begin{itemize}
        \item \textbf{I. Une conflictualité structurelle ancrée (Facteurs de permanence).} A. Ressource/Énergie (pétrole, Ormuz). B. Identités/Religion (Chiites/Sunnites, Kurdes). C. Géopolitique (Iran vs SA, Grandes Puissances).
        \item \textbf{II. Des dynamiques de transformation et de régulation (Facteurs d'évolution).} A. Épuisement modèles (guerres coûteuses, opinion publique US réticente). B. Nouveaux acteurs (Chine médiatrice, Turquie puissance régionale). C. Transition énergétique mondiale (pic pétrolier demande ?) $\rightarrow$ baisse rente future.
        \item \textbf{III. Scénarios pour l'avenir.} A. Statut quo géré (conflits gelés, proxy wars). B. Conflagration majeure (blocus Ormuz, guerre Iran-Israël/USA). C. Réorganisation régionale (Grand Moyen-Orient apaisé ? Peu probable court terme).
    \end{itemize}
    \item \textbf{Rédiger :} Intro (accroche, définition, pb, plan). Développement (arguments + exemples précis : dates, noms, chiffres). Conclusion (réponse nuancée, ouverture).
\end{enumerate}
\end{methode}

\begin{exercice}[Entraînement Bac - Sujet de réflexion]
\textbf{Sujet :} « Le golfe Persique, cœur énergétique de la planète, est-il voué à rester une zone de guerres permanentes ? »
\textbf{Consignes :} Vous rédigerez une réponse structurée (introduction, développement en deux ou trois parties équilibrées, conclusion) en mobilisant les connaissances du cours (1979, Iran-Irak, 1991, 2003, guerres par procuration, Kurdes, OPEP, grandes puissances) et en intégrant au moins un exemple chiffré précis.
\end{exercice}

\begin{corrige}[Exercice n°1 - Éléments de corrigé]
\textbf{Introduction :} Accroche sur prix pétrole / vie quotidienne. Définition espace (Golfe + arrière-pays). Problématique : tension entre intérêts économiques partagés (stabilité prix) et logiques de puissance/identitaires divergentes. Plan annoncé.
\textbf{Développement :}
\textbf{I. Les moteurs structurels de la conflictualité (Le « condamné »).} Ressource (75 \% prod OPEP, Ormuz 20 \% flux). Fracture chiite/sunnite (Iran/SA) + Kurdes. Ingérences (USA bases, Russie Syrie, Chine diplomatie). Exemple : Guerre Iran-Irak 1M morts, 1000 Mrd \$.
\textbf{II. Les freins à la guerre totale et pistes de régulation (Le « pas voué »).} Dissuasion nucléaire (seuil Iran), interdépendance économique (Chine client, USA producteur schiste), fatigue des opinions, diplomatie (Accords Pékin 2023, détente SA/Iran). Rôle OPEP+ comme forum de gestion de crise.
\textbf{III. L'incertitude majeure : la transition énergétique.} Baisse demande pétrole à horizon 2030-2050 $\rightarrow$ effondrement rente $\rightarrow$ soit ouverture/réformes, soit durcissement répressif/guerres de survie des régimes.
\textbf{Conclusion :} Pas « condamné » fatalement, mais \textbf{piégé} dans une violence chronique à bas bruit (proxy wars) avec risque d'embrasement majeur. Paix durable = refonte contrat social (rente vs citoyenneté) + architecture sécurité collective inclusive (difficile).
\end{corrige}

\coursec{Pourquoi des tensions « sans fin » ? Répercussions jusqu'au Cameroun}

\courssub{Transition : des guerres aux causes structurelles}

Les sections précédentes ont décrit la succession des conflits armés — guerres israélo-arabes, guerre Iran-Irak, guerres du Golfe, invasion de l'Irak en 2003 — qui ont ensanglanté la région depuis 1948. Mais au-delà des affrontements militaires, ce sont des \textbf{facteurs structurels} qui expliquent la permanence de l'instabilité. Ces facteurs s'entrelacent : le contrôle des ressources (pétrole, eau), les revendications territoriales et identitaires, le poids des religions, les ingérences des grandes puissances, la montée du terrorisme et la prolifération nucléaire. Comprendre ces moteurs profonds permet de saisir pourquoi la paix reste fragile et comment les secousses de cette région lointaine retentissent jusqu'aux pompes à essence de Yaoundé ou de Douala, et jusqu'aux choix diplomatiques du Cameroun à l'ONU.

\begin{popfigure}
\begin{tikzpicture}[every node/.style={font=\small}]
  % Centre
  \node[circle, draw=popblue, thick, fill=popblueL, minimum width=3cm, align=center] (C) at (0,0) {\textbf{Facteurs de}\\\textbf{tension}\\\textbf{permanents}};
  % Six branches positionnées manuellement
  \node[align=center, text width=2.6cm] (N1) at (90:4.2)  {\textbf{Pétrole}\\\emph{Rente, chocs, stratégie énergétique}};
  \node[align=center, text width=2.6cm] (N2) at (30:4.4)  {\textbf{Eau}\\\emph{Jourdain, Tigre, Euphrate, nappes}};
  \node[align=center, text width=2.8cm] (N3) at (-30:4.4) {\textbf{Territoires}\\\emph{Jérusalem, Golan, Cisjordanie}};
  \node[align=center, text width=2.6cm] (N4) at (-90:4.2) {\textbf{Religions}\\\emph{Lieux saints, identités, extrémismes}};
  \node[align=center, text width=2.8cm] (N5) at (210:4.4) {\textbf{Ingérences}\\\emph{USA, Russie, puissances régionales}};
  \node[align=center, text width=2.8cm] (N6) at (150:4.4) {\textbf{Terrorisme et nucléaire}\\\emph{Groupes armés, programme iranien}};
  \draw[popline, thick] (C) -- (N1);
  \draw[popline, thick] (C) -- (N2);
  \draw[popline, thick] (C) -- (N3);
  \draw[popline, thick] (C) -- (N4);
  \draw[popline, thick] (C) -- (N5);
  \draw[popline, thick] (C) -- (N6);
\end{tikzpicture}
\legende{Schéma en toile d'araignée : les six facteurs structurels qui alimentent l'instabilité chronique au Proche et au Moyen-Orient.}
\end{popfigure}

\courssub{Les six piliers de l'instabilité permanente}

\begin{definition}[Choc pétrolier]
Un \cle{choc pétrolier} est une rupture brutale de l'offre ou de la demande de pétrole qui provoque une flambée soudaine et durable des cours mondiaux, avec des répercussions économiques majeures (inflation, récession, transferts de revenus vers les pays producteurs).
\end{definition}

\begin{definition}[Ingérence]
L'\cle{ingérence} désigne l'intervention d'un État ou d'une coalition dans les affaires intérieures d'un autre État souverain — militaire, diplomatique, économique ou médiatique — pour y défendre ses intérêts stratégiques ou idéologiques.
\end{definition}

\begin{aretenir}[Les six facteurs structurels]
\begin{enumerate}
  \item \textbf{Pétrole et rente énergétique :} la région détient \emph{près de la moitié des réserves mondiales prouvées} de pétrole. Le pétrole finance les États et les armées ; il attire les convoitises extérieures.
  \item \textbf{Eau :} ressource rare et partagée (Jourdain, Tigre, Euphrate, nappes souterraines). Le contrôle de l'eau est un enjeu de survie et de souveraineté (barrages turcs sur le Tigre et l'Euphrate, partage des eaux du Jourdain entre Israël, la Jordanie et les Palestiniens).
  \item \textbf{Territoires et frontières :} héritage du partage colonial (accords Sykes-Picot, 1916), frontières contestées, statuts non réglés (Jérusalem-Est, plateau du Golan, Cisjordanie, question kurde).
  \item \textbf{Religions et identités :} lieux saints des trois monothéismes à Jérusalem ; rivalité sunnites/chiites instrumentalisée (Arabie saoudite contre Iran) ; nationalismes concurrents (arabe, juif, kurde, perse, turc).
  \item \textbf{Ingérences étrangères :} Guerre froide (USA/URSS), puis rôle dominant des États-Unis, retour de la Russie, ambitions régionales (Iran, Turquie, Arabie saoudite, Israël). Chaque crise locale devient une guerre par procuration.
  \item \textbf{Terrorisme et nucléaire :} émergence de groupes transnationaux (Al-Qaïda, Daech, Hezbollah, Houthis) ; programme nucléaire iranien perçu comme une menace par Israël et les monarchies du Golfe.
\end{enumerate}
\end{aretenir}

\courssub{Chocs pétroliers et volatilité des cours : l'économie mondiale sous tension}

\begin{exemplebox}[Premier choc pétrolier, 1973-1974]
En octobre 1973, lors de la guerre du Kippour, les pays arabes de l'OPEP décident un embargo sur leurs exportations vers les pays soutenant Israël et réduisent leur production. En quelques mois, le prix du baril passe d'environ \textbf{3 dollars} à environ \textbf{12 dollars} : il est \textbf{multiplié par quatre}. Les pays occidentaux entrent en crise (inflation et ralentissement économique) ; les pays pauvres non producteurs, dont de nombreux États africains, voient leur facture d'importation exploser.
\end{exemplebox}

\begin{exoresolu}[Calcul du taux de variation du prix du baril entre 1973 et 1974]
\textbf{Énoncé :} Le prix moyen du baril était de 3,00 dollars en septembre 1973 et de 11,65 dollars en janvier 1974. Calculer le taux de variation en pourcentage.
\tcblower
\textbf{Solution :}
\[
t = \frac{P_{\text{final}} - P_{\text{initial}}}{P_{\text{initial}}} \times 100
= \frac{11,65 - 3,00}{3,00} \times 100
= \frac{8,65}{3,00} \times 100
\approx 288,3 \%
\]
Le prix a été multiplié par environ \textbf{3,9} (soit quasi \textbf{× 4}), ce qui confirme l'expression « multiplié par quatre ».
\end{exoresolu}

\begin{popfigure}
\begin{tikzpicture}
\begin{axis}[
  width=\linewidth,
  height=8cm,
  xlabel={Année},
  ylabel={Prix du baril (dollars courants)},
  xmin=1970, xmax=2025,
  ymin=0, ymax=150,
  xtick={1973,1979,1991,2008,2022},
  xticklabels={1973,1979,1991,2008,2022},
  ytick={0,20,40,60,80,100,120,140},
  grid=major,
  grid style={popline},
  axis line style={popdark},
  tick style={popdark},
  label style={font=\small, popdark},
  tick label style={font=\footnotesize, popdark},
  legend style={at={(0.02,0.98)}, anchor=north west, font=\small, fill=popcream, draw=popline},
]
\addplot[popcurve, thick, mark=*, mark size=2, mark options={fill=popblue}] coordinates {
  (1970,3) (1973,3) (1974,12) (1978,14) (1979,18) (1980,35) (1985,27) (1990,23) (1991,32) (1998,12) (2000,28) (2008,140) (2009,60) (2014,100) (2016,40) (2020,40) (2022,120) (2024,80)
};
\addplot[only marks, mark=triangle*, mark size=3, poporange] coordinates {
  (1973,3) (1979,18) (1991,32) (2008,140) (2022,120)
};
\legend{Prix du baril, Crises majeures}
\node[anchor=west, font=\tiny, poporange] at (axis cs:1973,8) {Guerre du Kippour / embargo OPEP};
\node[anchor=west, font=\tiny, poporange] at (axis cs:1979,24) {Révolution iranienne};
\node[anchor=west, font=\tiny, poporange] at (axis cs:1991,38) {Guerre du Golfe};
\node[anchor=south, font=\tiny, poporange] at (axis cs:2008,140) {Pic de 2008};
\node[anchor=south, font=\tiny, poporange] at (axis cs:2022,120) {Guerre en Ukraine};
\end{axis}
\end{tikzpicture}
\legende{Évolution indicative du prix du baril de pétrole (dollars courants) et crises majeures signalées par des triangles. Sources : données historiques OPEP et BP Statistical Review (valeurs arrondies).}
\end{popfigure}

\begin{aretenir}[Chronologie des chocs et répercussions]
\begin{itemize}
  \item \textbf{1973 (1er choc) :} embargo de l'OPEP arabe → prix multiplié par 4 → récession mondiale, endettement du Sud.
  \item \textbf{1979 (2e choc) :} révolution iranienne puis guerre Iran-Irak → nouveau doublement des cours → politiques d'austérité (plans d'ajustement structurel du FMI dans les années 1980, y compris au Cameroun).
  \item \textbf{1990-1991 (guerre du Golfe) :} hausse brutale mais transitoire des cours après l'invasion du Koweït par l'Irak.
  \item \textbf{2008 :} pic historique proche de 147 dollars (forte demande asiatique, spéculation) → crise alimentaire mondiale, émeutes de la faim, y compris au Cameroun (émeutes de février 2008).
  \item \textbf{2022 :} guerre en Ukraine et tensions sur l'offre → baril autour de 120 dollars → inflation mondiale, hausse des carburants en Afrique.
\end{itemize}
\end{aretenir}

\courssub{Répercussions directes sur le Cameroun : de la pompe à l'assiette}

Le Cameroun, bien que producteur de pétrole brut (production modeste, en déclin, de l'ordre de quelques dizaines de milliers de barils par jour), est \textbf{importateur de produits raffinés} (essence, gasoil, kérosène), d'autant plus que la raffinerie de Limbé (SONARA) est à l'arrêt depuis l'incendie de mai 2019. Chaque crise du Golfe ou flambée des cours mondiaux se répercute donc sur les prix à la pompe, puis sur l'ensemble de la chaîne des transports et des prix alimentaires.

\begin{exemplebox}[Hausse des prix à la pompe à Yaoundé et Douala (exemples récents)]
\begin{itemize}
  \item \textbf{Avant 2023 :} prix administrés et fortement subventionnés (super autour de 630 FCFA/L, gasoil autour de 575 FCFA/L).
  \item \textbf{2022 (pic lié à la guerre en Ukraine) :} le baril dépasse 120 dollars ; l'État camerounais maintient les prix par une subvention massive qui pèse lourdement sur le budget public.
  \item \textbf{Février 2024 (réduction de la subvention) :} le super passe à 840 FCFA/L et le gasoil à 828 FCFA/L, soit des hausses de plus de 30 \%.
  \item \textbf{Conséquence immédiate :} hausse des tarifs de transport urbain (taxis, motos-taxis), répercutée sur le prix des denrées alimentaires sur les marchés (Mokolo, Marché Central, Nkolndongo).
\end{itemize}
\end{exemplebox}

\begin{attention}[Piège : « Le Cameroun est producteur, donc il gagne quand le pétrole monte »]
Raisonnement faux ou du moins très incomplet. Le Cameroun exporte du pétrole \textbf{brut} mais importe des produits \textbf{raffinés}, dont le prix suit les cours mondiaux augmentés du fret et des marges de raffinage. De plus, la subvention étatique destinée à amortir le choc pèse sur le budget (moins de moyens pour les écoles, les routes, la santé). Le gain sur les exportations de brut est donc largement compensé, voire annulé, pour les finances publiques.
\end{attention}

\courssub{Le Cameroun et la diplomatie onusienne : l'équilibre difficile}

Depuis son admission à l'ONU en 1960, le Cameroun participe aux votes des résolutions de l'Assemblée générale relatives au Proche-Orient, dans le respect du droit international et des positions communes de l'Organisation de l'unité africaine (OUA), devenue Union africaine (UA).

\begin{savaistu}[Diplomatie camerounaise : l'équilibre Israël / monde arabe]
Le Cameroun entretient \textbf{des relations diplomatiques à la fois avec Israël (relations établies en 1960, rompues en 1973 après la guerre du Kippour, rétablies en 1986) et avec les États arabes}. Cette position d'équilibre se traduit par :
\begin{itemize}
  \item le soutien à \textbf{la solution à deux États} (Israël et un État palestinien vivant côte à côte en sécurité), conformément aux résolutions 242 (1967) et 338 (1973) du Conseil de sécurité ;
  \item la condamnation de \textbf{la colonisation israélienne en Cisjordanie et à Jérusalem-Est} ;
  \item la condamnation des \textbf{attaques contre les civils}, qu'ils soient israéliens ou palestiniens ;
  \item le maintien de coopérations techniques avec Israël (agriculture, irrigation, santé) \textbf{et} avec les pays arabes (fonds koweïtien et saoudien, Banque islamique de développement).
\end{itemize}
C'est une illustration concrète de la \cle{diplomatie d'équilibre et de non-alignement} pratiquée par le Cameroun depuis l'indépendance.
\end{savaistu}

\begin{exoresolu}[Analyse d'un vote camerounais à l'Assemblée générale de l'ONU (octobre 2023)]
\textbf{Énoncé :} Après l'attaque du Hamas du 7 octobre 2023 et la riposte israélienne sur Gaza, l'Assemblée générale vote une résolution appelant à une « trêve humanitaire immédiate ». Le Cameroun choisit de \textbf{s'abstenir}. Justifier cette abstention au regard de la tradition diplomatique camerounaise.
\tcblower
\textbf{Solution :}
\begin{enumerate}
  \item \textbf{Position d'équilibre :} s'abstenir permet au Cameroun de ne rompre ni avec Israël (partenaire de coopération) ni avec le monde arabe et la majorité africaine qui soutiennent la résolution.
  \item \textbf{Respect du droit international humanitaire :} le Cameroun réaffirme par ailleurs la protection des civils des deux côtés et l'accès de l'aide humanitaire à Gaza.
  \item \textbf{Cohérence avec sa tradition :} depuis 1960, la diplomatie camerounaise évite les prises de position radicales sur ce conflit et privilégie la solution à deux États négociée.
  \item \textbf{Crédibilité :} pays contributeur aux opérations de maintien de la paix (notamment en Centrafrique avec la MINUSCA), le Cameroun se veut un acteur du dialogue plutôt qu'un partisan.
\end{enumerate}
\end{exoresolu}

\courssub{Parallèle sécuritaire : Boko Haram et le terrorisme international}

L'instabilité du Proche-Orient n'est pas qu'une affaire de frontières lointaines ; elle irrigue les idéologies et les réseaux qui frappent le Cameroun, en particulier dans la région de l'Extrême-Nord.

\begin{propriete}[Lien idéologique et opérationnel : Proche-Orient \textrightarrow{} bassin du Lac Tchad]
Boko Haram (secte née au Nigeria en 2002, passée à la lutte armée en 2009) et l'État islamique en Afrique de l'Ouest (ISWAP) sont des groupes locaux, mais ils puisent leur \textbf{référentiel idéologique} (salafisme djihadiste, projet de califat), leur \textbf{stratégie de propagande} (vidéos, serments d'allégeance au « calife » de Daech) et une partie de leurs \textbf{méthodes} dans les réseaux nés au Proche-Orient (Al-Qaïda en Irak, devenu Daech). La défaite territoriale de Daech en Syrie et en Irak (2017-2019) a favorisé une dispersion de combattants et de cadres vers le Sahel et le bassin du Lac Tchad.
\end{propriete}

\begin{experience}[Comparaison des modes opératoires : Daech (2014-2019) et Boko Haram/ISWAP (depuis 2009)]
\begin{center}
\begin{tabularx}{\linewidth}{|l|X|X|}\hline
\rowcolor{popblueL}
\textbf{Dimension} & \textbf{Daech (Syrie/Irak)} & \textbf{Boko Haram / ISWAP (bassin du Lac Tchad)} \\ \hline
\textbf{Contrôle territorial} & « Califat » avec administration, impôts, tribunaux (2014-2017) & Zones de contrôle fluides (forêt de Sambisa, îles du Lac Tchad), administration rudimentaire \\ \hline
\textbf{Recrutement} & Appel mondial (combattants étrangers), propagande numérique sophistiquée & Recrutement local (pauvreté, exclusion, idéologie), endoctrinement d'enfants \\ \hline
\textbf{Financement} & Pétrole (champs saisis), taxes, rançons, pillage d'antiquités & Rançons (enlèvements), taxes sur les marchés, la pêche et l'élevage, pillages \\ \hline
\textbf{Attaques emblématiques} & Attentats de Paris (2015), Bruxelles (2016) & Enlèvement des lycéennes de Chibok (2014), attaques de Fotokol, Kolofata, Mora (Extrême-Nord du Cameroun) \\ \hline
\textbf{Réponse militaire} & Coalition internationale (USA, France, Russie, forces locales) & Force multinationale mixte (Cameroun, Nigeria, Tchad, Niger, Bénin) avec appuis extérieurs (renseignement, formation) \\ \hline
\end{tabularx}
\end{center}
\legende{Tableau comparatif : analogies structurelles et adaptations locales du modèle djihadiste.}
\end{experience}

\begin{attention}[Erreur fréquente : « Boko Haram = Daech »]
Boko Haram a prêté allégeance à Daech en 2015, ce qui a donné naissance à l'ISWAP, mais une scission a lieu en 2016 : la faction historique (dirigée par Abubakar Shekau, mort en 2021) garde le nom « Boko Haram », tandis que l'ISWAP devient la branche officiellement reconnue par l'État islamique. Les deux groupes coexistent, s'affrontent parfois, mais partagent le même \textbf{objectif} : instaurer un État islamique dans la région du Lac Tchad. Ne pas les confondre, mais ne pas les dissocier totalement non plus.
\end{attention}

\courssub{Méthode : rédiger et justifier une conclusion de dissertation en Histoire}

\begin{methode}[Conclusion de dissertation : les trois temps obligatoires]
Pour réussir la conclusion (notée au Probatoire et au Baccalauréat technique), suivez cette structure en \textbf{trois temps distincts} :
\begin{enumerate}
  \item \textbf{Rappel de la problématique et réponse directe (bilan) :} reprenez les termes du sujet et donnez une réponse nuancée (ni « tout est perdu », ni « tout est réglé »). Exemple : « Le Proche et le Moyen-Orient demeurent une zone de tension structurelle car s'y concentrent de nombreux facteurs de conflit ; pourtant, des avancées diplomatiques prouvent que la guerre n'est pas une fatalité. »
  \item \textbf{Synthèse nuancée des arguments principaux :} reprenez les grandes parties du développement en une phrase chacune, en montrant leur articulation. Exemple : « Le poids du pétrole et des ingérences entretient les rivalités entre États ; la question palestinienne reste le cœur symbolique et politique du blocage ; le terrorisme djihadiste et le nucléaire iranien ajoutent une dimension transnationale. »
  \item \textbf{Ouverture maîtrisée (perspective) :} élargissez sans sortir du sujet et évitez les lieux communs (« l'avenir nous le dira »). Proposez une piste concrète : accords d'Abraham (2020), rôle possible d'une conférence internationale sous l'égide de l'ONU, effets de la transition énergétique mondiale sur les pétromonarchies.
\end{enumerate}
\end{methode}

\courssub{Conclusion rédigée et justifiée : une zone de tension durable, mais pas figée}

\begin{aretenir}[Conclusion type — Sujet : « Le Proche et le Moyen-Orient : zones de tension sans fin ? »]
\textbf{Bilan :} le Proche et le Moyen-Orient constituent bien une « poudrière » où se concentrent \textbf{de nombreux facteurs de conflit} : ressources stratégiques (pétrole, eau), territoires disputés, identités religieuses et nationales instrumentalisées, ingérences des grandes puissances, terrorisme transnational et risque nucléaire. Cette superposition explique la \textbf{récurrence des guerres} (1948, 1956, 1967, 1973, 1980-1988, 1991, 2003, 2006, 2014, 2023) et l'\textbf{échec répété des processus de paix} (Oslo, Camp David, Annapolis).

\textbf{Nuance :} pour autant, l'Histoire ne valide pas le déterminisme du « sans fin ». La \textbf{paix israélo-égyptienne (1979)} a tenu malgré les crises. Les \textbf{accords d'Abraham (2020)} ont normalisé les relations entre Israël et plusieurs États arabes (Émirats arabes unis, Bahreïn, Maroc, Soudan). La \textbf{détente saoudo-iranienne (mars 2023, sous médiation chinoise)} prouve que les rivalités régionales peuvent être gérées diplomatiquement.

\textbf{Ouverture :} l'avenir de la région se jouera sur trois tableaux : (1) la capacité de la communauté internationale à faire aboutir une \textbf{solution à deux États viable} ; (2) l'impact de la \textbf{transition énergétique mondiale} sur la rente pétrolière qui finance les États et parfois les groupes armés ; (3) le poids des sociétés civiles qui, de Téhéran à Tel-Aviv en passant par Le Caire et Beyrouth, réclament justice et libertés. Le Cameroun, par ses votes à l'ONU et sa lutte contre Boko Haram, est partie prenante de cette quête de stabilité.
\end{aretenir}

\begin{exercice}[Entraînement Bac — Dissertation]
\textbf{Sujet :} « Les conflits du Proche et du Moyen-Orient ont-ils des répercussions sur les pays africains ? » \\
\emph{Consignes :} mobilisez les connaissances du chapitre (chocs pétroliers, diplomatie onusienne, terrorisme) et l'exemple du Cameroun. Structurez votre devoir en : introduction (accroche, définitions, problématique, annonce du plan), développement en deux ou trois parties équilibrées, conclusion (méthode ci-dessus).
\end{exercice}

\begin{corrige}[Exercice n°1 — Éléments de corrigé]
\textbf{Problématique :} en quoi les conflits du Proche et du Moyen-Orient, par leur dimension mondialisée (énergie, idéologie, droit), retentissent-ils sur les économies, la sécurité et la diplomatie des États africains, au premier rang desquels le Cameroun ?

\textbf{Plan suggéré :}
\begin{enumerate}
  \item \textbf{Le choc économique : l'énergie comme vecteur de transmission.} Pétrole → cours mondiaux → prix à la pompe (Cameroun) → inflation des transports et des denrées → subventions budgétaires → arbitrages sociaux difficiles.
  \item \textbf{Le choc sécuritaire : le terrorisme djihadiste transnational.} Idéologie (Daech, Al-Qaïda) → groupes locaux (Boko Haram/ISWAP, Al-Shabaab, groupes du Sahel) → insécurité dans l'Extrême-Nord du Cameroun et le bassin du Lac Tchad → dépenses militaires, déplacés internes et réfugiés, crise humanitaire.
  \item \textbf{Le choc diplomatique : l'Afrique à l'ONU.} Votes des résolutions (droit international, solution à deux États) → position d'équilibre du Cameroun (relations avec Israël et avec le monde arabe) → cohérence panafricaine (UA) → crédibilité dans les opérations de maintien de la paix.
\end{enumerate}

\textbf{Conclusion attendue :} oui, les répercussions sont majeures et structurelles. L'Afrique n'est pas spectatrice mais actrice (diplomatie, Force multinationale mixte, voix à l'ONU). L'apaisement dépend de la résolution du nœud palestinien et de la diversification énergétique mondiale.
\end{corrige}

\coursec{TP guidés : construire pour comprendre}

\courssub{Transition et objectifs des activités}

Après avoir analysé les mécanismes profonds qui font du Proche et du Moyen-Orient une zone de tensions chroniques — des racines historiques du conflit israélo-palestinien aux enjeux pétroliers du Golfe en passant par les dynamiques confessionnelles et géopolitiques —, il est indispensable de passer de la compréhension théorique à la maîtrise des outils de l'historien et du géographe. Cette séance de travaux pratiques a pour but de vous faire \textbf{construire} vous-même les supports de la démonstration historique : un croquis de synthèse, un diagramme statistique et une frise chronologique à double entrée. Vous manipulerez également un dossier documentaire croisé, conformément aux exigences de l'épreuve du Probatoire / Baccalauréat technique. Chaque production devra être \textbf{justifiée} à l'oral, développant ainsi la compétence « Rédiger et justifier une démarche, une réponse ou une conclusion ».

\begin{methode}[Construire un croquis de synthèse en Histoire-Géographie]
\textbf{Étape 1 : Analyser le sujet et sélectionner l'information.} Identifier les mots-clés (dates, acteurs, dynamiques spatiales). Ne retenir que les éléments signifiants pour répondre à la problématique.
\textbf{Étape 2 : Préparer le fond de carte.} Tracer le contour (côte, frontières internationales), les repères hydrographiques majeurs (Jourdain, mer Morte, lac de Tibériade) et les grandes villes de référence (Jérusalem, Tel-Aviv, Haïfa, Gaza, Hébron, Ramallah). Indiquer l'échelle et l'orientation (Nord).
\textbf{Étape 3 : Hiérarchiser la légende \textit{avant} de colorier.} Organiser en 3 à 4 rubriques logiques (ex. : « Partage ONU 1947 », « Armistice 1949 », « Guerre des Six Jours 1967 », « Situation actuelle »). Choisir des signes cartographiques adaptés : \textbf{surfaces} (aplats de couleur) pour les territoires, \textbf{traits} (lignes) pour les frontières/ligne verte, \textbf{pointillés} pour les zones disputées, \textbf{flèches} pour les mouvements/attaques, \textbf{points} pour les villes/colonies.
\textbf{Étape 4 : Réaliser le croquis avec rigueur.} Colorier proprement (crayons de couleur), tracer les traits à la règle. Respecter la sémiologie : couleurs claires pour le passé/historique, couleurs vives pour le conflit actuel, hachures pour les zones contestées.
\textbf{Étape 5 : Soigner la présentation.} Titre complet (sujet + dates + espace), légende encadrée et ordonnée, échelle graphique, flèche du Nord, source(s), nom de l'auteur.
\end{methode}

\begin{attention}[Erreurs fréquentes en croquis — Pièges à éviter]
\begin{itemize}
    \item \textbf{Surcharge informationnelle :} Vouloir tout mettre (toutes les colonies, tous les check-points, toutes les guerres) rend la carte illisible. \textit{Solution :} Ne représenter que les éléments structurants de la légende.
    \item \textbf{Oubli du titre, de l'échelle ou du Nord :} Un croquis sans ces éléments n'est pas un document géographique valide et perd des points d'emblée.
    \item \textbf{Légende non hiérarchisée :} Une simple liste d'items sans rubriques ni logique visuelle (ex. mélanger surfaces et points sans distinction). \textit{Solution :} Regrouper par thèmes (Territoires, Frontières, Villes, Conflits).
    \item \textbf{Mauvaise sémiologie :} Utiliser une surface pour une ligne (frontière) ou un point pour une zone. Respecter : surface = étendue, trait = limite/axe, point = localisation ponctuelle.
    \item \textbf{Imprécision graphique :} Coloriage qui déborde, traits tremblants, écriture illisible. Le croquis est une \textbf{production soignée}, pas un brouillon.
\end{itemize}
\end{attention}

\courssub{TP 1 — Croquis de synthèse : « Israël-Palestine : du partage de 1947 aux territoires occupés de 1967 »}

\paragraph{Consigne.} À partir du fond de carte vierge ci-dessous et de vos connaissances, réalisez le croquis complet avec sa légende hiérarchisée. Vous devez faire apparaître : le plan de partage de l'ONU (résolution 181, 1947), les lignes d'armistice de 1949 (ligne verte), les territoires occupés lors de la guerre des Six Jours (1967) — Cisjordanie, Gaza, Golan, Sinaï (restitué), Jérusalem-Est —, l'annexion de Jérusalem-Est (1980) et du Golan (1981), la localisation des principaux blocs de colonies israéliennes en Cisjordanie, et le tracé schématique du mur de séparation (début 2000s).

\begin{popfigure}
\begin{tikzpicture}[scale=1.1]
% --- Fond de carte simplifié : contour côte + hydrographie ---
% Méditerranée
\fill[popblue!15] (-3.5,-1) rectangle (5.5,5.5);
% Contour côtier simplifié (Israël / Palestine / Liban)
\draw[popdark, thick] (-3.5,5.5) -- (-3.5,0.5) .. controls (-2,1) and (-1,2) .. (-0.5,3) .. controls (0,3.5) and (1,4) .. (1.5,4.5) -- (5.5,4.5) -- (5.5,-1) -- (-3.5,-1) -- cycle;
% Lac de Tibériade
\fill[popblue!30] (0.8,3.8) ellipse ({0.35} and {0.2});
\draw[popblue] (0.8,3.8) ellipse ({0.35} and {0.2});
\node[popdark, font=\tiny] at (0.8,3.4) {Lac de Tibériade};
% Jourdain
\draw[popblue, thick] (0.8,3.6) .. controls (0.5,3) and (0.2,2) .. (0.5,1.2) .. controls (0.7,0.5) and (1.2,0) .. (1.5,-0.5);
\node[popdark, font=\tiny, rotate=-20] at (0.2,2.5) {Jourdain};
% Mer Morte
\fill[popblue!30] (1.2,1.2) .. controls (1.8,1) and (2.2,0) .. (2.5,-0.5) .. controls (2.2,-1) and (1.5,-1.2) .. (1.2,-0.8) .. controls (0.8,-0.5) and (0.8,0.5) .. (1.2,1.2);
\draw[popblue] (1.2,1.2) .. controls (1.8,1) and (2.2,0) .. (2.5,-0.5) .. controls (2.2,-1) and (1.5,-1.2) .. (1.2,-0.8) .. controls (0.8,-0.5) and (0.8,0.5) .. (1.2,1.2);
\node[popdark, font=\tiny] at (2.0,0.2) {Mer Morte};
% Frontières internationales reconnues (traits fins continus)
% Frontière Égypte (Sud) - au sud de Gaza
\draw[popdark, thin] (-1.0,-0.5) -- (0.5,-0.5) -- (0.5,-1.0);
% Frontière Jordanie (Est) - suit Jourdain / Mer Morte / Arava
\draw[popdark, thin] (0.8,3.6) .. controls (1.0,3) and (1.5,2) .. (1.8,1.2) .. controls (2.2,0) and (2.5,-0.5) .. (2.5,-1);
% Frontière Liban (Nord) - ligne bleue approximative
\draw[popdark, thin] (-3.5,4.5) .. controls (-2,4.5) and (-0.5,4.5) .. (0.5,4.5);
% Frontière Syrie / Golan (Nord-Est) - ligne pourpre (ligne de cessez-le-feu 1974)
\draw[poppurple, thin, densely dashed] (0.5,4.5) -- (1.5,4.5);
% Ligne verte (Armistice 1949) - Trait discontinu épais vert
\draw[popgreen!70!black, thick, densely dashed] (-0.5,4.5) .. controls (0,4) and (0.5,3.5) .. (1.0,3.2) .. controls (1.5,2.5) and (1.8,1.5) .. (2.0,0.5) .. controls (2.2,0) and (2.5,-0.5) .. (2.5,-1);
\node[popgreen!70!black, font=\tiny, rotate=30] at (0.5,4.0) {Ligne verte (1949)};
% Villes repères (points)
\foreach \x/\y/\name in {
    1.5/4.2/Haïfa,
    2.5/3.5/Naplouse,
    1.8/2.8/Jérusalem,
    1.0/2.0/Bethléem,
    2.2/2.2/Hébron,
    0.5/1.0/Gaza,
    3.5/4.0/Jénine,
    4.0/3.0/Amman (Jordanie),
    -1.0/4.5/Beyrouth (Liban),
    3.0/1.5/Jericho,
    1.0/4.0/Nazareth,
    3.5/2.0/Ramallah} {
    \fill[popdark] (\x,\y) circle (1.5pt);
    \node[popdark, font=\tiny, anchor=west] at (\x+0.1,\y) {\name};
}
% Échelle graphique
\draw[popdark, thick] (-3.0,-0.5) -- (-2.0,-0.5);
\node[popdark, font=\tiny, anchor=north] at (-2.5,-0.6) {100 km};
% Flèche Nord
\draw[popdark, thick, ->] (-3.0,4.8) -- (-3.0,5.3);
\node[popdark, font=\small\bfseries] at (-3.0,5.5) {N};
% Grille de repérage discrète
\foreach \x in {-3,-2,-1,0,1,2,3,4,5} {
    \draw[popmuted, very thin] (\x,-1) -- (\x,5.5);
    \node[popmuted, font=\tiny] at (\x,-1.2) {\x};
}
\foreach \y in {-1,0,1,2,3,4,5} {
    \draw[popmuted, very thin] (-3.5,\y) -- (5.5,\y);
    \node[popmuted, font=\tiny, anchor=east] at (-3.7,\y) {\y};
}
\end{tikzpicture}
\legende{Trame vierge pour le croquis TP 1. Contours côtiers, hydrographie, frontières internationales (traits fins), ligne verte (armistice 1949) en pointillés verts, villes repères. L'élève doit ajouter : zones du plan de partage 1947 (aplats), territoires 1967 (hachures/couleurs), mur de séparation (trait rouge épais), colonies (points rouges), légende hiérarchisée, titre complet.}
\end{popfigure}

\begin{exoresolu}[Modèle de légende hiérarchisée pour le TP 1]
\textbf{Titre :} Israël-Palestine : du partage de 1947 aux territoires occupés de 1967 (Situation vers 2020)

\textbf{LÉGENDE}

\textbf{I. Le partage de la Palestine (Résolution 181, nov. 1947)}
\begin{itemize}
    \item \textcolor{popblue!60!white}{\rule{5mm}{3mm}} État juif prévu (56\% du territoire)
    \item \textcolor{poporange!70!white}{\rule{5mm}{3mm}} État arabe prévu (44\% du territoire)
    \item \textcolor{poppurple!70!white}{\rule{5mm}{3mm}} Corpus separatum (Jérusalem-Bethléem, internationalisé)
\end{itemize}

\textbf{II. L'armistice de 1949 et la « Ligne Verte »}
\begin{itemize}
    \item \textcolor{popgreen!70!black}{\rule[0.5ex]{5mm}{1.5pt}} Ligne d'armistice (frontière de facto 1949-1967)
    \item \textcolor{popblue!40!white}{\rule{5mm}{3mm}} Territoire israélien (78\% de la Palestine mandataire)
    \item \textcolor{poporange!40!white}{\rule{5mm}{3mm}} Cisjordanie (sous admin. jordanienne) \& Bande de Gaza (sous admin. égyptienne)
\end{itemize}

\textbf{III. La guerre des Six Jours (juin 1967) et ses conséquences}
\begin{itemize}
    \item \textcolor{popred!50!white}{\rule{5mm}{3mm}} Territoires occupés en 1967 (Cisjordanie, Gaza, Golan, Sinaï)
    \item \textcolor{popred}{\rule[0.5ex]{5mm}{2pt}} Annexion unilatérale : Jérusalem-Est (1980), Golan (1981)
    \item $\blacktriangleright$ \textcolor{popred}{Restitution du Sinaï à l'Égypte (1979-1982)}
\end{itemize}

\textbf{IV. Dynamiques territoriales actuelles (Processus d'Oslo / Mur)}
\begin{itemize}
    \item \textcolor{popred}{\textbullet} Principaux blocs de colonies israéliennes en Cisjordanie (Ariel, Maale Adoumim, Gush Etzion)
    \item \textcolor{popred}{\rule[0.5ex]{5mm}{2.5pt}} Tracé du mur de séparation (début construction 2002) — \textit{majoritairement en Cisjordanie}
    \item \textcolor{popblue}{\textbullet} Villes palestiniennes majeures (Zones A/B admin. palestinienne)
    \item \textcolor{popdark}{\textbullet} Villes israéliennes / capitale proclamée (Jérusalem-Ouest / Tel-Aviv)
\end{itemize}

\textbf{Repères :} \textcolor{popdark}{\textbullet} Villes repères \quad \textcolor{popblue}{\rule[0.5ex]{5mm}{1.5pt}} Frontières int. reconnues \quad \textcolor{popgreen!70!black}{\rule[0.5ex]{5mm}{1.5pt}} Ligne verte
\par\medskip
\textit{Source :} Cartes ONU, IGN, données B'Tselem / OCHA. \textit{Auteur :} [Votre Nom], Terminale Tech, 2026.
\end{exoresolu}

\courssub{TP 2 — Diagramme statistique : La production pétrolière du Golfe persique}

\begin{methode}[Exploiter un tableau de données et construire un diagramme en barres]
\textbf{Étape 1 : Lire et comprendre le tableau.} Identifier les variables (pays, années, unité : millions de barils/jour ou Mt/an), la période, la source. Repérer les valeurs extrêmes et les tendances.
\textbf{Étape 2 : Choisir le type de graphique adapté.} Diagramme en barres (verticales) pour comparer des valeurs discrètes (pays) à une date donnée ou l'évolution d'un pays. Diagramme en barres groupées pour comparer plusieurs pays sur plusieurs dates.
\textbf{Étape 3 : Préparer les axes.} Axe des abscisses (x) : modalités qualitatives (pays) ou temps. Axe des ordonnées (y) : grandeur quantitative (production). \textbf{Graduation régulière}, unité indiquée entre parenthèses. Origine (0) obligatoire sur l'axe des ordonnées.
\textbf{Étape 4 : Tracer les barres.} Largeur constante, espacement régulier. Code couleur cohérent avec légende si plusieurs séries. Valeur exacte éventuellement notée au sommet.
\textbf{Étape 5 : Habiller le graphique.} Titre informatif (Quoi ? Où ? Quand ? Source ?), légende si nécessaire, source précise.
\end{methode}

\paragraph{Données fournies (Production de pétrole brut — Moyenne 2023, source OPEP / BP Statistical Review).}

\begin{center}
\begin{tabularx}{\linewidth}{|l|X|X|X|}\hline
\textbf{Pays} & \textbf{Production (millions de barils/jour)} & \textbf{Statut OPEP 2023} & \textbf{Réserves prouvées (milliards de barils)} \\ \hline
Arabie Saoudite & 9,0 & Membre & 267 \\ \hline
Irak & 4,3 & Membre & 145 \\ \hline
Iran & 3,6 & Membre & 209 \\ \hline
Émirats Arabes Unis & 3,2 & Membre & 111 \\ \hline
Koweït & 2,7 & Membre & 102 \\ \hline
Qatar & 1,3 & Hors OPEP (depuis 2019) & 25 \\ \hline
\textbf{Total OPEP (13 membres)} & \textbf{$\sim$ 27,5} & \textbf{--} & \textbf{1 240} \\ \hline
\end{tabularx}
\end{center}

\paragraph{Consigne.} Construire un diagramme en barres verticales comparant la production des 5 pays du Golfe membres de l'OPEP en 2023, en ajoutant le Qatar (producteur majeur hors OPEP) pour comparaison. Utiliser le code \texttt{pgfplots} ci-dessous comme base, le compléter et l'insérer dans votre copie (ou le reproduire à la main sur papier millimétré avec la même rigueur).

\begin{popfigure}
\begin{tikzpicture}
\begin{axis}[
    ybar,
    bar width=10pt,
    width=\linewidth,
    height=9cm,
    enlarge x limits=0.15,
    ylabel={Production (millions de barils/jour)},
    ylabel style={font=\small},
    xlabel={Pays du Golfe - 2023},
    xlabel style={font=\small},
    symbolic x coords={Arabie Saoudite, Irak, Iran, EAU, Koweït, Qatar},
    xtick=data,
    xticklabel style={rotate=30, anchor=east, font=\tiny},
    ymin=0,
    ymax=10,
    ymajorgrids=true,
    grid style={popmuted, dashed},
    legend style={at={(0.5,-0.25)}, anchor=north, legend columns=-1, font=\tiny, /tikz/every even column/.append style={column sep=0.5cm}},
    nodes near coords,
    nodes near coords align={vertical},
    every node near coord/.append style={font=\tiny, popdark},
]
\addplot[popblue!70, fill=popblue!60] coordinates {
    (Arabie Saoudite, 9.0)
    (Irak, 4.3)
    (Iran, 3.6)
    (EAU, 3.2)
    (Koweït, 2.7)
};
\addlegendentry{Membres OPEP (Production 2023)}
\addplot[poporange!70, fill=poporange!60] coordinates {
    (Qatar, 1.3)
};
\addlegendentry{Hors OPEP (Qatar)}
% Ligne de référence moyenne OPEP par pays membre (optionnel)
%\addplot[popred, dashed, thick, forget plot] coordinates {(Arabie Saoudite, 5.5) (Qatar, 5.5)};
\end{axis}
\end{tikzpicture}
\legende{Diagramme en barres corrigé : Production pétrolière des pays du Golfe en 2023. Les 5 membres de l'OPEP (bleu) dominent ; le Qatar (orange) a quitté l'organisation en 2019. L'Arabie Saoudite hégémonise la production (9,0 Mb/j). Source : OPEP Annual Statistical Bulletin 2024, BP Statistical Review 2024.}
\end{popfigure}

\begin{exemplebox}[Analyse commentée du diagramme (modèle de restitution orale)]
« Ce diagramme met en évidence l'hégémonie de l'\textbf{Arabie Saoudite} (9,0 Mb/j), qui produit à elle seule près de \textbf{33\%} de la production des 5 membres du Golfe à l'OPEP (total ~22,8 Mb/j) et reste le \textit{swing producer} mondial. L'Iran, malgré des réserves immenses (209 Mrds de barils), est pénalisé par les sanctions internationales (production 3,6 Mb/j bien en deçà de son potentiel). Le Qatar (1,3 Mb/j), bien que hors OPEP depuis 2019, reste un acteur gazier majeur (GNL). Le diagramme illustre ainsi la \textbf{vulnérabilité géopolitique} de l'approvisionnement mondial : une crise dans le Golfe (blocus d'Ormuz, attaque d'installations saoudiennes) fait basculer l'économie mondiale. »
\end{exemplebox}

\courssub{TP 3 — Frise chronologique double entrée : « Guerres / Initiatives de paix » (1947–2024)}

\paragraph{Consigne.} Sur une feuille A4 paysage (ou une frise numérique), tracez une ligne du temps centrale graduée de 1947 à 2024 (pas de 5 ou 10 ans). Au-dessus : \textbf{Guerres et violences majeures} (barres rouges, durée). En-dessous : \textbf{Initiatives diplomatiques et accords de paix} (barres vertes, dates ponctuelles). Utilisez des codes couleurs et symboles distincts. Reportez les événements suivants (liste non exhaustive) :

\begin{center}
\begin{tabularx}{\linewidth}{|l|X|l|X|}\hline
\textbf{Période} & \textbf{GUERRES / VIOLENCES (Haut)} & \textbf{Année} & \textbf{PAIX / DIPLOMATIE (Bas)} \\ \hline
1947–1949 & Guerre de 1948 (Indépendance / Nakba) & 1947 & Résolution 181 (Partage ONU) \\ \hline
1956 & Crise de Suez & 1949 & Accords d'armistice (Rhodés) \\ \hline
1967 & Guerre des Six Jours (Juin) & 1967 & Résolution 242 (Territoires vs Paix) \\ \hline
1967–1970 & Guerre d'usure & 1973 & Conférence de Genève (échec) \\ \hline
1973 & Guerre du Kippour / Ramadan (Oct.) & 1978 & Accords de Camp David I (Égypte-Israël) \\ \hline
1975–1990 & Guerre civile du Liban & 1979 & Traité de paix Égypte-Israël \\ \hline
1982 & Guerre du Liban (Invasion israélienne) & 1988 & Fondation Hamas / Charte (radicalisation) \\ \hline
1987–1993 & Première Intifada & 1991 & Conférence de Madrid \\ \hline
1990–1991 & Guerre du Golfe (Koweït) & 1993 & Accords d'Oslo I (Reconnaissance mutuelle) \\ \hline
2000–2005 & Deuxième Intifada (Al-Aqsa) & 1994 & Traité de paix Jordanie-Israël \\ \hline
2006 & Guerre du Liban (Hezbollah) & 2000 & Sommet de Camp David II (Échec) \\ \hline
2008–2009 & Opération « Plomb durci » (Gaza) & 2002 & Initiative de paix arabe (Beyrouth) \\ \hline
2014 & Opération « Bordure protectrice » (Gaza) & 2007 & Conférence d'Annapolis (Échec) \\ \hline
2023–... & Guerre Israël-Hamas (7 oct. 2023) & 2020 & Accords d'Abraham (EAU, Bahreïn, Maroc, Soudan) \\ \hline
\end{tabularx}
\end{center}

\begin{aretenir}[Clé de lecture de la frise double entrée]
La superposition des deux lignes révèle une \textbf{asymétrie structurelle} : les phases de guerre sont longues, récurrentes et destructrices (barres rouges épaisses), tandis que les initiatives de paix sont ponctuelles, fragiles et souvent suivies d'une reprise du conflit (barres vertes fines, espacées). Seuls les accords Égypte-Israël (1979) et Jordanie-Israël (1994) ont abouti à des traités de paix durables \textit{bilatéraux}. Le cœur du conflit (palestinien) reste non résolu malgré Oslo (1993). Les Accords d'Abraham (2020) normalisent les relations avec des pays arabes \textit{sans} résoudre la question palestinienne.
\end{aretenir}

\courssub{Dossier du programme : Étude de documents croisés}

\paragraph{Consigne.} Étudiez les quatre documents ci-après. Répondez aux questions de mise en relation pour construire une explication synthétique : \textbf{« Comment la communauté internationale a-t-elle tenté de gérer la question palestinienne depuis 1917, et pourquoi le sort des réfugiés reste-t-il le point d'achoppement central ? »}

\begin{definition}[Document 1 — Extrait de la Déclaration Balfour (2 nov. 1917)]
\textit{« Le gouvernement de Sa Majesté voit d'un favorable œil l'établissement en Palestine d'un foyer national pour le peuple juif, et s'efforcera de faciliter la réalisation de cet objectif, il étant bien entendu que rien ne sera fait qui puisse porter préjudice aux droits civils et religieux des collectivités non juives existantes en Palestine, ni aux droits et au statut politique dont jouissent les Juifs dans tout autre pays. »} — Arthur Balfour, Secrétaire d'État britannique aux Affaires étrangères, adressé à Lord Rothschild.
\end{definition}

\begin{definition}[Document 2 — Extrait de la Résolution 181 (II) de l'Assemblée Générale de l'ONU, 29 nov. 1947]
\textit{« L'Assemblée générale [...] Recommande au Royaume-Uni, puissance mandataire pour la Palestine, et à tous les autres États membres des Nations Unies l'adoption et la mise en œuvre, en ce qui concerne le gouvernement futur de la Palestine, du Plan de partage avec union économique [...] : un État arabe, un État juif, et un régime international spécial pour la ville de Jérusalem. »}
\end{definition}

\begin{definition}[Document 3 — Photographie : Le mur de séparation à Bethléem (2010s)]
\textit{[Description de l'image à analyser :] Un mur de béton gris de 8 mètres de haut, surmonté de miradors et de fil barbelé, traverse un paysage urbain dense. De l'un côté, des habitations palestiniennes serrées, une tour de guet militaire. De l'autre, une route israélienne lisse. Sur le mur, des graffitis en anglais et arabe : « Justice », « Freedom », « Ich bin ein Berliner » (référence au mur de Berlin), une colombe percée d'une balle. Au premier plan, un enfant palestinien passe devant une fresque géante représentant Leila Khaled (icône du FPLP).}
\end{definition}

\begin{definition}[Document 4 — Données UNRWA : Réfugiés palestiniens enregistrés (janvier 2024)]
\begin{center}
\begin{tabularx}{\linewidth}{|l|X|X|}\hline
\textbf{Champ d'opération UNRWA} & \textbf{Réfugiés enregistrés (millions)} & \textbf{Part du total (\%)} \\ \hline
Jordanie & 2,4 & 40\% \\ \hline
Bande de Gaza & 1,5 & 25\% \\ \hline
Cisjordanie & 0,9 & 15\% \\ \hline
Syrie & 0,6 & 10\% \\ \hline
Liban & 0,5 & 8\% \\ \hline
\textbf{Total} & \textbf{5,9} & \textbf{100\%} \\ \hline
\end{tabularx}
\end{center}
\textit{Note : L'UNRWA (Office de secours et de travaux des Nations Unies pour les réfugiés de Palestine dans le Proche-Orient) créé en 1949 (Rés. 302). Définition du réfugié : « Personnes dont le lieu de résidence normal était la Palestine entre juin 1946 et mai 1948, qui ont perdu leurs maisons et leurs moyens de subsistance à la suite du conflit de 1948 », et leurs descendants.}
\end{definition}

\begin{exemplebox}[Données chiffrées clés : Le drame des réfugiés (UNRWA)]
L'UNRWA recense \textbf{plus de 5,9 millions de réfugiés palestiniens enregistrés} (chiffre 2024), répartis sur 5 champs d'opération. \textbf{La Jordanie accueille 40\%} d'entre eux (majorité naturalisée mais statut UNRWA maintenu). \textbf{Gaza concentre 25\%} dans une densité extrême (8 camps sur 365 km²). Le \textbf{droit au retour (Rés. 194, déc. 1948)} est la revendication centrale palestinienne ; Israël y oppose le « droit au retour » en Israël = fin de l'État juif (argument démographique). Ce contentieux alimente l'impasse d'Oslo (reporté aux « négociations de statut final » jamais tenues).
\end{exemplebox}

\begin{exercice}[Questions guidées pour le dossier croisé]
\begin{enumerate}
    \item \textbf{Doc 1 \& 2 :} En quoi la Déclaration Balfour (1917) et la Résolution 181 (1947) traduisent-elles une \textbf{évolution du cadre juridique} (promesse unilatérale $\rightarrow$ décision multilatérale ONU) mais une \textbf{continuité du paradoxe} (droit au foyer national vs droits des populations autochtones) ?
    \item \textbf{Doc 2 \& 3 :} Montrez que le mur de séparation (Doc 3) contredit l'esprit du Plan de partage (Doc 2) et des Résolutions 242/338. Quel \textbf{fait colonial} (colonies, annexion de Jérusalem-Est) le mur vient-il consolider sur le terrain ?
    \item \textbf{Doc 4 :} Calculez le nombre de réfugiés vivant \textbf{hors des territoires palestiniens} (Jordanie + Syrie + Liban). Quel pourcentage cela représente-t-il ? En déduire la dimension \textbf{régionale} (externe) de la question des réfugiés, au-delà du seul territoire israélo-palestinien.
    \item \textbf{Synthèse (15 lignes max) :} Rédigez une réponse argumentée à la problématique : « Comment la communauté internationale a-t-elle tenté de gérer la question palestinienne depuis 1917, et pourquoi le sort des réfugiés reste-t-il le point d'achoppement central ? » \textit{Utilisez les 4 documents et vos connaissances (Rés. 194, Oslo, Camp David, Accords d'Abraham).}
\end{enumerate}
\end{exercice}

\begin{corrige}[Exercice n° Dossier croisé — Éléments de réponse attendus]
\begin{enumerate}
    \item \textbf{Cadre juridique :} Balfour = acte unilatéral puissance mandataire (GB), engageant sa responsabilité morale/politique sans base de droit international contraignant. Rés. 181 = décision multilatérale (ONU, 33 voix pour), tentative de partage légal du territoire. \textbf{Paradoxe continu :} Les deux textes promettent un « foyer/État juif » tout en mentionnant la protection des « droits civils/religieux » (Balfour) ou en créant un État arabe distinct (ONU). Dans les deux cas, la \textbf{réalisation du projet sioniste se fait sur une terre habitée} (90\% arabes en 1917, 67\% en 1947), sans consultation effective des populations (rejet arabe de la 181).
    \item \textbf{Mur vs Droit international :} Le mur (Doc 3) ne suit pas la Ligne Verte (frontière 1967, base Rés. 242) mais s'enfonce en Cisjordanie (85\% de son tracé) pour inclure les grands blocs de colonies (Ariel, Maale Adoumim, Gush Etzion) et Jérusalem-Est annexée. Il \textbf{annexe de facto} 10\% de la Cisjordanie, fragmente le territoire palestinien (bantoustans), sépare les paysans de leurs terres. C'est une \textbf{frontière unilatérale} imposée par la force, contraire à l'interdiction d'acquisition de territoire par la force (Charte ONU, Rés. 242, Avis CIJ 2004).
    \item \textbf{Calcul réfugiés externes :} Jordanie (2,4) + Syrie (0,6) + Liban (0,5) = \textbf{3,5 millions}. Pourcentage : 3,5 / 5,9 $\approx$ \textbf{59\%}. \textbf{Conséquence :} La majorité des réfugiés vit dans les pays voisins. Leur sort conditionne la stabilité régionale (Jordanie : équilibre démographique fragile ; Liban : camps hors contrôle étatique, milices ; Syrie : guerre civile, destruction camps comme Yarmouk). Toute solution (retour, indemnisation, réinstallation, naturalisation) nécessite l'accord de ces États hôtes $\rightarrow$ dimension \textbf{régionale incontournable}.
    \item \textbf{Synthèse type :} Depuis 1917, la communauté internationale (GB puis ONU) a tenté de gérer le conflit par le \textbf{partage territorial} (Rés. 181) puis la \textbf{résolution du contentieux réfugiés/territoires} (Rés. 194, 242, 338, Processus d'Oslo 1993). Mais l'\textbf{asymétrie de puissance} (force militaire israélienne, soutien US) a permis la colonisation continue (Doc 3, mur) et le refus du droit au retour (Doc 4, 5,9 M). Les accords bilatéraux (Égypte 1979, Jordanie 1994, Abraham 2020) ont normalisé les relations arabes-israéliennes \textit{sans} résoudre le noyau dur palestinien (État, Jérusalem, réfugiés). L'impasse persiste car la \textbf{légalité internationale (ONU, CIJ 2004 sur le mur)} reste sans contrainte effective sur le terrain.
\end{enumerate}
\end{corrige}

\courssub{Restitution orale et évaluation des compétences}

\paragraph{Modalités.} Chaque groupe (3–4 élèves) présente \textbf{une} de ses productions (Croquis TP1, Diagramme TP2, Frise TP3, ou Synthèse Dossier) en \textbf{3 minutes maximum}, suivie de \textbf{2 minutes de questions} de la classe et du professeur.

\begin{methode}[Réussir sa restitution orale (Savoir-faire : Rédiger et justifier une démarche)]
\begin{enumerate}
    \item \textbf{Accroche (30 s) :} Annoncer le titre exact de la production, la problématique traitée, le choix des échelles/données.
    \item \textbf{Démonstration structurée (2 min) :} Suivre la logique de la légende (croquis), des barres (diagramme), des entrées (frise) ou des documents (dossier). \textbf{Justifier chaque choix graphique :} « J'ai mis cette couleur pour... », « J'ai regroupé ces pays car... », « J'ai placé cet événement ici car il marque une rupture... ».
    \item \textbf{Conclusion synthétique (30 s) :} Répondre explicitement à la question de départ : « Ce croquis montre que... », « Ce diagramme prouve l'hégémonie de... », « Cette frise révèle l'échec structurel des initiatives de paix... ».
    \item \textbf{Posture :} Regarder l'audience, voix posée, vocabulaire précis (sémiologie, géopolitique, démographie), notes interdites (sauf plan détaillé).
\end{enumerate}
\end{methode}

\begin{aretenir}[Grille d'auto-évaluation des TP (à remplir en fin de séance)]
\begin{center}
\begin{tabularx}{\linewidth}{|l|X|X|X|}\hline
\textbf{Critère} & \textbf{Maîtrisé (3)} & \textbf{En acquisition (2)} & \textbf{Non acquis (1)} \\ \hline
\textbf{Rigueur graphique} (titre, échelle, Nord, légende hiérarchisée, propreté) & & & \\ \hline
\textbf{Pertinence du contenu} (choix des infos, exactitude dates/chiffres, exhaustivité) & & & \\ \hline
\textbf{Qualité de l'analyse} (liens cause/effet, mise en perspective, réponse à la probl.) & & & \\ \hline
\textbf{Expression orale} (structure, justifications, vocabulaire, temps de parole) & & & \\ \hline
\end{tabularx}
\end{center}
\textit{Objectif :} Atteindre au moins 10/12 pour valider les compétences « Appliquer les notions » et « Rédiger et justifier » sur ce chapitre.
\end{aretenir}

\begin{savaistu}[Le Cameroun et l'UNRWA : une solidarité concrète]
Le Cameroun, membre de l'ONU depuis 1960, contribue régulièrement au budget de l'UNRWA (contributions volontaires). Des cadres camerounais (médecins, administrateurs, logisticiens) ont servi dans les missions de l'UNRWA à Gaza, en Cisjordanie, au Liban ou en Jordanie. Par ailleurs, la position diplomatique du Cameroun à l'ONU (vote résolutions, comité pour l'exercice des droits inaliénables du peuple palestinien) s'inscrit dans la continuité de son engagement panafricain et non-aligné pour le droit des peuples à l'autodétermination. Comprendre les chiffres de l'UNRWA (Doc 4), c'est aussi mesurer l'ampleur de l'action humanitaire internationale dont le Cameroun est un acteur, modeste mais réel.
\end{savaistu}

\newpage

\coursec{Activité d'intégration}

\coursec{Le Proche et le Moyen-Orient : zones de tension sans fin}

\courssub{Objectifs de l'activité}

Cette activité d'intégration mobilise l'ensemble des savoirs et savoir-faire de la leçon \og Le Proche et le Moyen-Orient : zones de tension sans fin \fg{}. À l'issue de l'activité, l'élève doit être capable de :
\begin{itemize}
\item localiser sur un croquis les enjeux stratégiques du Proche et du Moyen-Orient (gisements pétroliers, détroits, foyers de conflit) ;
\item construire et exploiter un diagramme en barres à partir de données chiffrées sur la production pétrolière ;
\item mobiliser la chronologie des conflits (1947-2023) pour expliquer la permanence des tensions ;
\item rédiger un texte argumenté de 15 à 20 lignes, en justifiant chaque argument par un document du dossier (savoir-faire : rédiger et justifier une démarche, une réponse ou une conclusion) ;
\item appliquer les notions de l'unité à une situation concrète de la vie courante : la hausse du prix du carburant à la pompe au Cameroun.
\end{itemize}

\courssub{Situation}

\textbf{Le baril de la discorde : du golfe Persique à la pompe de Yaoundé.}

En février 2024, dans un contexte de tensions persistantes au Proche-Orient depuis le déclenchement de la guerre entre Israël et le Hamas en octobre 2023, le prix du litre de super à la station d'Akwa (Douala) passe d'environ \num{730} à \num{840} FCFA. À Yaoundé, les conducteurs de taxis augmentent la course de 200 à 300 FCFA, et les commerçantes du marché Mokolo répercutent la hausse sur le prix des denrées transportées depuis l'Ouest. Le soir même, un journaliste de la CRTV explique au journal de 20 h : \og Cette hausse est une conséquence directe de la guerre au Proche-Orient, région qui fournit une grande partie du pétrole mondial. \fg{}

Dans la salle de Terminale du lycée technique de Nkol-Eton, les élèves s'interrogent : comment une guerre située à plus de \num{4000} km du Cameroun peut-elle faire monter le prix du carburant à la pompe de Yaoundé ? Pourquoi cette région du monde est-elle en permanence au cœur des tensions internationales ?

Pour répondre, la classe est répartie en groupes de 4 élèves. Chaque groupe reçoit un dossier de 5 documents :

\begin{itemize}
\item \textbf{Document 1} : une carte du Proche et du Moyen-Orient indiquant les principaux gisements pétroliers (Arabie Saoudite, Irak, Koweït, Iran, Émirats arabes unis) et les détroits stratégiques (Ormuz, Suez, Bab el-Mandeb).
\item \textbf{Document 2} : un tableau des réserves et de la production pétrolière du Golfe (données simplifiées, 2022) :

\begin{center}
\begin{tabularx}{\linewidth}{|l|X|X|}
\hline
\rowcolor{popblueL} \textbf{Pays} & \textbf{Réserves prouvées (milliards de barils)} & \textbf{Production (millions de barils/jour)} \\
\hline
Arabie Saoudite & 267 & 10,5 \\
\hline
Iran & 209 & 3,8 \\
\hline
Irak & 145 & 4,5 \\
\hline
Koweït & 102 & 2,7 \\
\hline
Émirats arabes unis & 113 & 3,4 \\
\hline
\end{tabularx}
\end{center}

\item \textbf{Document 3} : une frise des conflits de 1947 à 2023 (plan de partage de 1947 ; guerres de 1948, 1956, 1967, 1973 ; Intifadas de 1987 et 2000 ; guerre Iran-Irak 1980-1988 ; guerres du Golfe 1991 et 2003 ; guerre Israël-Hamas de 2023).
\item \textbf{Document 4} : une photographie du mur de séparation en Cisjordanie, construit par Israël à partir de 2002.
\item \textbf{Document 5} : un extrait d'article de presse camerounaise : \og Le Cameroun importe l'essentiel de ses produits pétroliers raffinés. Chaque flambée du cours du baril sur le marché mondial, fixé en grande partie selon l'offre du Moyen-Orient, se traduit par une hausse des prix à la pompe et une augmentation du coût des transports et des denrées. \fg{}
\end{itemize}

\courssub{Consignes}

\textbf{Travail par groupes de 4, puis restitution orale et correction collective au tableau.}

\begin{enumerate}
\item À l'aide du document 1, construire un \cle{croquis} simple intitulé \og Un espace stratégique \fg{} montrant : les gisements pétroliers, les trois détroits (Ormuz, Suez, Bab el-Mandeb) et les foyers de conflit. Ne pas oublier le titre, la légende et l'orientation.
\item À l'aide du document 2, construire un \cle{diagramme en barres} représentant la production pétrolière journalière des cinq pays du Golfe. Calculer ensuite la part de l'Arabie Saoudite dans la production totale des cinq pays (en \%).
\item À l'aide du document 3, dégager \textbf{deux} raisons pour lesquelles on peut parler de tensions \og sans fin \fg{} dans la région.
\item À l'aide des documents 4 et 5, expliquer en quelques lignes le lien entre le conflit israélo-palestinien, le marché mondial du pétrole et la vie quotidienne des Camerounais.
\item \textbf{Tâche finale} : rédiger un texte argumenté de 15 à 20 lignes répondant à la question : \og Pourquoi le Proche et le Moyen-Orient sont-ils des zones de tension sans fin, et quelles conséquences pour le Cameroun ? \fg{} Chaque argument doit être justifié par un document du dossier, cité entre parenthèses (doc. 1, doc. 2, etc.).
\end{enumerate}

\courssub{Ressources à mobiliser}

\begin{itemize}
\item \textbf{Savoirs} : les origines du problème israélo-palestinien (sionisme, déclaration Balfour de 1917, plan de partage de 1947, création d'Israël en 1948) ; les guerres israélo-arabes et les Intifadas ; les conflits du Moyen-Orient et du golfe Persique (guerre Iran-Irak, guerres du Golfe) ; le rôle stratégique du pétrole et des détroits.
\item \textbf{Savoir-faire} : construire un croquis avec titre, légende et orientation ; construire un diagramme en barres ; calculer une part en pourcentage :
\[ \text{Part (\%)} = \frac{\text{valeur de la partie}}{\text{valeur du total}} \times 100 \]
\item \textbf{Méthode du texte argumenté} : introduction (rappel du sujet), développement (arguments justifiés par les documents), conclusion (réponse claire à la question posée).
\end{itemize}

\courssub{Démarche et corrigé commenté}

\begin{exoresolu}[Corrigé commenté de l'activité]
\textbf{Tâches a) et b) :} construire le croquis stratégique et le diagramme en barres de la production pétrolière, puis calculer la part de l'Arabie Saoudite. \textbf{Tâche c) :} rédiger le texte argumenté de 15 à 20 lignes.
\tcblower
\textbf{Étape 1 : le croquis (document 1).} Le croquis doit faire apparaître trois idées-forces : la concentration du pétrole autour du golfe Persique, les passages obligés (détroits) et la superposition des foyers de conflit.

\begin{popfigure}
\begin{tikzpicture}[scale=0.92, every node/.style={font=\small}]
% Mers
\fill[popblue!25] (6.2,0.3) .. controls (7.4,0.8) and (7.6,2.6) .. (6.6,3.6)
 .. controls (5.8,4.4) and (5.2,3.2) .. (5.4,2.0) .. controls (5.6,1.0) and (5.6,0.2) .. (6.2,0.3);
\node[popblue] at (6.6,2.0) {\textbf{Golfe Persique}};
\fill[popblue!25] (-0.4,3.4) rectangle (2.2,4.4);
\node[popblue] at (0.9,3.9) {\textbf{Méditerranée}};
\fill[popblue!25] (1.6,0.2) rectangle (2.4,2.6);
\node[popblue,rotate=90] at (2.0,1.4) {\textbf{Mer Rouge}};
% Pays producteurs (jaune/or)
\draw[fill=popgold!40] (2.6,1.0) rectangle (5.2,3.0);
\node at (3.9,2.0) {\textbf{Arabie Saoudite}};
\draw[fill=popgold!40] (3.2,3.2) rectangle (5.2,4.4);
\node at (4.2,3.8) {\textbf{Irak}};
\draw[fill=popgold!40] (6.8,3.8) rectangle (8.6,5.0);
\node at (7.7,4.4) {\textbf{Iran}};
\draw[fill=popgold!40] (5.4,3.2) rectangle (6.6,4.0); % Koweït approx
\node[font=\footnotesize] at (6.0,3.6) {\textbf{Koweït}};
\draw[fill=popgold!40] (7.0,1.8) rectangle (8.2,2.6); % EAU approx
\node[font=\footnotesize] at (7.6,2.2) {\textbf{Émirats}};
% Foyer israélo-palestinien (orange)
\draw[fill=poporange!40] (1.4,2.8) rectangle (2.4,3.4);
\node[font=\footnotesize] at (1.9,3.1) {Israël/Pal.};
% Égypte (contexte Suez)
\draw[fill=popgold!20] (0.2,2.0) rectangle (1.4,3.2);
\node[font=\footnotesize] at (0.8,2.6) {Égypte};
% Détroits (rose)
\fill[poppink] (6.4,0.35) circle (0.09); \node[right,font=\footnotesize] at (6.5,0.35) {Détroit d'Ormuz};
\fill[poppink] (1.5,3.3) circle (0.09); \node[left,font=\footnotesize] at (1.4,3.5) {Canal de Suez};
\fill[poppink] (2.0,0.1) circle (0.09); \node[below,font=\footnotesize] at (2.0,-0.05) {Bab el-Mandeb};
% Gisements (triangles noirs)
\foreach \p in {(3.2,1.4),(4.6,1.6),(4.0,3.5),(7.4,4.1),(5.0,2.6),(6.0,3.6),(7.6,2.2)} \node[popdark] at \p {\textbf{$\blacktriangle$}};
% Nord
\draw[->,thick,popdark] (8.8,4.6) -- (8.8,5.4) node[above]{N};
\end{tikzpicture}
\legende{Croquis attendu : en jaune, les grands pays producteurs (les triangles noirs figurent les gisements) ; en rose, les trois détroits stratégiques ; en orange, le foyer israélo-palestinien. Le croquis est volontairement schématique.}
\end{popfigure}

\textbf{Étape 2 : le diagramme en barres (document 2).} On reporte la production journalière de chaque pays. \emph{Attention : dans le code \texttt{pgfplots}, les nombres décimaux s'écrivent avec un point.}

\begin{popfigure}
\begin{tikzpicture}
\begin{axis}[ybar, width=0.85\linewidth, height=6cm, bar width=0.7cm,
 symbolic x coords={Arabie S., Iran, Irak, Koweït, Émirats},
 xtick=data, ymin=0, ymax=12,
 ylabel={Production (millions de barils/jour)},
 nodes near coords, every node near coord/.append style={font=\footnotesize},
 x tick label style={font=\footnotesize}]
\addplot[fill=popblue] coordinates {(Arabie S.,10.5) (Iran,3.8) (Irak,4.5) (Koweït,2.7) (Émirats,3.4)};
\end{axis}
\end{tikzpicture}
\legende{Production pétrolière journalière des cinq grands producteurs du Golfe (données simplifiées, 2022).}
\end{popfigure}

Calcul de la part de l'Arabie Saoudite (les nombres sont saisis avec \texttt{\textbackslash num} pour respecter la typographie française en mode math) :
\begin{align*}
\text{Total} &= \num{10,5} + \num{3,8} + \num{4,5} + \num{2,7} + \num{3,4} = \num{24,9} \text{ millions de barils/jour} \\
\text{Part} &= \frac{\num{10,5}}{\num{24,9}} \times 100 \approx \num{42,2} \,\%
\end{align*}
L'Arabie Saoudite fournit à elle seule environ \num{42,2} \% de la production des cinq pays : toute crise dans le Golfe pèse donc lourd sur l'offre mondiale.

\textbf{Étape 3 : deux raisons des tensions \og sans fin \fg{} (document 3).} Premièrement, la succession ininterrompue des guerres depuis 1948 (1948, 1956, 1967, 1973, puis Intifadas de 1987 et 2000, guerre de 2023) montre qu'aucun règlement durable du conflit israélo-palestinien n'a été trouvé. Deuxièmement, les conflits s'étendent à tout le Moyen-Orient (guerre Iran-Irak 1980-1988, guerres du Golfe 1991 et 2003), car le pétrole attire les rivalités régionales et les interventions des grandes puissances.

\textbf{Étape 4 : le lien avec le Cameroun (documents 4 et 5).} Le mur de Cisjordanie (doc. 4) illustre l'échec de la paix entre Israéliens et Palestiniens ; or chaque flambée de violence inquiète les marchés pétroliers. Comme le Cameroun importe ses produits raffinés (doc. 5), la hausse du baril mondial se répercute à la pompe, puis sur les transports et les denrées.

\textbf{Étape 5 : texte argumenté (modèle de rédaction).}

\emph{Le Proche et le Moyen-Orient sont des zones de tension sans fin pour trois raisons principales. D'abord, le conflit israélo-palestinien, né du partage de la Palestine en 1947 et de la création d'Israël en 1948, oppose deux peuples revendiquant le même territoire ; aucune guerre (1948, 1967, 1973) ni aucun soulèvement (Intifadas de 1987 et 2000) n'a réglé le problème, et le mur de séparation construit en Cisjordanie depuis 2002 symbolise cet échec (doc. 3 et 4). Ensuite, la région concentre des richesses pétrolières immenses : les cinq pays du Golfe produisent environ 24,9 millions de barils par jour, dont 42 \% pour la seule Arabie Saoudite (doc. 2). Ce pétrole doit emprunter des détroits étroits — Ormuz, Suez, Bab el-Mandeb — faciles à menacer en cas de guerre (doc. 1). Enfin, ces richesses provoquent des conflits régionaux (guerre Iran-Irak, guerres du Golfe) et des interventions étrangères qui entretiennent l'instabilité (doc. 3). Les conséquences atteignent le Cameroun : importateur de produits pétroliers raffinés, il subit chaque flambée du baril, comme en février 2024 où le litre de super est passé d'environ 730 à 840 FCFA à Douala, entraînant la hausse des transports et des denrées (doc. 5). Ainsi, tant que la question palestinienne restera irrésolue et que le monde dépendra du pétrole du Golfe, cette région restera une poudrière dont les explosions se font sentir jusqu'à la pompe de Yaoundé.}
\end{exoresolu}

\courssub{Bilan de l'activité}

Cette activité a permis de mettre en évidence trois acquis essentiels de la leçon. D'une part, les élèves ont vérifié concrètement, par le croquis et le diagramme en barres, que la concentration du pétrole autour du golfe Persique et l'étroitesse des détroits font de cette région un espace \cle{stratégique} mondial. D'autre part, la frise 1947-2023 a montré que les conflits se succèdent sans solution durable : guerres israélo-arabes, Intifadas, guerres du Golfe, guerre de 2023. Enfin, le texte argumenté a relié ces tensions lointaines à la vie quotidienne des Camerounais : la hausse du carburant à Akwa ou à Yaoundé est une conséquence directe de l'instabilité du Proche et du Moyen-Orient. Les savoir-faire du programme — construire des documents graphiques, rédiger et justifier une réponse, appliquer les notions de l'unité à une situation concrète — ont ainsi été mobilisés dans une démarche complète, de l'analyse des documents à la restitution orale.

\newpage

\coursec{Méthodes et exercices résolus}

\courssub{Méthode 1 : Appliquer les notions de l'unité à une situation concrète}

\begin{methode}[Mobiliser les notions du cours face à une situation de la vie courante]
Au Probatoire et au Baccalauréat technique, le sujet peut présenter une \cle{situation concrète} (un fait d'actualité, une hausse des prix, une crise diplomatique) et demander de l'expliquer avec les savoirs de la leçon. Démarche en 5 étapes :
\begin{enumerate}
\item \textbf{Lire la situation} et repérer les mots-clés : lieu (Proche-Orient, golfe Persique), acteurs (Israël, Palestiniens, États arabes, grandes puissances), enjeu (territoire, pétrole, sécurité).
\item \textbf{Identifier la notion du cours} correspondante : problème israélo-palestinien, guerre du Kippour et chocs pétroliers, guerre Iran-Irak, guerre du Golfe, ingérences étrangères.
\item \textbf{Faire le lien explicite} entre la situation et la notion : « cette situation illustre... », « elle s'explique par... ». C'est ce lien qui est noté, pas la seule récitation du cours.
\item \textbf{Appuyer avec des faits précis} : dates (1948, 1967, 1973, 1991), accords (Oslo 1993), acteurs (OLP, OPEP).
\item \textbf{Conclure} en revenant à la situation de départ pour montrer qu'elle est comprise.
\end{enumerate}
\textbf{Réflexe à retenir :} Situation $\rightarrow$ Notion $\rightarrow$ Faits $\rightarrow$ Retour à la situation. Ne jamais réciter le cours sans l'accrocher à la situation proposée.
\end{methode}

\begin{exoresolu}[Expliquer une hausse du prix du carburant à Yaoundé]
\textbf{Énoncé.} En octobre 1973, le prix du baril de pétrole quadruple en quelques mois. Des décennies plus tard, chaque flambée de tension au Moyen-Orient se traduit par une hausse du prix du carburant à la pompe, y compris au Cameroun. À l'aide des notions de l'unité, expliquez pourquoi un conflit du golfe Persique peut affecter la vie quotidienne d'un élève camerounais. (Compétence : appliquer les notions à une situation concrète.)
\tcblower
\textbf{Solution détaillée.}
\begin{enumerate}
\item \textbf{Mots-clés de la situation :} pétrole, golfe Persique, hausse des prix, Cameroun. La notion du cours mobilisée est celle des \cle{conflits du Moyen-Orient et du golfe Persique} et de leur dimension économique mondiale.
\item \textbf{Notion identifiée :} le golfe Persique concentre une part essentielle des réserves mondiales de pétrole. La région est traversée par des conflits majeurs : guerre du Kippour (octobre 1973), guerre Iran-Irak (1980-1988), guerre du Golfe (1990-1991).
\item \textbf{Lien explicite :} lors de la guerre du Kippour (1973), les pays arabes producteurs réunis dans l'\cle{OPEP} utilisent le pétrole comme \emph{arme économique} : ils décrètent un embargo et augmentent les prix pour faire pression sur les pays occidentaux soutenant Israël. Le prix du baril est multiplié par quatre : c'est le \cle{premier choc pétrolier}.
\item \textbf{Application au Cameroun :} le pétrole est une marchandise dont le prix se fixe sur le marché mondial. Toute tension dans le golfe Persique (menace sur le détroit d'Ormuz, guerre, embargo) fait craindre une pénurie et fait monter le prix du baril partout. Le Cameroun, qui importe une partie de ses produits pétroliers raffinés, subit cette hausse : le carburant coûte plus cher à la pompe, le transport en taxi ou en bus augmente, et les prix des denrées transportées suivent.
\item \textbf{Conclusion :} la hausse du prix du carburant à Yaoundé illustre la mondialisation des conséquences des conflits du Moyen-Orient : une zone de tension lointaine agit directement sur la vie quotidienne camerounaise par le canal du marché mondial du pétrole.
\end{enumerate}
\end{exoresolu}

\courssub{Méthode 2 : Rédiger une réponse construite (dissertation ou question de synthèse)}

\begin{methode}[Construire et justifier une réponse argumentée en Histoire]
Pour une question de synthèse du type « Pourquoi le Proche-Orient est-il une zone de tension sans fin ? », suivre 6 étapes :
\begin{enumerate}
\item \textbf{Analyser le sujet} : dégager les mots-clés (« tension », « sans fin ») et la question posée (des \emph{causes} sont demandées).
\item \textbf{Délimiter} l'espace (Proche et Moyen-Orient, golfe Persique) et le temps (de 1948 à nos jours).
\item \textbf{Mobiliser un plan en 2 ou 3 parties} : par exemple I. Un conflit territorial et national (israélo-palestinien) ; II. Des enjeux économiques et stratégiques (pétrole, position géographique) ; III. Des ingérences extérieures et des rivalités régionales qui prolongent les tensions.
\item \textbf{Rédiger} : chaque partie = idée directrice + faits datés (1948 : création d'Israël et première guerre ; 1967 : guerre des Six Jours et occupation des territoires ; 1987 et 2000 : Intifadas ; 1993 : accords d'Oslo non appliqués ; 1980-1988 : guerre Iran-Irak ; 1991 : guerre du Golfe).
\item \textbf{Justifier chaque affirmation} par un fait précis : une affirmation sans date ni exemple vaut zéro.
\item \textbf{Conclure} en répondant exactement à la question posée, sans ouvrir sur un autre sujet.
\end{enumerate}
\textbf{Réflexe :} une bonne copie d'Histoire = un plan visible + des dates + des acteurs nommés + une conclusion qui répond au sujet.
\end{methode}

\begin{exoresolu}[Question de synthèse : les causes d'une tension permanente]
\textbf{Énoncé.} « Le Proche et le Moyen-Orient : une zone de tension sans fin. » En vous appuyant sur des faits précis, dégagez \textbf{trois causes} qui expliquent la persistance des conflits dans cette région depuis 1948. (Rédiger et justifier votre réponse.)
\tcblower
\textbf{Solution détaillée (corps de réponse rédigé).}

\textbf{Introduction.} Depuis la création de l'État d'Israël en 1948, le Proche et le Moyen-Orient n'ont jamais connu une paix durable. Trois causes principales expliquent cette persistance des tensions.

\textbf{Première cause : un conflit territorial et national non résolu.} Le partage de la Palestine décidé par l'ONU en 1947 est refusé par les Arabes. La proclamation de l'État d'Israël le 14 mai 1948 déclenche la première guerre israélo-arabe et l'exode de centaines de milliers de Palestiniens (la \emph{Nakba}, la « catastrophe »). En 1967, la guerre des Six Jours permet à Israël d'occuper la Cisjordanie, Gaza, Jérusalem-Est, le Sinaï et le Golan. Les Palestiniens, organisés autour de l'\cle{OLP} (Organisation de libération de la Palestine, créée en 1964), revendiquent un État ; les Intifadas de 1987 et de 2000 montrent que le problème reste entier, malgré les accords d'\cle{Oslo} de 1993.

\textbf{Deuxième cause : des enjeux économiques et stratégiques majeurs.} La région détient l'essentiel des réserves mondiales de pétrole. Le pétrole devient une arme : lors de la guerre du Kippour (1973), l'OPEP quadruple le prix du baril (premier choc pétrolier). Le contrôle des ressources et des voies de passage (détroit d'Ormuz, canal de Suez) alimente en permanence les rivalités.

\textbf{Troisième cause : les rivalités régionales et les ingérences étrangères.} La guerre Iran-Irak (1980-1988) oppose deux puissances régionales ; l'invasion du Koweït par l'Irak de Saddam Hussein provoque la guerre du Golfe (1991) menée par une coalition internationale sous égide américaine. Les grandes puissances (États-Unis, URSS puis Russie) interviennent en soutenant des camps opposés, ce qui internationalise et prolonge chaque conflit.

\textbf{Conclusion.} La tension est « sans fin » parce que trois facteurs se combinent : un problème national palestinien non réglé, la richesse pétrolière convoitée et l'implication constante de puissances extérieures. Chaque paix signée (Oslo, 1993) bute sur ces causes profondes.
\end{exoresolu}

\courssub{Méthode 3 : Commenter un document (texte ou carte)}

\begin{methode}[La méthode du commentaire de document en 5 étapes]
\begin{enumerate}
\item \textbf{Présenter le document} : nature (texte, carte, photographie), auteur ou source, date, thème. Exemple : « carte du partage de la Palestine proposé par l'ONU en 1947 ».
\item \textbf{Décrire avec méthode} : pour une carte, lire le titre, la légende, l'échelle, l'orientation ; décrire du général au particulier. Pour un texte, dégager les idées principales sans recopier.
\item \textbf{Analyser} : dégager ce que le document montre (exemple : un partage jugé injuste car l'État juif reçoit environ 55 \% du territoire alors que les Juifs sont minoritaires dans la population).
\item \textbf{Mettre en relation avec le cours} : le document illustre quelle notion ? (origines du problème israélo-palestinien, guerre du Golfe...).
\item \textbf{Dégager l'intérêt et les limites} : que nous apprend-il ? Que ne dit-il pas (point de vue de l'auteur, date ancienne) ?
\end{enumerate}
\textbf{Formule à retenir :} Présenter $\rightarrow$ Décrire $\rightarrow$ Analyser $\rightarrow$ Relier au cours $\rightarrow$ Conclure.
\end{methode}

\begin{exoresolu}[Commentaire d'un extrait de la résolution 181 de l'ONU]
\textbf{Énoncé.} Document : « L'Assemblée générale des Nations unies recommande le partage de la Palestine en deux États, l'un juif, l'autre arabe, Jérusalem étant placée sous contrôle international. » (Résolution 181, 29 novembre 1947.) Présentez ce document, dégagez son contenu et montrez en quoi il est à l'origine du problème israélo-palestinien.
\tcblower
\textbf{Solution détaillée.}
\begin{enumerate}
\item \textbf{Présentation :} il s'agit d'un extrait de la résolution 181 adoptée par l'Assemblée générale de l'ONU le 29 novembre 1947, au lendemain de la Seconde Guerre mondiale et de la Shoah. C'est un texte diplomatique qui propose un plan de partage de la Palestine, alors sous mandat britannique.
\item \textbf{Contenu :} le document prévoit trois choses : la création d'un État juif, la création d'un État arabe, et un statut international pour Jérusalem, ville sainte des trois religions monothéistes.
\item \textbf{Analyse :} ce partage est accepté par les dirigeants sionistes mais refusé par les Arabes palestiniens et les États arabes, qui jugent injuste d'attribuer environ 55 \% du territoire à une population juive alors minoritaire. La Grande-Bretagne se retire ; le 14 mai 1948, David Ben Gourion proclame l'État d'Israël.
\item \textbf{Relation avec le cours :} ce document est l'acte fondateur du problème israélo-palestinien : dès le lendemain de la proclamation, la première guerre israélo-arabe éclate (1948-1949). L'État arabe prévu par la résolution ne verra jamais le jour, et l'exode palestinien crée le problème des réfugiés, toujours non résolu.
\item \textbf{Conclusion :} la résolution 181 illustre l'origine internationale du conflit : une décision extérieure, appliquée sans l'accord de la population arabe locale, fait naître un contentieux territorial qui nourrit les tensions jusqu'à aujourd'hui.
\end{enumerate}
\end{exoresolu}

\courssub{Méthode 4 : Construire et exploiter une frise chronologique}

\begin{methode}[Bien construire une frise pour mémoriser et réviser]
\begin{enumerate}
\item \textbf{Choisir l'intervalle} : de 1947 (plan de partage) à nos jours.
\item \textbf{Sélectionner 8 à 10 dates essentielles} maximum : trop de dates rend la frise illisible.
\item \textbf{Classer par thème} si possible : conflit israélo-palestinien d'un côté, conflits du Golfe de l'autre.
\item \textbf{Vérifier l'ordre chronologique} et l'exactitude de chaque date.
\item \textbf{Exploiter la frise} : elle doit permettre de repérer les enchaînements (une guerre $\rightarrow$ ses conséquences $\rightarrow$ une tentative de paix).
\end{enumerate}
\textbf{Dates-clés à retenir :} 1947 (résolution 181) ; 1948 (création d'Israël, première guerre israélo-arabe) ; 1956 (crise de Suez) ; 1967 (guerre des Six Jours) ; 1973 (guerre du Kippour, premier choc pétrolier) ; 1978-1979 (accords de Camp David puis traité de paix Égypte-Israël) ; 1980-1988 (guerre Iran-Irak) ; 1987 (première Intifada) ; 1991 (guerre du Golfe) ; 1993 (accords d'Oslo) ; 2000 (seconde Intifada).
\end{methode}

\begin{exoresolu}[Exploiter une frise pour établir un enchaînement de causes et conséquences]
\textbf{Énoncé.} Voici quatre événements : (a) guerre du Kippour, octobre 1973 ; (b) embargo et hausse des prix de l'OPEP ; (c) premier choc pétrolier mondial ; (d) crise économique dans les pays importateurs, dont l'Afrique. Replacez ces événements dans l'ordre chronologique et logique, puis rédigez un court paragraphe justifiant cet enchaînement.
\tcblower
\textbf{Solution détaillée.}
\begin{enumerate}
\item \textbf{Ordre :} (a) $\rightarrow$ (b) $\rightarrow$ (c) $\rightarrow$ (d).
\item \textbf{Justification :} en octobre 1973, l'Égypte et la Syrie attaquent Israël le jour du Kippour pour récupérer les territoires perdus en 1967 (a). Pour punir les pays occidentaux qui soutiennent Israël, les États arabes producteurs de pétrole, réunis dans l'OPEP, décrètent un embargo et multiplient par quatre le prix du baril (b). Cette décision provoque le premier choc pétrolier, c'est-à-dire une brutale envolée du coût de l'énergie à l'échelle mondiale (c). Les pays importateurs, y compris les pays africains comme le Cameroun, voient leurs factures pétrolières exploser : inflation, ralentissement économique et endettement s'ensuivent (d).
\item \textbf{Conclusion :} cet enchaînement montre qu'un conflit militaire localisé au Proche-Orient produit, par l'arme pétrolière, des conséquences économiques planétaires : c'est exactement ce que signifie l'expression « zone de tension » aux répercussions mondiales.
\end{enumerate}
\end{exoresolu}

\courssub{Méthode 5 : Lire et décrire une carte historique}

\begin{methode}[Décrire une carte de la région en 4 étapes]
\begin{enumerate}
\item \textbf{Lire les éléments obligatoires} : titre, légende, échelle, orientation (flèche du nord). Une carte sans titre ni légende lus est une carte non comprise.
\item \textbf{Situer l'espace} : le Proche-Orient (Israël, Palestine, Liban, Syrie, Jordanie, Égypte) autour de la Méditerranée orientale ; le Moyen-Orient et le golfe Persique (Irak, Iran, Arabie Saoudite, Koweït) plus à l'est.
\item \textbf{Décrire du général au particulier} : d'abord l'ensemble (un espace carrefour entre Afrique, Asie et Europe), puis les détails (territoires occupés en 1967, champs pétroliers, détroits stratégiques comme Ormuz et Suez).
\item \textbf{Interpréter} : dégager ce que la carte révèle (position stratégique, concentration du pétrole, morcellement territorial source de conflits).
\end{enumerate}
\textbf{Réflexe :} toujours nommer les mers et les lieux précisément (mer Méditerranée, mer Rouge, golfe Persique, canal de Suez, détroit d'Ormuz).
\end{methode}

\begin{exoresolu}[Exploiter une carte schématique du golfe Persique]
\textbf{Énoncé.} Sur une carte du golfe Persique figurent : l'Irak, l'Iran, le Koweït, l'Arabie Saoudite, le détroit d'Ormuz et des zones pétrolières. Décrivez cette carte et expliquez pourquoi cet espace a été le théâtre de deux guerres majeures entre 1980 et 1991.
\tcblower
\textbf{Solution détaillée.}
\begin{enumerate}
\item \textbf{Description :} le golfe Persique est une mer presque fermée, bordée à l'est par l'Iran, à l'ouest par l'Irak, le Koweït et l'Arabie Saoudite. Il ne communique avec le golfe d'Oman et l'océan Indien que par le \cle{détroit d'Ormuz}, passage étroit et stratégique. Les zones pétrolières se concentrent sur ses rives et ses fonds.
\item \textbf{Interprétation :} cet espace réunit deux facteurs de tension : une richesse pétrolière considérable et un goulot d'étranglement (Ormuz) par lequel transite une grande partie du pétrole mondial.
\item \textbf{Explication des deux guerres :} d'une part, la \cle{guerre Iran-Irak} (1980-1988) éclate lorsque l'Irak de Saddam Hussein attaque l'Iran, notamment pour le contrôle du Chatt el-Arab (fleuve frontalier donnant accès au golfe) et des zones pétrolières ; c'est une guerre de position longue et meurtrière. D'autre part, en août 1990, l'Irak envahit le \cle{Koweït}, petit État riche en pétrole ; une coalition internationale dirigée par les États-Unis, autorisée par l'ONU, libère le Koweït lors de la \cle{guerre du Golfe} (janvier-février 1991).
\item \textbf{Conclusion :} la carte montre que la concentration du pétrole et l'étroitesse des voies de passage font du golfe Persique un espace convoité : qui contrôle ces rives contrôle une part décisive de l'énergie du monde, d'où la succession des conflits.
\end{enumerate}
\end{exoresolu}

\courssub{Méthode 6 : Traiter des données chiffrées (calculer et interpréter)}

\begin{methode}[Calculer un taux de variation et l'interpréter historiquement]
Pour un choc pétrolier ou toute évolution chiffrée :
\begin{enumerate}
\item \textbf{Identifier} la valeur de départ $V_0$ et la valeur d'arrivée $V_1$ (avec leurs dates).
\item \textbf{Calculer le taux de variation} en pourcentage :
\[ t = \frac{V_1 - V_0}{V_0} \times 100 \]
\item \textbf{Vérifier la cohérence} : si le prix quadruple, le taux est de $+300\,\%$ (et non $400\,\%$ : quadrupler, c'est ajouter trois fois la valeur initiale à la valeur de départ).
\item \textbf{Interpréter} le résultat par une phrase historique : qui décide la hausse, pourquoi, avec quelles conséquences.
\end{enumerate}
\textbf{Réflexe :} un calcul non interprété ne rapporte qu'une partie des points ; toujours terminer par le sens historique du chiffre.
\end{methode}

\begin{exoresolu}[Calculer l'ampleur du premier choc pétrolier]
\textbf{Énoncé.} En octobre 1973, le baril de pétrole coûte environ 3 dollars. En janvier 1974, après les décisions de l'OPEP, il atteint environ 12 dollars. 1) Calculez le taux de variation du prix du baril entre octobre 1973 et janvier 1974. 2) Interprétez ce résultat en le reliant aux événements de l'époque et à ses conséquences pour les pays importateurs comme le Cameroun.
\tcblower
\textbf{Solution détaillée.}
\begin{enumerate}
\item \textbf{Données :} $V_0 = 3$ dollars (octobre 1973) ; $V_1 = 12$ dollars (janvier 1974).
\item \textbf{Calcul :}
\[ t = \frac{V_1 - V_0}{V_0} \times 100 = \frac{12 - 3}{3} \times 100 = \frac{9}{3} \times 100 = 3 \times 100 = +300\,\%. \]
Le prix du baril a donc augmenté de \textbf{300 \%} : il a été \emph{multiplié par quatre} ($3 \times 4 = 12$). On vérifie : quadrupler correspond bien à $+300\,\%$, car la hausse (9 dollars) représente trois fois la valeur initiale.
\item \textbf{Interprétation :} cette hausse brutale est le \cle{premier choc pétrolier}. Elle est décidée par les pays arabes de l'OPEP pendant la guerre du Kippour (octobre 1973) pour faire pression sur les pays occidentaux soutenant Israël : le pétrole devient une \emph{arme économique et diplomatique}.
\item \textbf{Conséquences :} pour les pays importateurs, la facture énergétique explose, ce qui provoque inflation et ralentissement économique. Les pays africains, dont le Cameroun, subissent le renchérissement du carburant et des transports, puis s'endettent : un conflit du Proche-Orient pèse ainsi directement sur l'économie camerounaise.
\item \textbf{Conclusion :} le chiffre de $+300\,\%$ mesure concrètement la force de l'arme pétrolière : en quelques mois, les producteurs arabes ont transformé un conflit régional en crise économique mondiale.
\end{enumerate}
\end{exoresolu}

\begin{attention}[Les 5 erreurs qui coûtent des points au Bac]
\begin{enumerate}
\item \textbf{Confondre Proche-Orient et golfe Persique.} Le conflit israélo-palestinien se situe au Proche-Orient (Méditerranée orientale) ; les guerres Iran-Irak et du Golfe se déroulent autour du golfe Persique. Mélanger les deux espaces dans une même partie affaiblit toute la copie.
\item \textbf{Inverser les dates des guerres.} Erreurs classiques : placer la guerre des Six Jours en 1973 (c'est \textbf{1967}) ou la guerre du Kippour en 1967 (c'est \textbf{1973}) ; situer la guerre du Golfe en 1980 (c'est \textbf{1991} ; 1980-1988 est la guerre Iran-Irak). Apprendre la frise par cœur.
\item \textbf{Affirmer sans justifier.} Écrire « la région est une poudrière » sans citer un seul fait daté (1948, 1967, Oslo 1993...) ne rapporte rien : chaque idée doit être appuyée par un fait, une date et un acteur.
\item \textbf{Se tromper dans le calcul du choc pétrolier.} Un prix qui passe de 3 à 12 dollars augmente de $+300\,\%$ (multiplié par 4), et non de $400\,\%$. Oublier de soustraire $V_0$ au numérateur est l'erreur la plus fréquente.
\item \textbf{Réciter le cours sans répondre à la situation.} Quand le sujet part d'une situation concrète (hausse du carburant à Yaoundé, crise diplomatique), il faut explicitement relier la situation à la notion du cours et \emph{revenir à la situation en conclusion} : une copie qui ne parle que du cours, sans jamais citer la situation, perd les points du savoir-faire.
\end{enumerate}
\end{attention}

\newpage

\coursec{Exercices}

\courssub{Je vérifie mes connaissances}

\begin{exercice}[QCM : Les origines du conflit israélo-palestinien]
Parmi les propositions suivantes, une seule est exacte. Recopiez la lettre correspondante.
\begin{enumerate}
    \item La déclaration Balfour (1917) promet :
    \begin{itemize}
        \item[A] La création d'un État arabe uni au Proche-Orient.
        \item[B] L'établissement en Palestine d'un « Foyer national juif » tout en préservant les droits des communautés non juives.
        \item[C] Le partage de la Palestine en deux États indépendants.
        \item[D] Le retrait immédiat des troupes britanniques de Palestine.
    \end{itemize}
    \item Le plan de partage de la Palestine voté par l'ONU en novembre 1947 (Résolution 181) prévoyait :
    \begin{itemize}
        \item[A] La création d'un État unique binational.
        \item[B] La création d'un État juif, d'un État arabe et d'un corpus separatum pour Jérusalem.
        \item[C] L'annexion de la Palestine par la Jordanie.
        \item[D] Le maintien du mandat britannique pour 20 ans.
    \end{itemize}
    \item La guerre des Six Jours (juin 1967) a pour conséquence territoriale majeure :
    \begin{itemize}
        \item[A] La destruction de l'État d'Israël.
        \item[B] L'occupation par Israël de la Cisjordanie, de Gaza, du Sinaï et du Golan.
        \item[C] La signature d'un traité de paix définitif avec l'Égypte.
        \item[D] La création de l'OLP.
    \end{itemize}
    \item Les accords d'Oslo (1993) ont instauré :
    \begin{itemize}
        \item[A] La paix définitive et le tracé des frontières de 1967.
        \item[B] La reconnaissance mutuelle Israël/OLP et la mise en place d'une autonomie palestinienne intérimaire.
        \item[C] L'annexion de Jérusalem-Est par l'ONU.
        \item[D] Le retrait israélien du Liban sud.
    \end{itemize}
\end{enumerate}
\end{exercice}

\begin{exercice}[Vrai ou Faux justifié : Conflits du Golfe et Moyen-Orient]
Pour chaque affirmation, écrivez VRAI ou FAUX et justifiez brièvement votre réponse par une date, un fait ou un acteur précis.
\begin{enumerate}
    \item La guerre Iran-Irak (1980-1988) a été déclenchée par l'invasion du Koweït par l'Irak.
    \item La première guerre du Golfe (1990-1991) a vu une coalition internationale sous égide de l'ONU chasser l'Irak du Koweït.
    \item La révolution iranienne de 1979 a renversé la monarchie du Chah pour instaurer une République islamique dirigée par l'ayatollah Khomeyni.
    \item Les accords d'Abraham (2020) ont normalisé les relations entre Israël et l'Arabie saoudite.
\end{enumerate}
\end{exercice}

\begin{exercice}[Définitions et repères chronologiques]
Définissez avec précision les termes suivants en les situant dans le temps et l'espace (2 à 3 lignes chacun) :
\begin{enumerate}
    \item \cle{Sionisme} (théorisé par Theodor Herzl, congrès de Bâle 1897).
    \item \cle{Nakba} (« catastrophe », 1948).
    \item \cle{OLP} (Organisation de libération de la Palestine, créée en 1964, Yasser Arafat).
    \item \cle{Intifada} (première 1987-1993, seconde 2000-2005).
    \item \cle{Pétrodollars} (recyclage des revenus pétroliers, choc de 1973).
\end{enumerate}
\end{exercice}

\begin{exercice}[Lecture de données : Production pétrolière du Golfe (extrait 2022)]
Le tableau ci-dessous donne la production de pétrole brut (en millions de barils par jour) de quatre pays du Golfe en 2022.
\begin{center}
\begin{tabularx}{\linewidth}{|l|X|}\hline
\textbf{Pays} & \textbf{Production (Mb/j)} \\ \hline
Arabie saoudite & 10,5 \\ \hline
Irak & 4,5 \\ \hline
Koweït & 3,0 \\ \hline
Iran & 2,6 \\ \hline
\end{tabularx}
\end{center}
\begin{enumerate}
    \item Calculez la part de l'Arabie saoudite dans la production totale de ces quatre pays (arrondi au dixième de pourcent).
    \item Classez ces pays par ordre décroissant de production.
    \item Quelle conclusion géopolitique immédiate pouvez-vous tirer sur le poids de l'Arabie saoudite au sein de l'OPEP ?
\end{enumerate}
\end{exercice}

\courssub{Je m'entraîne}

\begin{exercice}[Carte mentale : Les acteurs du conflit israélo-palestinien]
Réalisez une carte mentale (schéma heuristique) structurée autour du nœud central « Conflit israélo-palestinien ». Faites apparaître quatre branches principales :
\begin{enumerate}
    \item \textbf{Acteurs directs} (Israël, Autorité palestinienne, Hamas, Jihad islamique, colons).
    \item \textbf{Acteurs régionaux} (Égypte, Jordanie, Liban, Syrie, Iran, Arabie saoudite, Turquie).
    \item \textbf{Acteurs internationaux} (USA, ONU, UE, Russie, Quartet).
    \item \textbf{Enjeux territoriaux et symboliques} (Jérusalem, réfugiés, frontières 1967, colonies, eau, sécurité).
\end{enumerate}
Utilisez des flèches pour montrer les relations (alliance, opposition, médiation).
\end{exercice}

\begin{exercice}[Chronologie commentée : De 1947 à nos jours]
Complétez la frise chronologique ci-dessous en associant à chaque date l'événement majeur correspondant, puis ajoutez une phrase explicative sur sa conséquence immédiate.
\begin{center}
\begin{tabularx}{\linewidth}{|c|X|X|}\hline
\textbf{Date} & \textbf{Événement} & \textbf{Conséquence immédiate} \\ \hline
1947 & & \\ \hline
1948 & & \\ \hline
1967 & & \\ \hline
1973 & & \\ \hline
1978-1979 & & \\ \hline
1993 & & \\ \hline
2000 & & \\ \hline
2005 & & \\ \hline
\end{tabularx}
\end{center}
\textit{Liste des événements à placer : Guerre du Kippour, Création de l'État d'Israël / Guerre de 1948, Plan de partage ONU (Rés. 181), Accords de Camp David (paix Égypte-Israël), Accords d'Oslo I, Guerre des Six Jours, Retrait israélien de Gaza (plan Sharon), Déclenchement de la Seconde Intifada.}
\end{exercice}

\begin{exercice}[Analyse de document : Extrait de la résolution 242 du Conseil de sécurité de l'ONU (novembre 1967)]
\textit{« Le Conseil de sécurité [...] affirme que l'acquisition de territoire par la guerre est inadmissible et qu'il faut œuvrer pour une paix juste et durable dans laquelle chaque État de la région peut vivre en sécurité. [...] Retrait des forces armées israéliennes des territoires occupés lors du récent conflit ; cessation de tout état de belligérance ; respect et reconnaissance de la souveraineté, de l'intégrité territoriale et de l'indépendance politique de chaque État de la région et de leur droit de vivre en paix à l'intérieur de frontières sûres et reconnues. »}
\begin{enumerate}
    \item Présentez le document (nature, date, auteur, contexte historique immédiat).
    \item Identifiez les deux principes fondamentaux posés par cette résolution.
    \item Expliquez pourquoi l'ambiguïté de la formulation anglaise (« withdrawal from territories » vs français « retrait des territoires ») a nourri le conflit ultérieur.
    \item Montrez en quoi cette résolution reste, encore aujourd'hui, la référence juridique de base pour toute négociation de paix.
\end{enumerate}
\end{exercice}

\begin{exercice}[Mise en relation : Crises du Golfe et économie camerounaise]
À partir de vos connaissances et du document statistique ci-après, répondez aux questions.
\begin{center}
\begin{tabularx}{\linewidth}{|l|X|X|}\hline
\textbf{Année} & \textbf{Prix moyen du baril de Brent (\$)} & \textbf{Prix à la pompe au Cameroun (Super, FCFA/L)} \\ \hline
2003 & 28,8 & 440 \\ \hline
2008 & 97,6 & 600 \\ \hline
2014 & 99,0 & 710 \\ \hline
2020 & 41,8 & 580 \\ \hline
2022 & 101,0 & 730 \\ \hline
\end{tabularx}
\end{center}
\begin{enumerate}
    \item Relevez les années où le prix du baril a dépassé 95 \$ et indiquez le prix à la pompe correspondant au Cameroun.
    \item Calculez le taux de variation du prix à la pompe entre 2003 et 2022 (arrondi au dixième de pourcent).
    \item Expliquez pourquoi le prix camerounais ne suit pas exactement la courbe du Brent (rôle de la subvention, taux de change, fiscalité, raffinage local vs importation).
    \item Donnez deux autres exemples concrets d'impact des tensions du Moyen-Orient sur le quotidien des Camerounais (ex : coût du fret maritime, prix des engrais/urée, pèlerinage à La Mecque, inflation importée).
\end{enumerate}
\end{exercice}

\courssub{Je me prépare au Bac}

\begin{exercice}[Dissertation type Baccalauréat Technique]
\textbf{Sujet :} « Le Proche et le Moyen-Orient : zones de tension sans fin. » Discutez cette affirmation en vous appuyant sur des exemples précis tirés du problème israélo-palestinien et des conflits du Golfe persique.

\textit{Consignes :} Vous rédigerez une introduction (accroche, définition des termes, délimitation spatio-temporelle, problématique, annonce du plan), un développement structuré en deux ou trois parties équilibrées avec des arguments illustrés par des faits, dates, acteurs, et une conclusion (bilan, ouverture).
\begin{itemize}
    \item \textbf{Partie 1 :} Les racines historiques et la dynamique du conflit israélo-palestinien (promesses contradictoires, 1948, 1967, processus de paix bloqué, colonisation, Jérusalem, Gaza).
    \item \textbf{Partie 2 :} Le Golfe persique, épicentre des rivalités régionales et internationales (pétrole, révolution iranienne 1979, guerre Iran-Irak, guerres du Golfe 1991/2003, tensions arabes/iraniennes, normalisation récente).
    \item \textbf{Partie 3 :} Les facteurs structurels de l'enlisement (ingérences extérieures, absence de sécurité collective, ressources énergétiques, identités nationales/religieuses, asymétrie des puissances).
\end{itemize}
\end{exercice}

\begin{exercice}[Commentaire de carte type Probatoire : Israël et les territoires palestiniens depuis 1967]
On vous propose une carte schématique (non fournie ici, à visualiser mentalement ou sur support) intitulée « Israël et les territoires palestiniens : évolution territoriale 1967-2024 ». Elle fait apparaître : la Ligne verte (frontière d'armistice 1949), l'annexion de Jérusalem-Est (1980), l'annexion du Golan (1981), les blocs de colonies en Cisjordanie, le tracé du « mur de séparation » (barrière de sécurité), les zones A, B, C (accords Oslo II), la bande de Gaza (blocus depuis 2007), le Sinaï (restitué à l'Égypte 1982).

À partir de cette carte, répondez aux questions suivantes :
\begin{enumerate}
    \item \textbf{Présentation du document :} Nature, titre, date, auteur (si connu), échelle, projection.
    \item \textbf{Analyse :}
    \begin{enumerate}
        \item Identifiez et nommez les trois territoires conquis par Israël en juin 1967 encore sous son contrôle direct ou indirect aujourd'hui.
        \item Montrez comment la carte illustre la fragmentation de l'espace palestinien (zones A/B/C, colonies, mur, blocus de Gaza).
        \item Expliquez en quoi le statut de Jérusalem et du Golan constitue un obstacle majeur au droit international (résolutions 242, 338, 478).
    \end{enumerate}
    \item \textbf{Synthèse :} En une dizaine de lignes, dégagez les « nœuds » géographiques et politiques qui rendent la solution « deux États » difficile à mettre en œuvre sur le terrain.
\end{enumerate}
\end{exercice}

\begin{exercice}[Étude de document + Construction graphique : Le pétrole, arme et enjeu du Golfe]
\textbf{Document 1 :} Extrait du discours du roi Fayçal d'Arabie saoudite (octobre 1973) : « Le pétrole des pays arabes ne coulera plus vers les pays qui soutiennent l'agression israélienne... Nous utiliserons l'arme du pétrole pour libérer nos territoires. »
\textbf{Document 2 :} Données de production de pétrole brut (millions de barils/jour) en 1973 et 2022.

\begin{center}
\begin{tabularx}{\linewidth}{|l|X|X|}\hline
\textbf{Pays} & \textbf{1973} & \textbf{2022} \\ \hline
Arabie saoudite & 7,6 & 10,5 \\ \hline
Iran & 5,8 & 2,6 \\ \hline
Irak & 2,0 & 4,5 \\ \hline
Koweït & 3,0 & 3,0 \\ \hline
\end{tabularx}
\end{center}

\begin{enumerate}
    \item \textbf{Étude du document 1 :} Présentez le document. Quel événement déclenche cet embargo ? Quelles en furent les conséquences économiques mondiales (choc pétrolier, stagflation) et politiques (reconnaissance de l'OLP, conférences de paix) ?
    \item \textbf{Construction graphique :} À partir du document 2, construisez un diagramme en barres groupées (deux séries : 1973 et 2022) comparant la production de ces quatre pays. Respectez les règles de la sémiologie graphique (titre, axes légendés, unités, échelle, légende, source).
    \item \textbf{Analyse croisée :} En vous appuyant sur le graphique et vos connaissances, tirez deux conclusions géopolitiques sur l'évolution du rapport de forces au sein de l'OPEP et du Golfe entre 1973 et 2022 (montée en puissance de l'Arabie saoudite, effondrement de la production iranienne post-révolution/guerres/sanctions, résilience koweïtienne, retour de l'Irak).
\end{enumerate}
\end{exercice}

\newpage

\coursec{Corrigés des exercices}

\begin{corrige}[Exercice 1 — QCM : Les origines du conflit israélo-palestinien]
\begin{enumerate}
    \item \textbf{B} — La déclaration Balfour (2 novembre 1917) promet l'établissement en Palestine d'un « Foyer national juif », tout en précisant que les droits civils et religieux des communautés non juives ne seront pas lésés.
    \item \textbf{B} — La résolution 181 (29 novembre 1947) prévoit un État juif, un État arabe et un \emph{corpus separatum} (zone internationale) pour Jérusalem.
    \item \textbf{B} — En juin 1967, Israël occupe la Cisjordanie (avec Jérusalem-Est), Gaza, le Sinaï égyptien et le Golan syrien.
    \item \textbf{B} — Oslo I (13 septembre 1993) : reconnaissance mutuelle Israël/OLP et autonomie palestinienne intérimaire de cinq ans.
\end{enumerate}
\textbf{Points de vigilance :} ne pas confondre la déclaration Balfour (promesse britannique, 1917) avec le plan de partage de l'ONU (1947) ; ne pas attribuer la création de l'OLP (1964) à la guerre des Six Jours (1967).
\end{corrige}

\begin{corrige}[Exercice 2 — Vrai ou Faux justifié]
\begin{enumerate}
    \item \textbf{FAUX} — La guerre Iran-Irak éclate le 22 septembre 1980 par l'attaque de l'Irak de Saddam Hussein contre l'Iran (litige sur le Chatt al-Arab). L'invasion du Koweït date du 2 août 1990 et déclenche la première guerre du Golfe.
    \item \textbf{VRAI} — Après l'invasion du Koweït (août 1990), une coalition de 34 pays dirigée par les États-Unis, autorisée par la résolution 678 de l'ONU, libère le Koweït lors de l'opération « Tempête du désert » (janvier-février 1991).
    \item \textbf{VRAI} — En février 1979, le Chah Mohammad Reza Pahlavi est renversé ; l'ayatollah Khomeyni proclame la République islamique d'Iran (1\textsuperscript{er} avril 1979).
    \item \textbf{FAUX} — Les accords d'Abraham (septembre 2020) normalisent les relations d'Israël avec les Émirats arabes unis et Bahreïn (puis le Maroc et le Soudan), mais pas avec l'Arabie saoudite.
\end{enumerate}
\textbf{Points de vigilance :} chaque justification doit comporter au moins une date ou un acteur précis ; un « vrai » ou « faux » non justifié ne rapporte pas de point.
\end{corrige}

\begin{corrige}[Exercice 3 — Définitions et repères chronologiques]
\begin{enumerate}
    \item \textbf{Sionisme} : mouvement politique national juif né en Europe à la fin du XIX\textsuperscript{e} siècle, théorisé par Theodor Herzl (\emph{Der Judenstaat}, 1896) et organisé lors du premier congrès sioniste de Bâle (1897). Il vise la création d'un foyer national, puis d'un État juif en Palestine, aboutissant à la création d'Israël en 1948.
    \item \textbf{Nakba} : mot arabe signifiant « catastrophe », désignant l'exode d'environ 700 000 Palestiniens fuyant ou expulsés de leurs terres lors de la guerre de 1948, à la suite de la proclamation de l'État d'Israël (14 mai 1948). Il est à l'origine de la question des réfugiés palestiniens, toujours non résolue.
    \item \textbf{OLP} : Organisation de libération de la Palestine, créée en 1964, dirigée par Yasser Arafat à partir de 1969. Reconnue en 1974 par l'ONU comme représentante du peuple palestinien, elle signe les accords d'Oslo (1993) qui instaurent l'Autorité palestinienne.
    \item \textbf{Intifada} : mot arabe signifiant « soulèvement ». La première Intifada (1987-1993) est une révolte populaire palestinienne (grèves, jets de pierres) dans les territoires occupés ; la seconde (2000-2005), déclenchée après la visite d'Ariel Sharon sur l'esplanade des Mosquées et l'échec de Camp David II, est beaucoup plus violente (attentats, répression).
    \item \textbf{Pétrodollars} : recettes en dollars tirées des exportations de pétrole par les pays du Golfe, massivement accrues après le premier choc pétrolier (1973-1974) et recyclées dans les banques occidentales, les achats d'armes et les investissements. Elles font du pétrole une arme économique et diplomatique.
\end{enumerate}
\textbf{Points de vigilance :} une définition complète comporte le terme, sa signification, sa date et son espace ; éviter les définitions circulaires (« l'Intifada est une intifada »).
\end{corrige}

\begin{corrige}[Exercice 4 — Lecture de données : Production pétrolière du Golfe]
\begin{enumerate}
    \item Production totale : $10,5 + 4,5 + 3,0 + 2,6 = 20,6$ Mb/j. Part de l'Arabie saoudite : $\dfrac{10,5}{20,6} \times 100 \approx 51,0\,\%$.
    \item Ordre décroissant : Arabie saoudite (10,5) $>$ Irak (4,5) $>$ Koweït (3,0) $>$ Iran (2,6).
    \item Avec environ la moitié de la production de ce groupe et les plus grandes capacités de réserve, l'Arabie saoudite est le « swing producer » de l'OPEP : elle peut moduler sa production pour influencer les cours mondiaux, ce qui lui confère un poids diplomatique considérable.
\end{enumerate}
\textbf{Points de vigilance :} vérifier que la somme des parts fait 100 \% ; arrondir au dixième comme demandé ($50,97\,\% \approx 51,0\,\%$).
\end{corrige}

\begin{corrige}[Exercice 5 — Carte mentale : Les acteurs du conflit israélo-palestinien]
\textbf{Correction type (structure attendue) :} le nœud central « Conflit israélo-palestinien » doit comporter quatre branches équilibrées, avec des relations fléchées et légendées.
\begin{itemize}
    \item \textbf{Acteurs directs :} Israël (État, armée, colons) ; Autorité palestinienne (Fatah, Cisjordanie) ; Hamas (Gaza depuis 2007) ; Jihad islamique. Relations : opposition armée Israël/Hamas ; rivalité Fatah/Hamas depuis 2007.
    \item \textbf{Acteurs régionaux :} Égypte et Jordanie (paix signée avec Israël en 1979 et 1994, rôle de médiateurs) ; Liban (Hezbollah soutenu par l'Iran) ; Syrie (Golan occupé) ; Iran (soutien au Hamas et au Hezbollah, « axe de la résistance ») ; Arabie saoudite et Turquie (soutien diplomatique aux Palestiniens).
    \item \textbf{Acteurs internationaux :} États-Unis (allié majeur d'Israël, médiateur des accords) ; ONU (résolutions 181, 242, 338 ; UNRWA pour les réfugiés) ; Union européenne et Russie (membres du Quartet avec ONU et USA, créé en 2002).
    \item \textbf{Enjeux :} statut de Jérusalem ; droit au retour des réfugiés ; frontières de 1967 ; colonisation ; partage de l'eau ; sécurité d'Israël.
\end{itemize}
\textbf{Barème indicatif (sur 5) :} nœud central et 4 branches présentes (1 pt) ; acteurs correctement classés (2 pts) ; relations fléchées légendées (1 pt) ; lisibilité et orthographe (1 pt).
\textbf{Points de vigilance :} ne pas mélanger les catégories (le Hamas n'est pas un État régional) ; chaque flèche doit porter une étiquette (alliance, soutien, opposition, médiation).
\end{corrige}

\begin{corrige}[Exercice 6 — Chronologie commentée]
\begin{center}
\begin{tabularx}{\linewidth}{|c|X|X|}\hline
\textbf{Date} & \textbf{Événement} & \textbf{Conséquence immédiate} \\ \hline
1947 & Plan de partage ONU (Rés. 181) & Accepté par les Juifs, refusé par les Arabes ; la guerre devient inévitable. \\ \hline
1948 & Création de l'État d'Israël / Guerre de 1948 & Victoire israélienne ; exode de 700 000 Palestiniens (Nakba). \\ \hline
1967 & Guerre des Six Jours & Israël occupe Cisjordanie, Gaza, Sinaï et Golan ; résolution 242 de l'ONU. \\ \hline
1973 & Guerre du Kippour & Embargo pétrolier arabe et premier choc pétrolier. \\ \hline
1978-1979 & Accords de Camp David & Paix Égypte-Israël ; restitution du Sinaï (1982). \\ \hline
1993 & Accords d'Oslo I & Reconnaissance mutuelle Israël/OLP ; autonomie palestinienne. \\ \hline
2000 & Seconde Intifada & Échec de Camp David II ; reprise des violences. \\ \hline
2005 & Retrait israélien de Gaza & Évacuation des colonies de Gaza ; le Hamas y prend le pouvoir en 2007. \\ \hline
\end{tabularx}
\end{center}
\textbf{Points de vigilance :} la conséquence demandée est \emph{immédiate} : ne pas écrire, pour 1967, « aboutit aux accords d'Oslo » ; respecter l'ordre chronologique et l'exactitude des dates.
\end{corrige}

\begin{corrige}[Exercice 7 — Analyse de document : la résolution 242 de l'ONU]
\begin{enumerate}
    \item \textbf{Présentation :} il s'agit d'une résolution du Conseil de sécurité de l'Organisation des Nations unies, adoptée à l'unanimité le 22 novembre 1967, cinq mois après la guerre des Six Jours (5-10 juin 1967) au cours de laquelle Israël a conquis la Cisjordanie, Jérusalem-Est, Gaza, le Sinaï et le Golan.
    \item \textbf{Deux principes fondamentaux :}
    \begin{itemize}
        \item l'inadmissibilité de l'acquisition de territoires par la guerre, d'où l'exigence du retrait israélien des territoires occupés ;
        \item le droit de chaque État de la région (y compris Israël) à vivre en paix dans des frontières sûres et reconnues : c'est la formule « la paix contre les territoires ».
    \end{itemize}
    \item \textbf{Ambiguïté linguistique :} le texte anglais, qui fait foi, dit « withdrawal from territories » (retrait de territoires, sans article défini), ce qui laisse entendre un retrait partiel ; la version française « retrait des territoires » implique un retrait total. Israël s'appuie sur la version anglaise pour conserver des territoires (Jérusalem-Est, blocs de colonies), tandis que les Palestiniens et les États arabes exigent le retrait complet de toutes les conquêtes de 1967.
    \item \textbf{Référence durable :} la résolution 242, complétée par la résolution 338 (1973), fonde tous les processus de paix ultérieurs : Camp David (1978), Madrid (1991), Oslo (1993), feuille de route du Quartet (2003). Elle reste la base juridique du principe « deux États dans des frontières sûres et reconnues », encore invoqué dans les négociations actuelles.
\end{enumerate}
\textbf{Barème indicatif (sur 10) :} présentation complète (2 pts) ; deux principes identifiés et cités (3 pts) ; explication de l'ambiguïté (3 pts) ; mise en perspective (2 pts).
\textbf{Points de vigilance :} citer le numéro et la date exacte de la résolution ; distinguer clairement les deux principes ; ne pas écrire que la résolution crée un État palestinien (elle ne le mentionne pas).
\end{corrige}

\begin{corrige}[Exercice 8 — Mise en relation : Crises du Golfe et économie camerounaise]
\begin{enumerate}
    \item Trois années dépassent 95 \$ : 2008 (97,6 \$ ; 600 FCFA/L), 2014 (99,0 \$ ; 710 FCFA/L) et 2022 (101,0 \$ ; 730 FCFA/L).
    \item Taux de variation : $\dfrac{730 - 440}{440} \times 100 = \dfrac{290}{440} \times 100 \approx 65,9\,\%$ de hausse entre 2003 et 2022.
    \item Le prix à la pompe ne suit pas exactement le Brent car il dépend aussi de la subvention de l'État (qui lisse les hausses, comme en 2008), du taux de change FCFA/dollar, des taxes et marges de distribution, et du fait que le Cameroun importe une partie de ses produits raffinés malgré la SONARA.
    \item Exemples : hausse du fret maritime (détour par le cap de Bonne-Espérance lors des tensions en mer Rouge) qui renchérit les importations ; flambée du prix des engrais (urée) liée au gaz ; coût du pèlerinage à La Mecque ; inflation importée sur les produits manufacturés.
\end{enumerate}
\textbf{Points de vigilance :} dans le calcul du taux, diviser par la valeur de départ (440), pas par la valeur d'arrivée ; bien distinguer corrélation (les deux prix montent) et causalité directe (le prix camerounais est administré).
\end{corrige}

\begin{corrige}[Exercice 9 — Dissertation type Baccalauréat Technique]
\textbf{Correction guidée (plan détaillé et éléments de rédaction).}

\textbf{Introduction (attendus) :}
\begin{itemize}
    \item \emph{Accroche :} depuis 1948, le Proche et le Moyen-Orient n'ont connu presque aucune année sans conflit ouvert ; la guerre entre Israël et le Hamas à Gaza (2023-2024) en est l'illustration récente.
    \item \emph{Définitions et délimitation :} le Proche-Orient désigne l'espace Israël/Palestine et ses voisins immédiats ; le Moyen-Orient inclut le golfe Persique (Iran, Irak, péninsule Arabique). Période étudiée : 1917 (déclaration Balfour) à nos jours.
    \item \emph{Problématique :} pourquoi cette région concentre-t-elle des conflits récurrents, et quels facteurs expliquent que la paix y demeure introuvable ?
    \item \emph{Annonce du plan :} nous verrons d'abord les racines et la dynamique du conflit israélo-palestinien, puis le Golfe comme épicentre des rivalités, enfin les facteurs structurels de l'enlisement.
\end{itemize}

\textbf{Partie 1 — Les racines historiques et la dynamique du conflit israélo-palestinien :} promesses contradictoires (Hussein-McMahon 1915, Sykes-Picot 1916, Balfour 1917) ; plan de partage de 1947 refusé par les Arabes ; guerre de 1948 et Nakba (700 000 réfugiés) ; guerre des Six Jours (1967) et occupation ; première et seconde Intifadas (1987, 2000) ; échecs de Camp David II (2000) ; colonisation continue, statut de Jérusalem, blocus de Gaza depuis 2007, prise du pouvoir du Hamas (2007) et guerres récurrentes à Gaza.

\textbf{Partie 2 — Le Golfe persique, épicentre des rivalités régionales et internationales :} richesse pétrolière (environ la moitié des réserves mondiales) ; révolution iranienne de 1979 et rupture avec l'Occident ; guerre Iran-Irak (1980-1988, environ un million de morts) ; invasion du Koweït (1990) et guerre du Golfe (1991) ; intervention américaine en Irak (2003) et chaos régional ; rivalité Iran/Arabie saoudite (Yémen, Liban, Syrie) ; accords d'Abraham (2020) comme tentative de recomposition.

\textbf{Partie 3 — Les facteurs structurels de l'enlisement :} ingérences des grandes puissances (USA, URSS puis Russie) ; absence de système de sécurité collective régional ; enjeu énergétique mondial ; conflits d'identités (nationales, religieuses : sunnites/chiites) ; asymétrie militaire et impasses du droit international (résolutions 242 et 338 non appliquées).

\textbf{Conclusion (attendus) :} bilan — les tensions sont « sans fin » car elles superposent un conflit territorial non résolu, des rivalités régionales et des intérêts énergétiques mondiaux ; ouverture possible — les accords d'Abraham ou la médiation Chine-Iran/Arabie saoudite (2023) montrent que des recompositions sont possibles, mais la question palestinienne reste le nœud central.

\textbf{Barème indicatif (sur 20) :} introduction complète avec problématique (4 pts) ; respect du plan et équilibre des parties (3 pts) ; exactitude des faits, dates, acteurs (6 pts) ; argumentation et articulation logique (4 pts) ; conclusion avec ouverture (2 pts) ; expression et orthographe (1 pt).
\textbf{Points de vigilance :} ne pas faire un récit chronologique sans problématique ; chaque argument doit être illustré par un fait daté ; éviter les jugements de valeur et les généralisations sur les religions ou les peuples.
\end{corrige}

\begin{corrige}[Exercice 10 — Commentaire de carte type Probatoire]
\begin{enumerate}
    \item \textbf{Présentation :} carte politique et territoriale intitulée « Israël et les territoires palestiniens : évolution territoriale 1967-2024 », de petite échelle (échelle régionale), réalisée à des fins pédagogiques ou documentaires ; elle couvre la période postérieure à la guerre des Six Jours.
    \item \textbf{Analyse :}
    \begin{enumerate}
        \item Trois territoires conquis en juin 1967 encore sous contrôle israélien direct ou indirect : la \textbf{Cisjordanie} (occupée et colonisée), \textbf{Jérusalem-Est} (annexée en 1980) et le \textbf{plateau du Golan} (annexé en 1981). La bande de Gaza, évacuée en 2005, reste sous blocus israélien depuis 2007 ; le Sinaï a été restitué à l'Égypte (1982).
        \item \textbf{Fragmentation de l'espace palestinien :} la carte montre la Cisjordanie morcelée en zones A (autonomie palestinienne), B (contrôle mixte) et C (contrôle israélien total, environ 60 \% du territoire) selon Oslo II (1995) ; les blocs de colonies reliés par des routes réservées ; le tracé du « mur de séparation » (commencé en 2002) qui empiète au-delà de la Ligne verte ; Gaza séparée physiquement de la Cisjordanie et soumise au blocus.
        \item \textbf{Jérusalem et le Golan face au droit international :} l'annexion de Jérusalem-Est (1980) est déclarée « nulle et non avenue » par la résolution 478 de l'ONU ; l'annexion du Golan (1981) n'est reconnue par presque aucun État (les États-Unis l'ont reconnue en 2019, à titre exceptionnel). Ces annexions violent le principe de la résolution 242 : inadmissibilité de l'acquisition de territoires par la guerre.
    \end{enumerate}
    \item \textbf{Synthèse (une dizaine de lignes) :} la solution « deux États » se heurte à plusieurs nœuds : la colonisation continue a créé des faits accomplis qui rendent le tracé d'une frontière viable presque impossible ; le statut de Jérusalem, revendiquée comme capitale par les deux peuples, bloque tout compromis ; la séparation physique et politique entre Gaza (Hamas) et la Cisjordanie (Autorité palestinienne) prive les Palestiniens d'un interlocuteur unique ; le mur et les zones C fragmentent la continuité territoriale indispensable à un État palestinien viable ; enfin, la question des réfugiés et la demande israélienne de sécurité absolue restent sans solution négociée. La carte montre ainsi que l'écart entre le droit international (résolutions 242, 338) et les réalités du terrain n'a cessé de croître depuis 1967.
\end{enumerate}
\textbf{Barème indicatif (sur 20) :} présentation du document (3 pts) ; identification des territoires (3 pts) ; analyse de la fragmentation (5 pts) ; statut de Jérusalem et du Golan (4 pts) ; synthèse structurée (4 pts) ; expression (1 pt).
\textbf{Points de vigilance :} en commentaire de carte, toujours commencer par décrire avant d'expliquer ; citer les éléments de la légende (Ligne verte, zones A/B/C, mur) ; ne pas confondre annexion (Jérusalem-Est, Golan) et occupation (Cisjordanie).
\end{corrige}

\begin{corrige}[Exercice 11 — Étude de document + Construction graphique]
\begin{enumerate}
    \item \textbf{Étude du document 1 :} il s'agit d'un extrait du discours du roi Fayçal d'Arabie saoudite, prononcé en octobre 1973, pendant la guerre du Kippour (déclenchée le 6 octobre 1973 par l'attaque surprise de l'Égypte et de la Syrie contre Israël). Pour faire pression sur les Occidentaux soutenant Israël, les pays arabes de l'OPEP décrètent un embargo pétrolier et des réductions de production. Conséquences économiques : quadruplement du prix du baril (de 3 à 12 \$), premier choc pétrolier, crise énergétique, stagflation (inflation + stagnation) dans les pays industrialisés. Conséquences politiques : prise de conscience de la dépendance énergétique (création de l'Agence internationale de l'énergie, 1974), reconnaissance de l'OLP comme représentante du peuple palestinien à l'ONU (1974), relance des efforts de paix menant à Camp David (1978).
    \item \textbf{Construction graphique (corrigé du diagramme attendu) :}
\begin{popfigure}
\begin{tikzpicture}
\begin{axis}[
    ybar, bar width=0.35cm,
    width=0.9\linewidth, height=6.5cm,
    ylabel={Production (millions de barils/jour)},
    symbolic x coords={Arabie saoudite, Iran, Irak, Kowe\"it},
    xtick=data, x tick label style={rotate=15, anchor=east, font=\small},
    ymin=0, ymax=12,
    legend style={at={(0.5,1.02)}, anchor=south, legend columns=2},
    nodes near coords, every node near coord/.append style={font=\scriptsize},
]
\addplot[fill=popblue] coordinates {(Arabie saoudite,7.6) (Iran,5.8) (Irak,2.0) (Kowe\"it,3.0)};
\addplot[fill=poporange] coordinates {(Arabie saoudite,10.5) (Iran,2.6) (Irak,4.5) (Kowe\"it,3.0)};
\legend{1973, 2022}
\end{axis}
\end{tikzpicture}
\legende{Production de pétrole brut de quatre pays du Golfe, 1973 et 2022 (source : OPEP/BP Statistical Review).}
\end{popfigure}
    \item \textbf{Analyse croisée (deux conclusions) :}
    \begin{itemize}
        \item \emph{Montée en puissance de l'Arabie saoudite :} sa production passe de 7,6 à 10,5 Mb/j (+38 \%), confirmant son rôle de « swing producer » et de leader de l'OPEP, capable d'influencer les cours mondiaux.
        \item \emph{Effondrement relatif de l'Iran :} de 5,8 Mb/j en 1973 (deuxième producteur du groupe), l'Iran tombe à 2,6 Mb/j en 2022, conséquence de la révolution de 1979, de la guerre Iran-Irak (1980-1988) et des sanctions internationales ; l'Irak (2,0 à 4,5 Mb/j) retrouve le second rang malgré les guerres, tandis que le Koweït maintient sa production (3,0 Mb/j), signe de résilience après la guerre de 1991.
    \end{itemize}
\end{enumerate}
\textbf{Barème indicatif (sur 20) :} présentation et analyse du discours (6 pts) ; graphique conforme aux règles de sémiologie (8 pts : titre 1, axes et unités 2, échelle 1, barres exactes 3, légende/source 1) ; deux conclusions argumentées (5 pts) ; expression (1 pt).
\textbf{Points de vigilance :} l'axe des ordonnées doit partir de 0 pour ne pas tromper le lecteur ; les valeurs doivent être reportées avec exactitude (Iran 2022 : 2,6 et non 6,2) ; relier explicitement le graphique au discours (le pétrole comme arme suppose des volumes de production importants).
\end{corrige}

\newpage

\coursec{Fiche bilan}

\begin{popfigure}
\centering
\begin{tikzpicture}[
  mindmap,
  concept color=popblue,
  level 1 concept/.append style={level distance=4.5cm, sibling angle=60, font=\small\bfseries},
  level 2 concept/.append style={level distance=3cm, sibling angle=45, font=\footnotesize},
  every node/.style={concept, execute at begin node=\hskip0pt},
  branch1/.style={concept color=poporange, edge from parent path={(\tikzparentnode) -- (\tikzchildnode)}},
  branch2/.style={concept color=popgreen, edge from parent path={(\tikzparentnode) -- (\tikzchildnode)}},
  branch3/.style={concept color=poppurple, edge from parent path={(\tikzparentnode) -- (\tikzchildnode)}},
  branch4/.style={concept color=poppink, edge from parent path={(\tikzparentnode) -- (\tikzchildnode)}},
  branch5/.style={concept color=popgold, edge from parent path={(\tikzparentnode) -- (\tikzchildnode)}},
  branch6/.style={concept color=popblue, edge from parent path={(\tikzparentnode) -- (\tikzchildnode)}}
]
\node[concept, text width=3.5cm, align=center, font=\normalsize\bfseries] {Proche \& Moyen-Orient : Zones de tension sans fin}
  child[branch1] { node[concept, text width=2.8cm, align=center] {Origines du conflit israélo-palestinien\\ (Sionisme, Mandat, 1948)} }
  child[branch2] { node[concept, text width=2.8cm, align=center] {Guerres, Intifadas \& Processus de paix\\ (1956, 1967, 1973, 1987, 2000, Oslo)} }
  child[branch3] { node[concept, text width=2.8cm, align=center] {Conflits du Golfe persique\\ (Iran-Irak 1980-88, Koweït 1990-91, Irak 2003)} }
  child[branch4] { node[concept, text width=2.8cm, align=center] {Acteurs \& Enjeux majeurs\\ (Pétrole, Religion, Territoire, Puissances)} }
  child[branch5] { node[concept, text width=2.8cm, align=center] {Géopolitique régionale\\ (Iran, Turquie, Arabie Saoudite, Syrie, Liban)} }
  child[branch6] { node[concept, text width=2.8cm, align=center] {Répercussions internationales \& Cameroun\\ (Terrorisme, Prix énergie, Diplomatie)} };
\end{tikzpicture}
\legende{Carte mentale de synthèse : Le Proche et le Moyen-Orient, zones de tension sans fin (Terminale Technique).}
\end{popfigure}

\begin{bilan}
\courssub{Définitions \& Concepts clés}
\begin{itemize}
  \item \cle{Sionisme} : Mouvement nationaliste juif (fin XIXe s.) visant à créer un foyer national en Palestine (Théodore Herzl, Congrès de Bâle 1897).
  \item \cle{Déclaration Balfour} (1917) : Lettre du gouvernement britannique (Lord Balfour) favorable à un « foyer national juif » en Palestine, tout en préservant les droits des communautés non-juives.
  \item \cle{Mandat britannique} (1922-1948) : Administration de la Palestine par la Grande-Bretagne sous égide SDN ; tensions entre immigration juive et population arabe.
  \item \cle{Plan de partage ONU} (Résolution 181, nov. 1947) : Partage de la Palestine en deux États (juif et arabe) et internationalisation de Jérusalem ; rejeté par les Arabes, accepté par l'Agence juive.
  \item \cle{Nakba} (« Catastrophe », 1948) : Exode de 700 000 à 750 000 Palestiniens lors de la guerre de 1948-1949 ; création de l'État d'Israël (14 mai 1948).
  \item \cle{OLP} (Organisation de Libération de la Palestine, 1964) : Représentant légitime du peuple palestinien (Yasser Arafat) ; reconnaissance ONU 1974.
  \item \cle{Intifada} : Soulèvement populaire palestinien (1ère : 1987-1993 ; 2ème : 2000-2005) contre l'occupation israélienne des Territoires (Cisjordanie, Gaza).
  \item \cle{Accords d'Oslo} (1993, 1995) : Reconnaissance mutuelle Israël/OLP, création de l'Autorité palestinienne, transfert partiel de compétences ; processus inachevé (statut final : Jérusalem, réfugiés, colonies, frontières).
  \item \cle{Pétrole} : Ressource stratégique du Golfe (2/3 des réserves mondiales) ; enjeu de puissance pour les USA, l'URSS/Russie, la Chine, l'Europe.
  \item \cle{Chiisme / Sunnisme} : Fracture confessionnelle majeure (Iran chiite vs Arabie Saoudite sunnite) structurant les alliances régionales (axe de la résistance vs axe modéré).
\end{itemize}

\courssub{Repères chronologiques indispensables (à situer sur une frise)}
\begin{center}
\begin{tabularx}{\linewidth}{|c|X|X|}\hline
\rowcolor{popblueL}\textbf{Date} & \textbf{Événement majeur} & \textbf{Conséquence clé} \\ \hline
1916 & Accords Sykes-Picot & Partage ottoman (zones d'influence FR/UK) \\ \hline
1917 & Déclaration Balfour & Promesse foyer juif en Palestine \\ \hline
1947 & Résolution 181 (ONU) & Plan de partage rejeté par Arabes \\ \hline
1948 & Création État d'Israël / 1ère guerre & Nakba, 700k réfugiés, Armistice 1949 \\ \hline
1956 & Crise de Suez (Guerre 56) & Échec franco-britannique, hégémonie US/URSS \\ \hline
1967 & Guerre des Six Jours & Israël occupe Sinaï, Golan, Cisjordanie, Gaza, Jérusalem-Est (Résolution 242) \\ \hline
1973 & Guerre du Kippour / Choc pétrolier & Rééquilibrage militaire, diplomatie US (Kissinger), 1er choc pétrole \\ \hline
1979 & Révolution iranienne / Traité Camp David & République islamique Iran ; Paix Égypte-Israël (Sinaï restitué) \\ \hline
1980-88 & Guerre Iran-Irak & 1 million de morts, statu quo, affaiblissement des deux \\ \hline
1987-93 & 1ère Intifada & Internationalisation cause palestinienne, Conférence Madrid 1991 \\ \hline
1990-91 & Guerre du Golfe (Koweït) & Coalition internationale (ONU/US), Irak sanctionné \\ \hline
1993-95 & Accords Oslo I \& II & Reconnaissance mutuelle, Autorité palestinienne, prix Nobel paix \\ \hline
2000-05 & 2ème Intifada (Al-Aqsa) & Échec Camp David 2000, mur de séparation, Hamas au pouvoir (Gaza 2007) \\ \hline
2003 & Guerre d'Irak (2ème) & Renversement Saddam Hussein, chaos régional, montée Daech \\ \hline
2011- & Printemps arabes / Guerre Syrie & Recomposition alliances, intervention russes/iraniennes/turques \\ \hline
2023- & Guerre Israël-Hamas (7 oct.) & Crise humanitaire Gaza, risque embrasement régional (Hezbollah, Iran, Houthis) \\ \hline
\end{tabularx}
\end{center}

\courssub{Méthodes express (Bac : Commentaire de doc / Dissertation)}
\begin{methode}[Analyser un document sur le conflit]
  \begin{enumerate}
    \item \textbf{Identifier} : Nature (carte, discours, traité, photo, caricature, stats), Auteur, Date, Contexte précis.
    \item \textbf{Structurer} : 2 ou 3 parties logiques (ex: Origines / Déclenchement / Conséquences ; ou Acteurs / Enjeux / Issu).
    \item \textbf{Exploiter} : Citer le doc (guillemets), croiser avec connaissances (dates, acteurs, résolutions ONU).
    \item \textbf{Conclure} : Répondre à la problématique, ouvrir sur l'actualité (impasse actuelle, rôle puissances).
  \end{enumerate}
\end{methode}

\begin{methode}[Dissertation : « Pourquoi tensions sans fin ? »]
  \begin{enumerate}
    \item \textbf{Analyser le sujet} : Définir « Proche-Orient » (Levant) vs « Moyen-Orient » (Golfe + Iran), « sans fin » = chronicité, récurrence, échec solutions durables.
    \item \textbf{Problématiser} : En quoi la superposition de conflits (national, religieux, pétrolier, proxy) empêche la stabilité ?
    \item \textbf{Plan type} :
      \begin{itemize}
        \item I. Un foyer conflictuel historique : la question de Palestine (nationalismes rivaux, droit international vs faits accomplis).
        \item II. Un système régional instable : rivalités de puissance (Iran/Arabie/Turquie), ressource pétrole, ingérences extérieures (US, Russie).
        \item III. L'impasse des règlementations : limites du droit international (résolutions non appliquées), échec processus paix, radicalisation (islamisme, nationalisme religieux).
      \end{itemize}
  \end{enumerate}
\end{methode}
\end{bilan}

\begin{aretenir}[Checklist avant le Bac \texteuro{} Histoire]
\begin{itemize}
  \item $\square$ Je sais situer \textbf{Proche-Orient} (Levant : Israël, Palestine, Liban, Syrie, Jordanie) et \textbf{Moyen-Orient} (Golfe : Irak, Koweït, Arabie, Iran, Émirats, Qatar, Bahreïn, Oman + Yémen) sur une carte muette.
  \item $\square$ Je maîtrise la chronologie : \textbf{1917 Balfour} $\rightarrow$ \textbf{1947 Partage} $\rightarrow$ \textbf{1948 Israël/Nakba} $\rightarrow$ \textbf{1967 Six Jours} $\rightarrow$ \textbf{1973 Kippour} $\rightarrow$ \textbf{1993 Oslo} $\rightarrow$ \textbf{2000 2ème Intifada} $\rightarrow$ \textbf{2023 7 Octobre}.
  \item $\square$ J'explique les \textbf{3 piliers du sionisme} (Terre, Peuple, Souveraineté) et la \textbf{revendication palestinienne} (Droit au retour R194, État viable frontières 67, Jérusalem-Est capitale).
  \item $\square$ Je cite les \textbf{résolutions ONU clés} : 181 (partage), 242 (retrait territoires occupés 67), 338 (cessez-le-feu 73), 194 (droit retour réfugiés).
  \item $\square$ Je distingue les \textbf{guerres interétatiques} (1948, 56, 67, 73, Iran-Irak 80-88, Golfe 91, Irak 2003) des \textbf{conflits asymétriques} (Intifadas, Gaza, Hezbollah, Houthis, Daech).
  \item $\square$ J'analyse le \textbf{rôle du pétrole} : Chocs 1973/1979, arme politique, dollar/pétrole, enjeu transition énergétique actuelle.
  \item $\square$ Je comprends la \textbf{rivalité Iran / Arabie Saoudite} (chiisme/sunnisme, influence régionale, guerre par procuration Yemen/Syrie/Liban/Irak).
  \item $\square$ J'évalue l'\textbf{impact des puissances extérieures} : US (soutien Israël, sécurité Golfe), Russie (retour Syrie 2015), Chine (médiation Iran/SA 2023, routes soie), UE (aide, processus paix).
  \item $\square$ Je relie les tensions à la \textbf{situation du Cameroun} : Coût hydrocarbures (budget, transport), Sécurité (Boko Haram/Daech lien idéologique), Diplomatie (vote ONU, relations Israël/Palestine), Diaspora.
  \item $\square$ Je sais rédiger une \textbf{intro de dissertation} (accroche, définitions, sujet, problématique, annonce plan) et une \textbf{conclusion} (bilan, ouverture géopolitique).
  \item $\square$ Je sais faire un \textbf{commentaire de carte} (légende, titres, flèches, zones, gradient, analyse spatiale : frontières, colonies, ressources, alliances).
\end{itemize}
\end{aretenir}

\findecours

\end{document}
